% concatenated from document.tex
\documentclass[a4paper, 12pt]{amsart}
\usepackage{amsfonts, amsthm, amssymb, amsmath, stmaryrd}
\usepackage{mathrsfs,array}
\usepackage{xy}
\usepackage{verbatim}
\usepackage{amscd} 
\usepackage{tikz}
\usepackage{tikz-cd}
\usetikzlibrary{matrix,arrows,decorations.pathmorphing}
\input xy
\xyoption{all}
\usepackage[unicode]{hyperref}



\setlength{\textwidth}{6.5in}
\setlength{\oddsidemargin}{-0.1in}
\setlength{\evensidemargin}{-0.1in}

\setcounter{tocdepth}{2}


\newtheorem*{thmn}{Theorem} \newtheorem{ithm}{Theorem}
\newtheorem{athm}{Theorem}[section]
\newtheorem{acorollary}[athm]{Corollary}
\newtheorem{acor}[athm]{Corollary}
\newtheorem{alemma}[athm]{Lemma}
\newtheorem{aprop}[athm]{Proposition}
\newtheorem{thm}{Theorem}[subsection]
\newtheorem{corollary}[thm]{Corollary}
\newtheorem{cor}[thm]{Corollary}
\newtheorem{lemma}[thm]{Lemma}
\newtheorem{lem}[thm]{Lemma} 
\newtheorem{prop}[thm]{Proposition}
\newtheorem{conj}[thm]{Conjecture} \theoremstyle{definition}
\newtheorem{conjecture}[thm]{Conjecture} \theoremstyle{definition}
\newtheorem{defn}[thm]{Definition} \theoremstyle{remark}
\newtheorem{rem}[thm]{\bf Remark}
\newtheorem{remark}[thm]{\bf Remark}
\newtheorem{para}[thm]{\bf} 
\newtheorem{observation}[thm]{Observation} 
\newtheorem{assumption}[thm]{\bf Assumption}

\renewcommand{\Bbb}{\mathbb}
\renewcommand{\bold}{\mathbf}


\DeclareMathOperator{\Id}{Id}
\DeclareMathOperator{\Hom}{Hom}
\DeclareMathOperator{\End}{End}
\DeclareMathOperator{\Ext}{Ext}
\DeclareMathOperator{\Tor}{Tor}
\DeclareMathOperator{\Op}{Op}
\DeclareMathOperator{\Stop}{Stop}
\DeclareMathOperator{\POp}{P}
\DeclareMathOperator{\Ker}{ker}
\DeclareMathOperator{\im}{im}
\DeclareMathOperator{\coker}{coker}
\DeclareMathOperator{\modd}{mod}
\DeclareMathOperator{\diag}{diag}
\DeclareMathOperator{\Spa}{Spa}
\DeclareMathOperator{\Spec}{Spec}
\DeclareMathOperator{\Spm}{Spm}
\DeclareMathOperator{\Spf}{Spf}
\DeclareMathOperator{\Fil}{Fil}
\DeclareMathOperator{\gr}{gr}
\DeclareMathOperator{\Gal}{Gal}
\DeclareMathOperator{\Proj}{Proj}
\DeclareMathOperator{\Lie}{Lie}
\DeclareMathOperator{\Frob}{Frob}
\DeclareMathOperator{\Res}{Res}
\DeclareMathOperator{\Ver}{Ver}
\DeclareMathOperator{\Sym}{Sym}
\DeclareMathOperator{\Set}{Set}
\DeclareMathOperator{\hht}{ht}
\DeclareMathOperator{\Aut}{Aut}
\DeclareMathOperator{\ad}{ad}
\DeclareMathOperator{\Supp}{Supp}


\def\cont{\mathrm{cont}}
\def\cris{\mathrm{cris}}
\def\crys{\mathrm{crys}}
\def\rig{\mathrm{rig}}
\def\dR{\mathrm{dR}}
\def\id{\mathrm{id}}
\def\HT{\mathrm{HT}}
\def\St{\mathrm{St}}
\def\uSym{\underline{\Sym}}
\def\Ind{\mathrm{Ind}}
\def\Mfc{\mathrm{Mod^{fl}_{comp}(O_E)}}
\def\Fix{\mathrm{Fix}}
\def\patch{\mathrm{p}}
\def\ps{\mathrm{ps}}
\def\loc{\mathrm{loc}}
\def\ord{\mathrm{ord}}
\def\depth{\mathrm{depth}}
\def\Irr{\mathrm{Irr}}
\def\sm{\mathrm{sm}}
\def\Ban{\mathrm{Ban}}
\def\adm{\mathrm{adm}}
\def\Mod{\mathrm{Mod}}

\newcommand{\N}{\mathbb{N}}
\newcommand{\Z}{\mathbb{Z}}
\newcommand{\F}{\mathbb{F}}
\newcommand{\Q}{\mathbb{Q}}
\newcommand{\R}{\mathbb{R}}
\newcommand{\T}{\mathbb{T}}
\newcommand{\A}{\mathbb{A}}
\newcommand{\bC}{\mathbb{C}}
%\newcommand{\G}{\mathbb{G}}
\newcommand{\fR}{\mathbf{R}}
\newcommand{\II}{\mathcal{I}}
\newcommand{\tr}{\mathrm{tr}}
\newcommand{\GL}{\mathrm{GL}}
\newcommand{\SL}{\mathrm{SL}}
\newcommand{\PGL}{\mathrm{PGL}}
\newcommand{\GSp}{\mathrm{GSp}}
\newcommand{\Sp}{\mathrm{Sp}}
\newcommand{\Fl}{{\mathscr{F}\!\ell}}
\newcommand{\cA}{\mathcal{A}}
\newcommand{\cC}{\mathcal{C}}
\newcommand{\cI}{\mathcal{I}}
\newcommand{\cO}{\mathcal{O}}
\newcommand{\cL}{\mathcal{L}}
\newcommand{\defeq}{\stackrel{\mathrm{def}}{=}}
\newcommand{\PP}{\mathbb{P}}
\newcommand{\indkg}{\Ind_{\GL_2(\Z_p)\Q_p^\times}^{\GL_2(\Q_p)}}
\newcommand{\ebf}{\mathbf{e}}
\newcommand{\cOf}{C_\cO^\mathrm{f}}
\newcommand{\kq}{\mathfrak{q}}
\newcommand{\kp}{\mathfrak{p}}
\newcommand{\kkp}{k(\kp)}
\newcommand{\kQ}{\mathfrak{Q}}
\newcommand{\kP}{\mathfrak{P}}
\newcommand{\kC}{\mathfrak{C}}
\newcommand{\kB}{\mathfrak{B}}
\newcommand{\kF}{\mathfrak{F}}
\newcommand{\ka}{\mathfrak{a}}
\newcommand{\kb}{\mathfrak{b}}
\newcommand{\km}{\mathfrak{m}}
\newcommand{\mm}{\mathrm{m}}
\newcommand{\DAi}{(D\otimes_F \A_F^\infty)^\times}
\newcommand{\AFi}{(\A_F^\infty)^\times}
\newcommand{\X}{\mathcal{X}}

\newcommand{\RNum}[1]{\uppercase\expandafter{\romannumeral #1\relax}}

\newcommand{\overbar}[1]{\mkern 1.5mu\overline{\mkern-1.5mu#1\mkern-1.5mu}\mkern 1.5mu}

\begin{document}
	\title{Zariski density of modular points in the Eisenstein case}
	\author{Xinyao Zhang}
	\address{Graduate School of Mathematical Sciences, University of Tokyo, Komaba, Meguro, Tokyo 153-8914, Japan}
	\email{zhangxy96@g.ecc.u-tokyo.ac.jp}
	\begin{abstract}
		In this article, we study the Zariski closure of modular points in the two-dimensional universal deformation space when the residual Galois representation is reducible. Unlike the previous approaches in the residually irreducible case from Gouv\^ea-Mazur, B\"ockle and Allen, our method relies on local-global compatibility results, potential pro-modularity arguments and a non-ordinary finiteness result between the local deformation ring at $p$ and the global deformation ring. This allows us to construct sufficiently many non-ordinary regular de Rham points whose modularity is guaranteed by the recent progress on the Fontaine-Mazur conjecture. Also, we will discuss some applications of our main results, including the equidimensionality of certain big Hecke algebras and big $R=\T$ theorems in the residually reducible case.
	\end{abstract}
	
	\maketitle
	
	\tableofcontents
	
	\section{Introduction}
	 
	 Let $E$ be a number field and $p$ be a prime. Let $\Sigma$ be a finite set of finite places of $E$ containing all those lying above $p$. Let $\bar{\rho} : \Gal(\overline{E}/E) \to \GL_2(\overline{\F_p})$ be a continuous Galois representation unramified outside $\Sigma$. By classical Galois deformation theory (see \cite[Proposition 2.3.1, page 42]{Berger_2013}), there exists a universal pseudo-deformation ring $R_E^\ps$ parametrizing all pseudo-representations of $\tr \bar{\rho}$ (for the profinite group $G_{E, \Sigma}$). The geometry of the set of modular/automorphic points in the universal deformation space $\Spec R_E^\ps$ has attracted much interest among number theorists. For example, we wonder whether it is Zariski dense.
	 
	 \subsection{The previous work}
	 We first recall some previous progress in the residually irreducible case. Note that when $\bar{\rho}$ is irreducible, $R_E^\ps$ is just the universal deformation ring of $\bar{\rho}$ (see \cite[Theorem 2.4.1, page 44]{Berger_2013}). 
	 
	 In \cite{gouvea}, Gouv\^ea and Mazur first studied this question in the case $E=\Q$, $\Sigma = \{p\}$, and $\bar{\rho}$ is irreducible, unobstructed and modular (it is known by Serre's modularity conjecture in \cite{Khare_2009a} and \cite{Khare_2009} at present). Under these settings, they showed that modular points are Zariski dense in $\Spec R_\Q^\ps$, and $ R_\Q^\ps$ is isomorphic to some big Hecke algebra, using the theory of eigencurves (the so-called \textit{infinite fern}).
	
		Later in \cite{B_ckle_2001}, B\"ockle refined Gouv\^ea-Mazur's strategy, allowing a much broader generalization. He first showed that each irreducible component of $\Spec R_\Q^\ps$ contains a smooth modular point under some technical assumptions, using explicit calculations of some local deformation rings (in the Fontaine-Laffaille case or the ordinary case) and small $R=\T$ theorem in the sense of Taylor-Wiles. The existence of a smooth point implies that the irreducible component has the correct Krull dimension. Then combining the infinite fern arguments, we can show that modular points are Zariski dense in each irreducible component and hence in $\Spec R_\Q^\ps$. In \cite{emerton11}, Emerton generalized B\"ockle's result to obtain a big $R=\T$ theorem (see \cite[Theorem 1.2.3]{emerton11}) and further proved the two-dimensional Fontaine-Mazur conjecture in the non-ordinary case under the Taylor--Wiles hypothesis.
	
	 Following B\"ockle's strategy, Allen proved similar results in the context of polarized Galois representations for CM fields and in the Hilbert modular case for some totally real fields in \cite{Allen_2019}. The main new idea there is to use the finiteness of the universal global (polarized) deformation ring $R^{\operatorname{pol}}$ over the local framed deformation ring $R^{\operatorname{loc}}$ at $p$ (see \cite{Allen_2014}). After assuming some local conditions on $\bar{\rho}|_{G_{E_v}}$ for $v|p$, the local framed deformation ring will be a regular local ring. Then by intersection theory (see \cite[Lemma 1.1.2]{Allen_2019}), Allen showed that each irreducible component of $ \Spec R^{\operatorname{pol}}$ contains a modular point with prescribed $p$-adic Hodge theoretic type, which is also smooth by using an argument of the vanishing of the geometric Bloch-Kato Selmer group (see \cite{Allen_2016}). The Zariski density follows from Chenevier's generalization of the infinite fern result \cite{chenevier2011}.
	
	Recently in \cite{Deo2023}, Deo proved some big $R=\T$ theorem for $E=\Q$ in the residually reducible case under some technical assumptions, especially the cyclicity of certain cohomology group. He also followed B\"ockle's strategy and used Pan's ordinary finiteness result \cite[Theorem 5.1.1]{pan2022fontaine} instead of small $R=\T$ theorems to find a smooth point in each irreducible component of $\Spec R_\Q^\ps$. 
	
	We emphasize here that in each work mentioned above, an infinite fern argument is necessary. One can see \cite{gouvea}, \cite{chenevier2011}, \cite{Emerton2011}, \cite{Hellmann_2022} and \cite{Hernandez:2022aa} for more discussions on this topic.
	
	\subsection{Main results}
	Now we state our main results in this paper.
	
	We first fix some notations. Let $p$ be an odd prime. Let $K$ be a $p$-adic local field with ring of integers $\cO$, a uniformizer $\varpi$ and the residue field $\F$.
	
	Let $F$ be an abelian totally real field in which $p$ splits completely. Let $\Sigma$ be a finite set of finite places of $F$ containing $\Sigma_p$, the set of all places lying above $p$. Let $\bar{\chi}: G_{F, \Sigma} \to \F^\times$ be a continuous, totally odd character and assume that $\bar{\chi}$ can be extended to a character of $G_\Q$. Let $\chi: G_{F, \Sigma} \to \cO^\times $ be a continuous character lifting $\bar{\chi}$ such that $\chi|_{G_{F_v}}$ is de Rham at each place $v|p$. Write $\omega $ for the mod $p$ cyclotomic character.
	
	Note that $\bar{T}=1+\bar{\chi}$ is a pseudo-representation of $G_{F, \Sigma}$. Let $R_F^\ps$ be the universal pseudo-deformation ring of $\bar{T} $, and let $R_F^{\ps, \chi}$ be the universal pseudo-deformation ring of the pseudo-representation $\bar{T} $ with determinant $\chi$. We say that a point of $\Spec R_F^\ps$ (or $\Spec R_F^{\ps, \chi}$) is \textit{modular} if it corresponds to a Galois representation arising from a regular algebraic cuspidal automorphic representation of $\GL_2(\A_{F})$. We say a modular point is \textit{global crystalline} (or just \textit{crystalline}) if the corresponding Galois representation is crystalline at each place $v|p$.
	
	\begin{thm} [Theorem \ref{main thm} \& Theorem \ref{main thm without det}]\label{main}
		Assume that $\bar{\chi}|_{G_{F_v}} \ne \omega$ for any $v|p$ if $p=3$.
		
		1) The Zariski closure of the set of all modular points in $\Spec R_F^{\ps, \chi}$ is the union of all irreducible components of dimension $1+2[F:\Q]$. If further $\chi$ is crystalline at each $v|p$, then we can replace modular points by crystalline points in the previous sentence.
		
		2) The Zariski closure of the set of all crystalline points in $\Spec R_F^\ps$ is the union of all irreducible components of dimension $2+2[F:\Q]$.
	\end{thm}
	

		%Here we explain the structure of $\Spec R_F^{\ps, \chi}$. We can show that for each irreducible component of $\Spec R_F^{\ps, \chi}$,  either it is of dimension $1+2[F:\Q]$, or it is contained in the reducible locus, which is of dimension at most $2$. Hence, our theorem states that if an irreducble component contains a point which is not a sum of two charcters, then modular points are Zariski dense in it. The structure of $\Spec R_F^\ps$ is similar. See Corollary \ref{structure of ps} and Proposition \ref{finite same dim}.

     Before explaining the strategy of the proof of our main theorems, we indicate some obstructions to the question in the residually reducible case. 
     
     First, the universal pseudo-deformation ring in the residually reducible case is much more complicated and mysterious than that in the residually irreducible case, or equivalently, the universal deformation ring. We do not know how to use Galois cohomology to describe the dimension of the tangent space and the minimal number of quotient relations as in \cite[Proposition 1.5.1, page 30]{Berger_2013}. In general, the universal pseudo-deformation ring is not a complete intersection ring. See \cite{Deo2022} for some calculations.
	
	Second, we have hardly any small $R=\T$ theorems in the residually reducible case so far. The patching results in \cite{Skinner_1999}, \cite{pan2022fontaine} are established only for certain special primes (called ``nice primes") under restrictive conditions, but these suffice to deduce the automorphy via soluble base change. Also see \cite{Thorne_2014} and \cite{Allen_2020} in the higher dimensional case. However, we do not know how to find such a special prime in general.
	
	In addition, although the arguments of the vanishing of the Bloch-Kato Selmer group and infinite fern are likely available in the residually reducible case for most cases (see \cite{NewtonThorne} and \cite{Hernandez:2022aa} respectively), we hope to add fewer assumptions in our context.
	
	To avoid these difficulties, we adopt a new strategy instead of B\"ockle's. Our main idea is to find sufficiently many non-ordinary regular de Rham points in the universal pseudo-deformation space. Unlike the Galois theoretic method used in \cite{Fakhruddin2022}, ours is based on the proof of the Fontaine-Mazur conjecture (or automorphy lifting theorems in general). Here we explain the main steps of the proof of Theorem \ref{main}.
	
	\textbf{Step 1.  Local-global compatibility results}
	
	Write $\T$ for the big Hecke algebra for completed cohomology of quaternionic forms and let $\km$ be a maximal ideal of $\T$. Then it is classical that there exists a pseudo-representation $T_\km: G_{F, \Sigma} \to  \T_{\km}$ sending $\Frob_v$ to the Hecke operator $T_v$ for $v \notin \Sigma$. Write $\overline{T}_\km: G_{F, \Sigma} \to  \T_{\km}/\km = \F$ for the residual pseudo-representation. For each place $v|p$, let $R_v^\ps$ be the universal pseudo-deformation ring of the pseudo-representation $\overline{T}_\km|_{G_{F_v}}$ for $G_{F_v}$ with a fixed determinant. Recall that by the local-global compatibility arguments, we can show that the natural map $\widehat{\otimes}_{v|p} R_v^\ps \to \T_{\km}$ induced by the universal property is finite (see \cite[Corollary 3.5.8]{pan2022fontaine} and \cite[Corollary 3.3.7]{Zhang:2024aa}). In this part, we will prove a stronger finiteness result.
	
	Consider the map $R_v^\ps \to \T_{\km} $ and let $I_v$ be its kernel. Recall that the local-global compatibility argument (see \cite[Theorem 3.5.3]{pan2022fontaine} and \cite[Theorem 3.3.3]{Zhang:2024aa}) shows that the two actions of $R_v^\ps$ on the Hecke module $\mm_v$, the Galois side $\tau_{\Gal}$ and the automorphic side $\tau_{\mathrm{Aut}}$, are actually the same. Reviewing this proof, we can also conclude that the ideal $I_v$ is the intersection of all \textit{global regular de Rham} primes. Here for a global regular de Rham prime, we mean that it is the inverse image of a maximal ideal $\kp$ of $ \T_{\km}[1/p]$ via the composite map $ R_v^\ps \to \T_{\km} \to \T_{\km}[1/p]/\kp$, satisfying that the corresponding Galois representation $\rho(\kp)|_{G_{F_v}}$ is de Rham of distinct Hodge-Tate weights. We emphasize that a global regular de Rham prime of $R_v^\ps$ not only carries the information of $p$-adic Hodge theoretic type, but also comes from a restriction of a global Galois representation.
	
	Write $R_v :=R_v^\ps/I_v $. The ring $R_v$ was first studied by Pašk\=unas in \cite{Paskunas_2021}. (Here we point out that $R_v$ is not Kisin's potential crystalline/semi-stable deformation ring, but it might be interesting to compare $R_v$ with those objects defined in \cite{Kisin2008}. See Remark \ref{rem} for more discussions.)  If $\overline{T}_\km|_{G_{F_v}} \ne \eta+\eta\omega$ for any character $\eta$ when $p=3$, Pašk\=unas showed that the Krull dimension of $R_v$ is at most $3$ (see \cite[Corollary 5.13]{Paskunas_2021}). Note that the map $\widehat{\otimes}_{v|p} R_v^\ps \to \T_{\km}$ factors through $\widehat{\otimes}_{v|p} R_v$. Then we get a finite map $\widehat{\otimes}_{v|p} R_v \to \T_{\km}$, which also implies that the Krull dimension of $\T_{\km} $ is bounded by $\dim \widehat{\otimes}_{v|p} R_v \le 1+2[F:\Q]$. Combining \cite[Theorem 3.6.1]{pan2022fontaine}, we obtain the following byproduct, which may give a positive answer to the question in \cite[page 255-256]{CE12} in our setting.
	
	\begin{thm}[Theorem \ref{dim of T}]\label{main T}
		Assume that $\overline{T}_\km|_{G_{F_v}} \ne \eta+\eta\omega$ for any character $\eta$ for any $v|p$ when $p=3$. Then the big Hecke algebra $\T_{\km} $ is equidimensional of dimension $1+2[F:\Q] $. In particular, $\T$ is equidimensional of Krull dimension $1+2[F:\Q] $ when $p \ge 5$. 
	\end{thm}
	
    \textbf{Step 2. Potential pro-modularity}
    
    We say that a prime of $ R_F^{\ps, \chi}$ is \textit{pro-modular} if it comes from a prime of some big Hecke algebra $\T_{\km} $ via the natural map $R_F^{\ps, \chi} \twoheadrightarrow \T_{\km} $. As we mentioned before, the patching arguments in the residually reducible case are established only for some nice primes, which are difficult to find in general. In the author's previous work \cite{Zhang:2024ab} and \cite{Zhang:2024aa},  we show that there is a practicable way to find a nice prime in each irreducible component of $\Spec R_F^{\ps, \chi}$ of dimension at least $1+2[F:\Q]$ when some assumptions on $F, \Sigma$ and $\chi$ hold. Combining Pan's patching argument \cite[Theorem 4.1.4]{pan2022fontaine} (also \cite[Theorem 4.1.4]{Zhang:2024aa}) and Grunwald-Wang's theorem, for the totally real field $F$, we can find a finite abelian extension $F^1/F$ such that $p$ splits completely in $F^1$ and for any irreducible component of $\Spec R_{F^1}^{\ps, \chi} $ of dimension at least $1+2[F^1: \Q]$, it is pro-modular. Write $\Spec R_{F^1, 0}^{\ps, \chi} $ for the union of all such irreducible components with reduced scheme structure. Then by \cite[Theorem 3.6.1]{pan2022fontaine} and \cite[Corollary 3.3.10]{Zhang:2024aa}, we are able to prove a \textit{potential pro-modularity} argument, i.e., $\Spec R_{F^1, 0}^{\ps, \chi} = \Spec \T_{\km} $ for the number field $F^1$.
    
    Combining the argument in \textbf{Step 1}, for the number field $F^1$, we can show that the composite map $ \widehat{\otimes}_{v|p} R_v^\ps \to R_{F^1}^{\ps, \chi} \twoheadrightarrow R_{F^1, 0}^{\ps, \chi}$ factors through $\widehat{\otimes}_{v|p} R_v $. In other words, there exists a well-defined finite map $\widehat{\otimes}_{v|p} R_v \to R_{F^1, 0}^{\ps, \chi} $. We emphasize that it is unclear that the map $R_v \to R_{F^1, 0}^{\ps, \chi}$ is well-defined without our potential pro-modularity argument.
	
	\textbf{Step 3. A non-ordinary finiteness result}
	
	Recall that in the ordinary case, Pan showed that there exists a finite map from the Iwasawa algebra $\Lambda_F$ to the universal ordinary pseudo-deformation ring $R^{\ps, \ord}$ (see \cite[Theorem 5.1.1]{pan2022fontaine} and \cite[Theorem 6.1.1]{pan2022fontaine} for the precise statements). In this step, we prove a non-ordinary analogue to Pan's finiteness results. More precisely, we show that there exists a finite map $\widehat{\otimes}_{v|p} R_v \to R_{F, 0}^{\ps, \chi}$, where $R_{F, 0}^{\ps, \chi} $ is the maximal reduced quotient of $R_{F}^{\ps, \chi}$ such that each irreducible component of it has dimension at least $1+2[F:\Q]$.  
	
	Note that we have got a well-defined finite map $\widehat{\otimes}_{v|p} R_v \to R_{F^1, 0}^{\ps, \chi} $ for the number field $F^1$. This implies that there also exists a well-defined finite map $\widehat{\otimes}_{v \in \Sigma_p} R_v \to R_{F, 0}^{\ps, \chi} $, if we argue in a similar way as the proof of \cite[Theorem 5.1.1]{pan2022fontaine}. Here the ring $\widehat{\otimes}_{v \in \Sigma_p} R_v$ plays a similar role as the Iwasawa algebra $\Lambda_F$ in the ordinary case and it may depend on the choice of $F^1$.
	
	As the Krull dimension of $\widehat{\otimes}_{v|p} R_v$ is at most $1+2[F:\Q]$, our non-ordinary finiteness result gives an upper bound of the Krull dimension of the universal pseudo-deformation ring.
	
	\begin{thm}[Proposition \ref{finite same dim}]
		Assume that $\bar{\chi}|_{G_{F_v}} \ne \omega$ for any $v \in \Sigma_p$ if $p=3$.
		
		1) The ring $R_{F, 0}^{\ps, \chi}$ is equidimensional of dimension $1+2[F:\Q]$. In particular, $\dim R_{F}^{\ps, \chi} =1+2[F:\Q] $.
		
		2) We have $\dim R_{F}^{\ps} =2+2[F:\Q]$.
	\end{thm}
	
	\begin{rem}
		We can further show that for each minimal prime of $R_{F, 0}^{\ps, \chi}$, it is of characteristic zero and non-ordinary at each $v|p$. See Lemma \ref{char 0} and Lemma \ref{min is non-ord}.
		
		In addition, we can show that for each irreducible component of $\Spec R_{F}^{\ps, \chi} $, either it is of dimension $1+2[F:\Q]$ or it is contained in its reducible locus, which is of dimension at most $2$ by Leopoldt's conjecture. This can be proved by the classical Galois deformation theory and a trick of connectedness dimension (see \cite[Chapter 19]{localcoh} for details). In some cases (see \cite{Deo2022} and \cite{Deo2023}), $R_{F}^{\ps, \chi} $ is equidimensional, but such a result is unknown in general (at least to the author).
	\end{rem}
	
	As the Krull dimension of the big Hecke algebra is bounded by the Krull dimension of the universal pseudo-deformation ring, our result also implies Emerton's conjecture \cite[Conjecture 2.9]{Emerton2011} in most cases. (See Subsection \ref{emerton's conj} for the definition of Emerton's big Hecke algebra $\T$.)
	
	\begin{thm}[Theorem \ref{main thm for emerton}]
		The big Hecke algebra $\T$ (in the sense of Emerton) is equidimensional of Krull dimension $4$ when $p \ge 5$.
	\end{thm}
	
	\begin{rem}
		1) Unlike Theorem \ref{dim of T}, the lower bound of $\dim \T$ in this result is proved by Gouv\^ea-Mazur's infinite fern argument rather than the local-global compatibility results.
		
		2) Actually, our arguments in these three steps are also valid in the residually irreducible case. In particular, the potential pro-modularity argument based on \cite[Section 9]{pan2022fontaine} allows us to relax the Taylor-Wiles hypothesis in the residually irreducible case. See Appendix \ref{A} for the discussions.
	\end{rem}
	
	\textbf{Step 4. Going-up}
	
	Note that we have obtained a well-defined finite map $\widehat{\otimes}_{v \in \Sigma_p} R_v \to R_{F, 0}^{\ps, \chi} $, and both have the same Krull dimension. From the local-global compatibility results in \textbf{Step 1}, global non-ordinary regular de Rham primes of $\widehat{\otimes}_{v \in \Sigma_p} R_v$ are Zariski dense. Here for a global non-ordinary regular de Rham prime of $\widehat{\otimes}_{v \in \Sigma_p} R_v$, we mean that its inverse image via the composite map $R_v^\ps \twoheadrightarrow R_v \to \widehat{\otimes}_{v \in \Sigma_p} R_v $ is a global non-ordinary regular de Rham prime for each $v \in \Sigma_p $. By the going-up property, we can prove the Zariski density of non-ordinary regular de Rham liftings in $\Spec R_{F, 0}^{\ps, \chi}$, which implies our main result by using the Brauer-Nesbitt theorem and the non-ordinary Fontaine-Mazur conjecture (see \cite[Theorem 7.1.1]{pan2022fontaine} and \cite[Theorem 5.1.1]{Zhang:2024aa}).
	
	\begin{rem}
		In \cite{Allen_2014}, Calegari and Allen proved a finiteness result between the local framed deformation ring at $p$ and the global universal deformation ring in the residually irreducible case. Such a result is used in Allen's work \cite{Allen_2019}, provided that the local framed deformation ring at $p$ is regular. Following Calegari-Allen's method, we also show a finiteness result from the local pseudo-deformation ring at $p$ and the global pseudo-deformation ring in the residually reducible case with a wider range of applications than our finiteness result in \textbf{Step 3}. See Proposition \ref{finiteness1}. However, as the local pseudo-deformation ring at $p$ has Krull dimension $1+3[F:\Q] > \dim R_{F, 0}^{\ps, \chi}$, we cannot use it for \textbf{Step 4}.
		
		From our strategy, it is not necessary to know more properties of $R_v$ and $R_F^{\ps, \chi}$. Some bounds of their Krull dimensions are sufficient for our purpose in  \textbf{Step 4}. Hence, we do not need more local or global conditions  in our setup.
		
		In addition, we do not use the arguments of infinite fern and the vanishing of Bloch-Kato Selmer group in our four steps. 
	\end{rem}
	
	As a direct consequence of Theorem \ref{main}, we can prove some big $R=\T$ theorem in the residually reducible case. Here we give a sample result as an analogue to \cite[Theorem 1.2.3]{emerton11} in the Eisenstein case.
	
	\begin{thm}[Theorem \ref{big R=T modular curve}]
		Assume that $\bar{\chi}|_{G_{\Q_p}} \ne \omega$ if $p=3$. For a prime $\kp$ of $R_\Q^\ps$, if $\kp$ is contained in an irreducible component of Krull dimension $4$, then it is pro-modular, i.e., it comes from a prime of the big Hecke algebra of the completed cohomology of modular curves.
	\end{thm}
	
	

    \begin{rem}
    	Here we explain the assumptions in our setting.
    	
    	The most crucial one is that $p$ splits completely in $F$. This is because Pašk\=unas theory ($p$-adic Langlands correspondence) is only established for $\GL_2(\Q_p)$. In our strategy of the proof, Pašk\=unas' results in \cite{Pa_k_nas_2013}, \cite{Pa_k_nas_2021} and \cite{Paskunas_2021} are used in each step.
    	
    	The assumption that $\bar{\chi}|_{G_{F_v}} \ne \omega$ for any $v|p$ if $p=3$ is only prepared for the upper bound of $\dim R_v$, coming from \cite[Corollary 5.13]{Paskunas_2021}. The proof of Pašk\=unas' result relies on the explicit description of the ring $E_\kB$ for the block $\kB$ calculated in \cite{Pa_k_nas_2013} (see Subsection \ref{pas sec} for precise definitions of $E_\kB$ and $\kB$ ). In the exceptional case, the behavior of the block $\kB$ is quite different (see \cite[Appendix A.2]{Zhang:2024aa} for the list of $\kB$) so that the computation of $E_\kB$ is much more involved. Once Pašk\=unas' result can be proved in that case, the assumption can be removed directly based on the author's previous work \cite{Zhang:2024aa}.
    	
    	The other assumptions on $F$ and $\bar{\chi}$ come from the known case of Fontaine-Mazur conjecture for $\GL_2$, particularly the work \cite{Skinner_1999} and \cite{pan2022fontaine}. Also, it is sometimes convenient to use Leopoldt's conjecture for abelian totally real fields.
    	
    	In addition, our strategy is also valid in the residually irreducible case. See some discussions in Appendix \ref{A}. If one can prove analogous finiteness results to \cite[Theorem 5.1.1 \& Theorem 6.1.1]{pan2022fontaine} in the residually irreducible case, following the strategy in this paper, we may obtain similar results without Taylor-Wiles hypothesis. See Remark \ref{rem irr}.
    \end{rem}

   The author expects to see more arithmetic applications of the main results in this paper, especially the residually reducible analogue to \cite{Nakamura2023} and \cite{Fouquet:2021aa} in Iwasawa theory.

    This paper is organized as follows. In Section 2, we prepare some results for the universal pseudo-deformation rings. In Section 3, we recall Pašk\=unas theory in \cite{Pa_k_nas_2013}, \cite{Pa_k_nas_2021} and \cite{Paskunas_2021}, and give some corollaries based on the local-global compatibility results at $p$. In Section 4, we recall the known cases of Fontaine-Mazur conjecture for $\GL_2$ and potential pro-modularity arguments in \cite{Zhang:2024ab} and \cite{Zhang:2024aa}. In Section 5, we give some non-ordinary finiteness results and further prove our main theorem. In Section 6, we present some applications of our results, such as the equidimensionality of some big Hecke algebras and big $R=\T $ theorems. Furthermore, we conjecture the Zariski density of modular points in some irreducible components of the universal deformation ring in the residually reducible case there. In Appendix A, we discuss \textbf{Step 2} and \textbf{Step 3} in the residually irreducible case, which is prepared for the proof of Emerton's conjecture.
    
	
	\subsection{Notations}
	
		Let $p$ be a prime. We always denote a $p$-adic local field by $K$. We also denote its uniformizer by $\varpi$, its ring of integers by $\cO$ and the residue field by $\F$. We fix an embedding of $K$ into $\overbar{\Q_p}$, some algebraic closure of $\Q_p$ and an isomorphism $\iota_p:\overbar{\Q_p}\simeq \bC$.
	
	We use $\mathcal{C}_{\mathcal{O}}$ to denote the category of complete noetherian local $\cO$-algebras with residue field $\F$. For an $\cO$-module $M$, we use $M^\vee$ to denote its Pontryagin dual $\Hom_{\cO}^{\cont}(M, K/\cO)$, and we use $M^d$ to denote its Schikhof dual $\Hom_{\cO}^{\cont}(M, \cO)$. 
	
	For a prime ideal $\kp$ of a commutative ring $R$, we denote its residue field by $k(\kp)$. For the dimension of the prime $\kp$, we mean the Krull dimension of $R/\kp$. Let $R_\kp$ be the localization at $\kp$. We write $\widehat{R_\kp}$ for its $\kp$-adic completion. We say a commutative ring $R$ is a CNL ring if it is a complete noetherian local ring, and we usually denote its maximal ideal by $\mathfrak{m}_R$ unless otherwise stated. We use $R^{\operatorname{red}}$ to denote the maximal reduced quotient of the ring $R$.
	
	Suppose $F$ is a number field with maximal order $O_F$. We write $F_+$ for the set of totally positive elements in $F$. For any finite place $v$, we write $F_v$ (resp. $\cO_{F_v}$) for the completion of $F$ (resp. $\cO_F$) at $v$, $\varpi_v$ for a uniformizer of $F_v$, $k(v)$ for the residue field, $\operatorname{Nm}(v)$ for the norm of $v$ (in $\Q$),  $G_{F_v}$ for the decomposition group above $v$, $I_{F_v}$ (or just $I_v$ if there is no confusion for the number field $F$) for its inertia group and $\Frob_v$ for a geometric Frobenius element in $G_{F_v}/I_{F_v}$. If $l$ is a rational prime, then we denote $O_F\otimes_\Z \Z_l$ by $\cO_{F,l}$. The adele ring of $F$ will be denoted by $\A_F$ and $\A_F^\infty$ will be the ring of finite adeles. Suppose $\Sigma$ is a finite set of places of $F$. We use $G_{F, \Sigma}$ for the Galois group of the maximal extension of $F$ unramified outside $\Sigma$ and all infinite places. The absolute Galois group of $F$ is denoted by $G_F=\Gal(\overbar{F}/F)$. 
	
	We use $\varepsilon$ to denote the $p$-adic cyclotomic character and $\omega$ to denote the mod $p$ cyclotomic character. Our convention for the Hodge-Tate weight of $\varepsilon$ is $-1$. 
	
	
	\subsection{Acknowledgements}
	The author would like to thank his supervisor Professor Takeshi Saito for his constant help and encouragement. 
	
	\section{Galois deformation theory}
	
	In this section, we collect some results about Galois deformation theory.
	
	\subsection{Preliminaries in commutative algebra}
	
	In this subsection, let $p$ be an odd prime and $K$ be a $p$-adic local field with its ring of integers $\mathcal{O}$, a uniformizer $\varpi$ and its residue field $\mathbb{F}$.
	
	\begin{defn} \cite[Remark 19.2.5]{localcoh}
		For a positive integer $r$, denote by $S(r)$ the set of all ordered pairs $(A,B)$ of non-empty subsets of $\{1, ..., r\}$ for which $A \amalg B=\{1, ..., r\}$. Let $R$ be a Noetherian ring with minimal primes $\mathfrak{p}_1, ..., \mathfrak{p}_r$. The \textit{connectedness dimension} $c(R)$ of $R$ is defined to be $$
		c(R) = \min_{(A,B) \in S(r)} \{\dim R/((\cap_{i \in A} \mathfrak{p}_i)+(\cap_{j \in B} \mathfrak{p}_j))\}.$$ 
	\end{defn}
	
	It is easy to see that the connectedness dimension can also be defined as $$c(R)=\min_{Z_1,Z_2}\{\dim (Z_1\cap Z_2)\}, $$
	where $Z_1,Z_2$ are non-empty unions of irreducible components of $\Spec R$ such that $Z_1\cup Z_2=\Spec R$.
	
	\begin{prop}\label{connected dim}
		 Let $(R, \mathfrak{m}_R) $ be a Noetherian local ring, and let $\widehat{R}$ be its $ \mathfrak{m}_R$-adic completion. Then we have $c(R) \ge c(\widehat{R})$.
	\end{prop}
	
	\begin{proof}
		See \cite[Lemma 19.3.1]{localcoh}.
	\end{proof}
	
	Write $\mathcal{C}_{\mathcal{O}}$ for the category of CNL $\mathcal{O}$-algebras with residue field $\mathbb{F}$. The following result collects some useful facts about the completed tensor product of finitely many objects.
	
	\begin{prop}\label{com alg}
		Let $n $ be a positive integer. For any $1 \le i \le n $, assume that $R_i \in \mathcal{C}_{\mathcal{O}}$ is $\cO$-flat. (Here we say that an ideal $I$ of $R_i$ is $\cO$-flat if the ring $R_i/I$ is $\cO$-flat.)
		
		1) The completed tensor product ${\widehat{\otimes}_\cO} R_i$ is $\cO$-flat.
		
		2) If all $R_i $ are $\cO$-flat domains, $1 \le i \le n$, then $\dim {\widehat{\otimes}_\cO} R_i = 1+ \sum_{i=1}^{n} (\dim R_i -1)$.
		
		3) For any ideal $\kp_i$ of $R_i $ such that $R_i/\kp_i$ is $\cO$-flat, we have $$ 
		\widehat{\otimes}_\cO (R_i/\kp_i) =(\widehat{\otimes}_\cO R_i)/(\kp_1, ..., \kp_n). $$
		
		4) For each $i$, assume that $R_i$ is a domain of dimension $1$.  Then ${\widehat{\otimes}_\cO} R_i $ is reduced.
		
		5) Let $\{\kq_i\}$ be a set (possibly infinite) of $\cO$-flat primes of $R_1$. Then we have $(\cap \kq_i) \widehat{\otimes}_\cO R_2 = \cap (\kq_i \widehat{\otimes}_\cO R_2)$.
	\end{prop}
	
	\begin{proof}
		The first assertion follows from \cite[Lemma 1.3]{Thorne_2014}. The second one follows from 1) and \cite[Lemma 1.4]{Thorne_2014}. The third one follows from 1) and \cite[Lemma 4.2.4]{pan2022fontaine}.
		
		Here we prove the fourth one. As $R_i$ is an $\cO$-flat CNL domain of dimension $1$, it is finite over $\cO$. The natural map $R_i \to R_i \otimes_\cO K$ is an injection, where $K$ is the fraction field of $\cO$. Hence, $R_i \otimes_\cO K $ is a finite separable extension of $K$.
		
		By the first assertion, to show our claim, we only need to show that $({\widehat{\otimes}_\cO} R_i) \otimes_\cO K = ({\otimes}_\cO R_i) \otimes_\cO K = \otimes_\cO (R_i \otimes_\cO K)$ is reduced. As each $R_i \otimes_\cO K $ is finite separable over $K$, $\otimes_\cO (R_i \otimes_\cO K)$ is a finite \'etale $K$-algebra, and hence a product of fields of characteristic $0$. This shows the result.
		
		%As $R_i$ is an $\cO$-flat CNL domain of dimension $1$, we know that $R_i$ is actually a complete DVR and finite over $ \cO$ (for example, see \cite[28.P, page 217]{matsumura80}). We may assume $R_i=\cO[[X_i]]/{f_i(X_i)}$, where $f_i$ is a distinguished polynomial with distinct roots in $\overline{\cO}$. Then by 3), $R:= {_{i}\widehat{\otimes}_\cO} R_i = \cO[[X_1,..., X_n]]/(f_1(X_1) ,..., f_n(X_n))$. Consider an $\cO$-flat complete DVR $\cO'$ containing all the roots of $f_i$. By 1), the natural map $R \to R \otimes_\cO \cO'$ is an injection, and hence we just need to show that $R \otimes_\cO \cO'$ is reduced.
		
		%Note that $R \otimes_\cO \cO' = \cO'[[X_1,..., X_n]]/(f_1(X_1) ,..., f_n(X_n))$. If $x_i$ is a root of $f_i$, then $\mathfrak{p}:=(X_1-x_1, ..., X_n-x_n)$ is a minimal prime of $R \otimes_\cO \cO'$. As the ideal $(f_1(X_1) ,..., f_n(X_n)) $ is the intersection of all such $\kp$, we conclude that it is a radical ideal. 
		
		For the last one, we write $A_i= R_1/\kq_i$ with maximal ideal $\km_i$.  As for any integer $k>0$, both $A_i/\km_i^k$ and $R_2/\km_{R_2}^k$ are of finite length (hence finitely presented), and by \cite[Lemma 10.89.3]{stacks-project}, we have $(\prod A_i/\km_i^k) \otimes_{\cO} R_2/\km_{R_2}^k \cong \prod (A_i/\km_i^k \otimes_{\cO} R_2/\km_{R_2}^k) $. Taking the inverse limit, we conclude that $(\prod A_i)\widehat{\otimes}_\cO R_2 = \prod(A_i \widehat{\otimes}_\cO R_2)$ (using \cite[Lemma 4.14.10]{stacks-project}). Consider the exact sequence
		$$ 0 \to \cap \kq_i \to R_1 \to \prod A_i,$$
		and we tensor it with $R_1 \widehat{\otimes}_\cO R_2 $ over $R_1$.
		Combining the third assertion and \cite[Lemma 4.2.4 (1) \& (2)]{pan2022fontaine}, we get the exact sequence
		$$ 0 \to \cap (\kq_i) \widehat{\otimes}_\cO R_2 \to R_1 \widehat{\otimes}_\cO R_2 \to \prod(A_i \widehat{\otimes}_\cO R_2).$$
		This shows our result.
	\end{proof}
	
	\subsection{Pseudo-representations}\label{sec for pseudo def}
	
	In this subsection, we recall the definition of two-dimensional pseudo-representations. For a general theory, one can see \cite{bc09}.
	
	\begin{defn} \label{pseudo}
		For a profinite group $G$ and a topological commutative ring $R$ in which $2$ is invertible, a two-dimensional \textit{pseudo-representation} is a continuous function $T: G \to R$ such that 
		
		1) $T(1)=2$,
		
		2) $T(\sigma \tau)=T(\tau \sigma)$ for all $ \sigma, \tau \in G$,
		
		3) $T(\gamma \delta \eta)+T(\gamma \eta \delta)-T(\gamma \eta)T(\delta)-T(\eta \delta)T(\gamma)-T(\delta \gamma)T(\eta)+T(\gamma)T(\delta)T(\eta)=0,$
		for any $\delta, \gamma,\eta\in G$.
		
		The \textit{determinant} $\operatorname{det}(T)$ of a pseudo-representation $T$ is a character of $G$ defined by $$ \operatorname{det}(T): G \to R^\times, ~
		\operatorname{det}(T)(\sigma)=\frac{1}{2}(T(\sigma)^2-T(\sigma^2)), \sigma \in G.$$
		
	\end{defn}
	
	Sometimes, we may take $G$ as some Galois group of a totally real number field. In this case, pseudo-representations will have more properties as a complex conjugation exists in $G$. We make the following assumption:
	
	\begin{assumption}\label{odd ass}
		There exists an order $2$ element $c \in G$ such that $T(c)=0$.
	\end{assumption}
		
	\begin{defn}\label{a,d}
		Suppose Assumption \ref{odd ass} holds for $G$. For a pseudo-representation $T$ and $\sigma, \tau \in G$,  we define:
		
		1) $a(\sigma)=\frac{1}{2}(T(c\sigma)+T(\sigma))$,
		
		2) $d(\sigma)=\frac{1}{2}(-T(c\sigma)+T(\sigma))$,
		
		3) $y(\sigma, \tau)=a(\sigma \tau)-a(\sigma)a(\tau).$
	\end{defn}
	
	We can check that this definition is equivalent to the one in \cite[Section 2.4]{Skinner_1999}.
	
	\begin{defn}\label{construction}\cite[2.1.5]{pan2022fontaine}
		Suppose Assumption \ref{odd ass} holds for $G$. Assume that $R$  is either a field or a DVR. For a pseudo-representation $T: G\to R$, a representation $\rho$ associated to $T$ is in the following form.
		
		1) $ \rho(\sigma)=
		\begin{pmatrix}
			a(\sigma) & ~\\
			~& d(\sigma)
		\end{pmatrix}$, if all $y(\sigma, \tau)=0$. We call this case \textit{reducible}.
		
		2) If $y(\sigma, \tau) \ne 0 $ for some $ \sigma, \tau \in G$, choose $\sigma_0, \tau_0$ such that $\frac{y(\sigma, \tau)}{y(\sigma_0, \tau_0)} \in R$ for any $ \sigma, \tau \in G$. Define $ \rho(\sigma)=
		\begin{pmatrix}
			a(\sigma) & \frac{y(\sigma, \tau_0)}{y(\sigma_0, \tau_0)}\\
			y(\sigma_0, \sigma)& d(\sigma)
		\end{pmatrix}$.  We call this case \textit{irreducible}.
	\end{defn}
	
	Assume the profinite group $G$ satisfies Mazur's finiteness condition $\Phi_p$ (see \cite[2.1.2 Assumption 1]{pan2022fontaine}). Let $\mathcal{O}$ be a complete DVR of characteristic $0$ with residue field $\mathbb{F}$ of characteristic $p>2$. Let $T_\F: G \to \mathbb{F}$ be a two-dimensional pseudo-representation. In our case, the functor sending each object $R$ of $\mathcal{C}_{\mathcal{O}}$ to the set of pseudo-representations $T: G \to R$ which lift $T_\F$ is pro-represented by a CNL $\cO$-algebra $R^{\ps}_{T_\F}$. This is the universal pseudo-deformation ring of the residual pseudo-representation $T_\F$.
	
	\subsection{Local deformation rings}
	
	In this subsection, we study certain local universal pseudo-deformation rings.
	
	Let $p$ be an odd prime and let $K$ be a $p$-adic local field with its ring of integers $\mathcal{O}$, a uniformizer $\varpi$ and its residue field $\mathbb{F}$. Let $\chi :G_{\Q_p} \to \mathcal{O}^\times$ be a continuous character with its residual character $\bar{\chi}$.
	
	Let $T: G_{\Q_p} \to \F$ be a two-dimensional pseudo-representation with determinant $\bar{\chi}$. As the profinite group $ G_{\Q_p}$ satisfies Mazur's finiteness condition $\Phi_p$ (for example, see \cite[Theorem 7.5.12]{Neukirch_2008}), there exists a universal pseudo-deformation ring $R_{p, \chi}^{\ps}$ parametrizing all the pseudo-deformations of $T$ with determinant $\chi$. The ring $R_{p, \chi}^{\ps}$ is a CNL $\cO$-algebra.
	
	\begin{prop}\label{local pseudo}
		The ring $R_{p, \chi}^{\ps}$ is an $\cO$-flat integral domain of Krull dimension $4$.
	\end{prop}
	
	\begin{proof}
		If $T=\mathbf{1}+\mathbf{1}$ is the sum of two trivial characters, then by \cite[Corollary 9.13]{Pa_k_nas_2013}, $R_{p, \chi}^{\ps}$ is a power series ring over $\cO$ of relative dimension $3$. Thus, our result holds in this case.
		
		Now we assume $T$ is multiplicity-free (in the sense of \cite[Definition 1.4.1]{bc09}). By \cite[Corollary 5.11]{B_ckle_2023}, $R_{p, \chi}^{\ps}$ is an $\cO$-flat integral domain. We only need to show that $R_{p, \chi}^{\ps}$ is of Krull dimension $4$.
		
		If $T$ is reducible, then $\operatorname{Ext}_{\F[G_{\Q_p}]}^1 (\bar{\chi}, \mathbf{1})$ is non-trivial (for example, see \cite[Theorem 7.3.2]{Neukirch_2008}). We can choose a non-split extension $0 \to \mathbf{1} \to \bar{\rho} \to \bar{\chi} \to 0$. If $T$ is absolutely irreducible, it defines an absolutely irreducible Galois representation $\bar{\rho} $ of $ G_{\Q_p}$. In each case, there exists a universal deformation ring $R_{\bar{\rho}}^{\chi}$ parametrizing all deformations of $\bar{\rho}$ with determinant $\chi$. Combining \cite[Theorem 1.5]{B_ckle_2023} and \cite[Proposition 2.1]{Khare_2009}, we know that  $R_{\bar{\rho}}^{\chi}$ is $\cO$-flat of Krull dimension $4$. 
		
		Choose an $\cO'$-irreducible lifting $\rho$ of $\bar{\rho}$, where $\cO'$ is some finite normal extension of $\cO$ with fraction field $K'$. Then we have natural maps $R_{p, \chi}^{\ps} \to R_{\bar{\rho}}^{\chi} \to \cO' $, and suppose $\rho$ defines the prime $\kp$ (resp. $\kp'$) of $ R_{p, \chi}^{\ps}$ (resp. $ R_{\bar{\rho}}^{\chi}$). Consider the universal deformation ring $R_{\rho}$ of the representation $\rho \otimes K'$ with determinant $\chi$, and by \cite[Corollary 2.2.2]{pan2022fontaine}, $R_{\rho}$ is isomorphic to $ \widehat{(R_{p, \chi}^{\ps})_{\kp}}$. By \cite[Proposition 9.5]{kisin2003overconvergent}, $R_{\rho}$ is isomorphic to $\widehat{(R_{\bar{\rho}}^{\chi}[\frac{1}{p}])_{\kp'}} \otimes K' $, which is of Krull dimension $3$. As $R_{p, \chi}^{\ps}$ is an $\cO$-flat integral domain (hence irreducible), we can conclude that $R_{p, \chi}^{\ps}$ is of Krull dimension $4$.
	\end{proof}
	
	\begin{rem}
		 In this article, we are concerned with the case $T$ is reducible. If we are not in a special case $\bar{\chi}=\omega$ and $p=3$, the ring $R_{p, \chi}^{\ps}$ has an explicit description. See \cite[Corollary 9.13, B.5,  Proposition B.17 and Remark B.28]{Pa_k_nas_2013} for more details.
		
	\end{rem}
	
	\subsection{Global deformation rings}
	In this subsection, we collect some results about global pseudo-deformation rings.
	
	Let $p$ be an odd prime and let $K$ be a $p$-adic local field with its ring of integers $\mathcal{O}$, a uniformizer $\varpi$ and its residue field $\mathbb{F}$. 
	
	Let $F$ be an abelian totally real field. For simplicity, we suppose that $p$ splits completely in $F$ (but it may not be necessary in this subsection). Let $\Sigma_p$ be the set of places of $F$ lying above $p$ and let $\Sigma$ be a finite set of finite places of $F$ containing $\Sigma_p$. Let $\chi :G_{F, \Sigma} \to \mathcal{O}^\times$ be a continuous totally odd character with its residual character $\bar{\chi}$.
	
	Let $T$ be the pseudo-representation $1+\bar{\chi}$. As the profinite group $G_{F, \Sigma} $ satisfies Assumption \ref{odd ass} and Mazur's finiteness assumption (see \cite[Corollary 10.11.15]{Neukirch_2008}), there exists a universal pseudo-deformation $R_{F}^{\ps}$ parametrizing all the pseudo-deformations of $T$. Let $R_{\bar{\chi}}$ be the universal deformation ring parametrizing all the liftings of $\bar{\chi}$. Write $ R_{F}^{\ps, \chi} = R_{F}^{\ps} \widehat{\otimes}_{R_{\bar{\chi}}, \chi} \cO$, where the map $R_{\bar{\chi}} \to \cO $ is given by $\chi$. Then $ R_{F}^{\ps, \chi}$ is the universal pseudo-deformation ring parametrizing all the pseudo-deformations of $T$ with determinant $\chi$.
		
    \begin{prop}\label{det}
    	1) Let $\Gamma$ be the maximal pro-$p$ abelian quotient of $G_{F, \Sigma} $. Then there exists a CNL $\cO$-algebra isomorphism $R_{F}^{\ps} \to R_{F}^{\ps, \chi} \widehat{\otimes}_\cO \cO[[\Gamma]] $.
    	
    	2) Let $\chi': G_{F,\Sigma} \to \cO^\times$ be a character lifting $\bar{\chi}$. Let $R_{F}^{\ps, \chi'}$ be the universal pseudo-deformation ring parametrizing all the pseudo-deformations of $T$ with determinant $\chi'$.Then there exists a CNL $\cO$-algebra isomorphism $ R_{F}^{\ps, \chi} \to R_{F}^{\ps, \chi'}$.
    \end{prop}
	
	\begin{proof}
		This is an analogue to \cite[Lemma 6.1.2]{Allen_2019}. We sketch the proof here.
		
		Let $T^\chi: G_{F, \Sigma} \to R_{F}^{\ps, \chi}$ be the universal pseudo-deformation of $T$ with determinant $\chi$ and let $D: G_{F, \Sigma} \to \cO[[\Gamma]]^\times $ be the universal deformation of the trivial character $\mathbf{1}$. As $p$ is odd, there exists a unique character $\tilde{\chi}$ of $G_{F, \Sigma}$ lifting $\mathbf{1}$ satisfying $(\tilde{\chi})^2=\chi\tilde{\bar{\chi}}^{-1}$, where $ \tilde{\bar{\chi}}$ is the Teichm\"uller lifting of $ \bar{\chi}$. Then $T^\chi \otimes (\tilde{\chi}D): G_{F, \Sigma} \to R_{F}^{\ps, \chi} \widehat{\otimes}_\cO \cO[[\Gamma]]$ defines a pseudo-deformation of $T$.
		
		For a pseudo-deformation $T_A: G_{F, \Sigma} \to A$ lifting $T$, $A \in \mathcal{C}_{\mathcal{O}}$, there exists a unique character $d_A$ of $G_{F, \Sigma}$ (lifting $\mathbf{1}$) satisfying $(d_A)^2= \det(T_A)\chi^{-1}$. Then $T_A \otimes (d_A)^{-1} : G_{F, \Sigma} \to A$ defines a pseudo-deformation of $T$ with determinant $\chi$. Thus, we get natural maps $R_{F}^{\ps, \chi} \to A $ and $ \cO[[\Gamma]] \to A$ (induced by $\tilde{\chi}^{-1}d_A$) by the universal property. By the uniqueness of the universal object, we have  $R_{F}^{\ps} \cong R_{F}^{\ps, \chi} \widehat{\otimes}_\cO \cO[[\Gamma]] $. This shows the first claim.
		
		For the second one, we may consider the pseudo-representation $ T^\chi \otimes \mu : G_{F, \Sigma} \to R_{F}^{\ps, \chi}$, where $\mu : G_{F, \Sigma} \to \cO^\times$ is the unique character lifting $\mathbf{1}$ satisfying $\mu^2 = \chi'\chi^{-1}$. This defines a pseudo-deformation of $T$ with determinant $\chi'$, and we get the map $ R_{F}^{\ps, \chi'} \to R_{F}^{\ps, \chi}$ by the universal property. Similarly, we can also obtain the map $ R_{F}^{\ps, \chi} \to R_{F}^{\ps, \chi'}$. The composite maps $R_{F}^{\ps, \chi'} \to R_{F}^{\ps, \chi} \to R_{F}^{\ps, \chi'} $ and $R_{F}^{\ps, \chi} \to R_{F}^{\ps, \chi'} \to R_{F}^{\ps, \chi} $ are both identities. This shows the second claim.
		
			\end{proof}
	
	The following result is crucial to our main results.
	
	\begin{prop}\label{irreducible point}
		Let $\kp$ be a one-dimensional prime of $ R_{F}^{\ps, \chi}$ such that the corresponding Galois representation is absolutely irreducible in the sense of Definition \ref{construction}. Here we do not require $\varpi \notin \kp$.
		
	   1)	 We have $c((R_{F}^{\ps, \chi})_{\kp}) \ge 2[F: \Q]-1$.
		
		2) For any irreducible component $C$ of $ \Spec R_{F}^{\ps, \chi}$ containing $\kp$, we have $\dim C \ge 2[F:\Q]+1$.
	\end{prop}
	
	\begin{proof}
		For the first assertion, it is proved in \cite[Lemma 7.4.6]{pan2022fontaine}. Here we only point out that the inequality $c((R_{F}^{\ps, \chi})_{\kp}) \ge  c(\widehat{(R_{F}^{\ps, \chi})_{\kp}})$ follows from Proposition \ref{connected dim}.
		
		For the second assertion, it is a direct consequence of the first one. See \cite[Remark 3.2.8]{Zhang:2024ab}.
	\end{proof}
	
	\begin{cor}\label{structure of ps}
		For any irreducible component $C$ of $ \Spec R_{F}^{\ps, \chi}$, $C$ is either of dimension at least $1+2[F:\Q]$ or at most $2$.
	\end{cor}
	
	\begin{proof}
		By Proposition \ref{irreducible point}, if $C$ is of dimension less than $1+2[F:\Q]$, then $C$ only contains reducible pseudo-deformations. Let $R_{F}^{\operatorname{ps, red}}$ be the reducible locus of $ R_{F}^{\ps, \chi}$. Then there exists a natural surjection $\cO[[G_{F, \Sigma}^{\operatorname{ab}}(p)]] \twoheadrightarrow R_{F}^{\operatorname{ps, red}}$. As $F$ is abelian, by Leopoldt's conjecture, we know that $R_{F}^{\operatorname{ps, red}}$ is of dimension at most  $2$, and so is $C$. This proves the result.
	\end{proof}
	
    Assume that there exists a one-dimensional irreducible pseudo-deformation of $T$ with determinant $\chi$. Let $I_F$ be the intersection of all minimal primes of $R_{F}^{\ps, \chi} $ of dimension at least $2[F:\Q]+1 $, and write $ R_{F,0}^{\ps, \chi} := R_{F}^{\ps, \chi}/I_F$. Then from our definition, we can conclude the following properties of $ R_{F,0}^{\ps, \chi}$ directly.
    
    \begin{prop}\label{def of 0}
    	1) $ R_{F,0}^{\ps, \chi}$ is reduced. Each irreducible component of $ R_{F,0}^{\ps, \chi}$ is of Krull dimension at least $1+2[F:\Q]$.
    	
    	2) For any irreducible prime $\kp$ of $R_{F}^{\ps, \chi} $, the natural surjection $ R_{F}^{\ps, \chi} \twoheadrightarrow R_{F}^{\ps, \chi}/\kp$ factors through $ R_{F,0}^{\ps, \chi}$.
    \end{prop}
    
    Now let $F_1/F$ be a totally real extension and suppose $F_1$ is also abelian. Let $\Sigma_1$ be the set of primes of $F_1$ lying above $\Sigma$. Let $ R_{F_1}^{\ps, \chi}$ is the universal pseudo-deformation ring parametrizing all the pseudo-deformations of $T|_{G_{F_1, \Sigma_1}}$ with determinant $\chi|_{G_{F_1, \Sigma_1}}$. By the universal property, there is a natural map $ R_{F_1}^{\ps, \chi} \to R_{F}^{\ps, \chi}$.
    
    \begin{prop}\label{finiteness of ps}
      The natural map $ R_{F_1}^{\ps, \chi} \to R_{F}^{\ps, \chi}$ is finite.
    \end{prop}
	
	\begin{proof}
		Write $ \km_1$ for the maximal ideal of $R_{F_1}^{\ps, \chi}$. We only need to show that $R_{F}^{\ps, \chi}/ (\km_1)$ is an Artinian ring. If not, we can find a one-dimensional prime $\kq$ of $R_F^{\ps, \chi}/ (\km_1)$. Consider the Galois representation $\rho(\kq) : G_{F, \Sigma} \to \GL_2(k(\kq))$ in the sense of Definition \ref{construction} (here we fix $\sigma_0, \tau_0 \in G_{F, \Sigma}$).
		
		  We claim that $\rho(\kq)|_{G_{F_1, \Sigma_1}}$ is either upper triangular or lower triangular. From our construction, we have $y(\alpha, \beta)=0$, $\alpha, \beta \in G_{F_1, \Sigma_1}$. If it is not lower triangular, i.e., $y(\alpha, \tau_0) \ne 0$ for some $\alpha \in G_{F_1, \Sigma_1}$, then $0=y(\alpha, \beta)y(\sigma_0, \tau_0)=y(\alpha, \tau_0)y(\sigma_0, \beta)$ (by \cite[2.1.4]{pan2022fontaine}), which implies $ y(\sigma_0, \beta)=0$. Then by Definition \ref{construction}, $\rho(\kq)|_{G_{F_1, \Sigma_1}}$ is 
		  upper triangular.
		  
		  We assume $\rho(\kq)|_{G_{F_1, \Sigma_1}}$ is upper triangular (the other case will be similar). Note that $\bar{\chi}$ is totally odd. For any complex conjugation $c \in G_{F_1, \Sigma_1}$, by our definition, we have $ \rho(\kq)(c)=\begin{pmatrix}
			1 & *\\
			~& -1
		\end{pmatrix}$ (as $a|_{G_{F_1, \Sigma_1}}=\mathbf{1}, d|_{G_{F_1, \Sigma_1}}=\bar{\chi}|_{G_{F_1, \Sigma_1}}$ for functions $a, d$ in Definition \ref{a,d}). For any $x \in G_{F, \Sigma}$, we have $\rho(\kq)(xcx^{-1}) = \rho(\kq)(c')$, where $c, c'$ are complex conjugations of $G_{F_1, \Sigma_1} $. Then by an easy calculation, we conclude that $\rho(\kq)(x)$ is upper-triangular. This shows that $\rho(\kq) $ is reducible. Consequently, we know that $\rho(\kq)$ is a direct sum of two characters. Here we write $\rho(\kq) = \psi_1 \oplus \psi_{2}$, where $\psi_1$ is a lifting of $\mathbf{1}$. By the definition of $\kq$, we know that $\psi_1|_{G_{F_1, \Sigma_1}}=\mathbf{1}$. As $F_1/F$ is finite, we have $\psi_1=\mathbf{1}$ and $\psi_{2} =\bar{\chi}$ as $ \psi_1\psi_{2}=\bar{\chi}$. It is contrary to the assumption that $\kq$ is of dimension one. This proves our result.
	\end{proof}
	
	\begin{cor}\label{lower}
		Assume that $\bar{\chi}$ can be extended to a character of $G_\Q$. Then we have $\dim R_{F}^{\ps, \chi} \ge 1+2[F:\Q]$. In other words, $\Spec R_{F, 0}^{\ps, \chi}$ is not empty.
	\end{cor}
	
	\begin{proof}
		Note that $\bar{\chi} $ can be extended to a character of $G_{\Q}$. By the arguments in \cite[Lemma 7.3]{Fakhruddin2022}, the pseudo-deformation space $\Spec R_\Q^{\ps, \chi'}$ (of the pseudo-representation $1+\bar{\chi}$ with determinant $\chi'$) has a one-dimensional irreducible point associated to some cusp eigenform $f$ of tame level $N_0$ (the Artin conductor of the representation $1 \oplus \bar{\chi}$), where $\chi'$ is the corresponding determinant lifting $\bar{\chi}$. Then by Proposition \ref{irreducible point}, we have $\dim R_\Q^{\ps, \chi'} \ge 3$. Consider the natural map $ R_{F}^{\ps, \chi'} \to R_\Q^{\ps, \chi'}$. By Proposition \ref{finiteness of ps}, we have $\dim R_{F}^{\ps, \chi'} \ge \dim R_\Q^{\ps, \chi'} >2$, and hence $\dim R_{F}^{\ps, \chi'} \ge 1+2[F:\Q]$ by Corollary \ref{structure of ps}. Using 2) of Proposition \ref{det}, we get $\dim R_{F}^{\ps, \chi} \ge 1+2[F:\Q]$. In particular, we have $\Spec R_{F, 0}^{\ps, \chi}$ is not empty.
	\end{proof}
	
	Similarly to $ R_{F,0}^{\ps, \chi}$, we can also define $ R_{F_1,0}^{\ps, \chi}:= R_{F_1}^{\ps, \chi}/I_{F_1}$ as the union of all irreducible components of $ R_{F_1}^{\ps, \chi}$ dimension at least $1+2[F_1: \Q]$. 
	
	\begin{cor}\label{finite for 0}
		The composite map $ R_{F_1}^{\ps, \chi} \to R_{F}^{\ps, \chi} \twoheadrightarrow R_{F,0}^{\ps, \chi}$ factors through $ R_{F_1,0}^{\ps, \chi}$. In other words,  there exists a well-defined and finite map $ R_{F_1,0}^{\ps, \chi} \to R_{F,0}^{\ps, \chi}$ induced by the universal property.
	\end{cor}
	
	\begin{proof}
		For the first assertion, we need to show that the kernel $I:= \operatorname{ker} \{R_{F_1}^{\ps, \chi} \to R_{F,0}^{\ps, \chi} \}$ contains $I_{F_1} $. Combining Proposition \ref{def of 0} and Proposition \ref{finiteness of ps}, we know that each minimal prime of $I$ is of dimension at least $1+2[F: \Q] >2$. By Corollary \ref{structure of ps}, we have $I_{F_1} \subset I$. The second one follows from Proposition \ref{finiteness of ps}.
	\end{proof}
	
	\section{Local-global compatibility}
	
	In this section, we recall some local-global compatibility results and deduce some corollaries. Our references are \cite{Pa_k_nas_2013}, \cite{Paskunas_2021}, \cite{Pa_k_nas_2021} and \cite{pan2022fontaine}.
	
	Throughout this section, we fix $p$ an odd prime, and $F$ a totally real field in which $p$ splits completely. We denote by $K$  a $p$-adic local field with its ring of integers $\cO$, a uniformizer $\varpi$ and its residue field $\F$.
	
	\subsection{Completed cohomology and Hecke algebra}\label{sec for T}
	In this subsection, we recall some definitions about the completed cohomology for quaternionic forms and the $p$-adic big Hecke algebra. Our reference is \cite[Section 3]{pan2022fontaine}.
	
	Let $D$ be a quaternion algebra with centre $F$ such that $D$ is ramified at all infinite places of $F$ and unramified at all places lying above $p$. 
	
	Let $A$ be a topological $\cO$-algebra and $U = \prod_{v \nmid \infty} U_v$ be an open compact subgroup of $(D \otimes_F \mathbb{A}_F^\infty)^\times$ such that $U_v \subseteq \operatorname{GL}_2(\cO_{F_v})$ for any $v \mid p$. We write $U^p= \prod_{v \nmid p} U_v$ (the tame level) and $U_p =\prod_{v \mid p} U_v$. Let $\psi: (\mathbb{A}_F^\infty)^\times/F_+^\times \to A^\times$ be a continuous character, where $F_+$ is the set of totally positive elements in $F$.
	
    Let $S_{\psi}(U,A)$ be the space of continuous functions:
    \[f:D^{\times}\setminus \DAi\to A\]
    such that for any $g\in \DAi,u\in U,z\in \AFi$, we have
    \begin{itemize}
    	\item $f(gu)=f(g)$,
    	\item $f(gz)=\psi(z)f(g)$.
    \end{itemize}
    
     Let $\Sigma$ be a finite set of places of $F$ containing all those lying above $p$ and places $v$ where either $D$ is ramified or $U_v$ is not a maximal open subgroup. For any $v \notin \Sigma$, we define the Hecke operator $T_v \in \operatorname{End}(S_{\psi}(U,A))$ to be the double coset action $[U_v \begin{pmatrix}
    	\varpi_v & ~\\
    	~ & 1
    \end{pmatrix}U_v]$. Precisely, if we write $U_v \begin{pmatrix}
    	\varpi_v & ~\\
    	~ & 1
    \end{pmatrix}U_v  = \coprod_i \gamma_i U_v $, then $$
    (T_v \cdot f) (g) = \sum_i f(g\gamma_i), ~~f \in S_{\psi}(U,A). $$
    
    We define the Hecke algebra $\mathbb{T}_{\psi} ^\Sigma (U,A) \subseteq \operatorname{End}(S_{\psi}(U,A))$ as the $A$-subalgebra generated  by all Hecke operators $T_v, v \notin \Sigma$. In the case that $A= \mathcal{O}/ \varpi^n$, $ U$ is sufficiently small (in the sense of \cite[page 1046]{pan2022fontaine}) and $\psi \mid_{U \cap (\mathbb{A}_F^\infty)^\times} $ is trivial modulo $ \varpi^n$, we can check that the Hecke algebra $\mathbb{T}_{\psi}^\Sigma (U, \mathcal{O}/ \varpi^n)$ is independent of $\Sigma$, and we drop $\Sigma$ afterwards.
    
    For a tame level $U^p$, we write $S_{\psi}(U^p,A):=\varinjlim_{U_p} S_\psi(U^pU_p,A) $ with discrete topology, where $U_p=\prod_{v|p}U_v$ ranges over all open compact subgroups $U_v$ of $\GL_2(F_v)$. 

    
    \begin{defn} \label{Completed coh}
    	The \textit{completed cohomology} of level $U^p$ is defined to be
    	\[S_\psi(U^p):=\Hom_\cO(K/\cO,S_{\psi}(U^p,K/\cO))\]
    	equipped with $\varpi$-adic topology, and the \textit{completed homology} is defined to be
    	\[M_\psi(U^p):=S_\psi(U^p,K/\cO)^\vee=\Hom_\cO^{\cont}(S_{\psi}(U^p,K/\cO),K/\cO)\]
    	equipped with compact-open topology.
    \end{defn}
    
    The following proposition is clear from the definition.
    
    \begin{prop}\label{3.1.2}
    	1) $ S_\psi(U^p)$ is $\varpi$-torsion free.
    	
    	2) $ S_\psi(U^p) \cong \Hom^{\mathrm{cont}}_{\cO}(M_\psi(U^p),\cO)$.
    \end{prop}
    
   	\begin{defn} \label{big}
   	Let $U^p$ be a tame level and $\psi: (\mathbb{A}_F^\infty)^\times/ (U^p \cap (\mathbb{A}_F^\infty)^\times)F_+^\times \to \mathcal{O}^\times$ be a continuous character. Define the \textit{big Hecke algebra} $$
   	\mathbb{T}_{\psi}(U^p) = \varprojlim_{(n, U_p) \in \mathcal{I}} \mathbb{T}_{\psi}(U^pU_p, \mathcal{O}/\varpi^n), $$
   	where $\mathcal{I}$ is the set of pairs $(n, U_p)$ with $U_p \subseteq K_p :=\prod_{v \mid p} \operatorname{GL}_2(\cO_{F_v})$ and $n$ a positive integer such that $ \psi \mid_{U_p \cap \mathcal{O}_{F,p}^\times} \equiv 1~ \operatorname{mod}~\varpi^n$.
   	
   \end{defn}
    
    The following result is classical.
    
    \begin{prop}\label{pseudo of T}
    	1) $\mathbb{T}_{\psi}(U^p) $ is a semi-local ring. Write $\km_1,\cdots,\km_r$ for all the maximal ideals of $\T_\psi(U^p)$. Then we have $\T_\psi(U^p)\cong \T_\psi(U^p)_{\km_1}\times\cdots\times \T_\psi(U^p)_{\km_r}$ and each $\T_\psi(U^p)_{\km_i}$ is $\km_i$-adically complete and separated. 
    	
    	2) For each maximal ideal $\km_r$, there exists a two-dimensional pseudo-representation $T_{\km_i}: G_{F, \Sigma} \to \T_\psi(U^p)_{\km_i}$ sending $\Frob_v$ to $T_v$ for any $v \notin \Sigma$ with determinant $\psi\varepsilon^{-1}$.
    \end{prop}
    
    \subsection{Pašk\=unas theory}\label{pas sec}
    
    In this subsection, we summarize some basic definitions and results of Paskunas theory for the $p$-adic analytic group $D_p^\times:=\prod_{v|p} \GL_2(F_v)$. For more details and precise definitions, one can see \cite{Pa_k_nas_2013}, \cite{Pa_k_nas_2021} or a summary in \cite[Appendix A]{Zhang:2024aa} for the case $G=\GL_2(\Q_p)$. We also keep the notations and terminologies there.
    
    Let $G=\prod_{v|p} \GL_2(F_v) \cong \prod_{v|p}\GL_2(\Q_p)$ and $Z(G)\simeq \prod_{v|p} F_v^\times$ be its centre. Let $\zeta: Z(G) \to \cO^\times $ be a continuous character.
    
    Let $\operatorname{Mod}^{\mathrm{sm}}_{G, \zeta}(\mathcal{\cO})$ be the category of smooth 
    representations of $G$ on $\mathcal{O}$-torsion modules with central character $\zeta$, i.e. the centre $Z$ acts via $\zeta$. Let $\mathrm{Mod}^{\mathrm{l\,fin}}_{G, \zeta}(\cO)$ be the full subcategory of $\mathrm{Mod}_{G, \zeta}^{\mathrm{sm}}(\cO)$ consisting of representations locally of finite length and $\mathrm{Mod}^{\mathrm{l\, adm}}_{G, \zeta} (\cO)$ the full subcategory consisting of locally admissible representations. By \cite[Lemma 3.4.1]{pan2022fontaine}, $\mathrm{Mod}^{\mathrm{l\,fin}}_{G, \zeta}(\cO)$ and $\mathrm{Mod}^{\mathrm{l\, adm}}_{G, \zeta} (\cO)$ are actually the same.
    
    Let $H \subset G$ be an open compact subgroup and $\cO[[H]]$ be the completed group algebra of $H$ over $\cO$. Let $\mathrm{Mod}_{G, \zeta}^{\mathrm{pro\,aug}}(\cO)$ be the category of profinite linearly topological $\cO[[H]]$-modules with action of $\cO[G]$ and central character $\zeta $ such that both actions agree on $\cO[H]$, with morphisms given by $G$-equivariant continuous homomorphisms of topological
    $\cO[[H]]$-modules. This category is independent of the subgroup $H$. 
    
    For an object $M$ of $\operatorname{Mod}^{\mathrm{sm}}_{G, \zeta}(\mathcal{\cO})$, the Pontryagin dual $M \mapsto M^\vee$ induces an anti-equivalence between categories $\operatorname{Mod}^{\mathrm{sm}}_{G, \zeta}(\mathcal{\cO})$ and $\mathrm{Mod}_{G, \zeta}^{\mathrm{pro\,aug}}(\cO)$. Write $\kC_{G,\zeta}(\cO)$ for the full subcategory of $\mathrm{Mod}_{G, \zeta}^{\mathrm{pro\,aug}}(\cO)$ with objects isomorphic to $\pi^\vee$ for some object $\pi$ in $\mathrm{Mod}^{\mathrm{l\, adm}}_{G,\zeta}(\cO)$.
    
    Let $\mathrm{Irr}_{G,\zeta}$ be the set of irreducible representations in $\mathrm{Mod}_{G,\zeta}^{\mathrm{sm}}(\cO)$. Recall that a \textit{block} is an equivalence class of the relation $\sim$, where $\pi \sim \pi'$ if there exists $\pi_1, \cdots, \pi_n \in \Irr_{G, \zeta}$ such that $\pi \cong \pi_1$, $\pi' \cong \pi_n$, and for $1 \leq i \leq n-1$, $\pi_i \cong \pi_{i+1}$ or $\Ext^1_{G}(\pi_i, \pi_{i+1}) \neq 0$ or $\Ext^1_{G}(\pi_{i+1}, \pi_i) \neq 0$.
    Under this equivalence relation, we have a decomposition of the category $\mathrm{Mod}^{\mathrm{l\, adm}}_{G, \zeta} (\cO)$:
    \begin{eqnarray*}
    	\mathrm{Mod}^{\mathrm{l\, adm}}_{G,\zeta}(\cO)\cong \prod_{\mathfrak{B}\in\mathrm{Irr}_{G,\zeta}/\sim}\mathrm{Mod}^{\mathrm{l\, adm}}_{G,\zeta}(\cO)^\kB,
    \end{eqnarray*}
    where $\mathrm{Mod}^{\mathrm{l\, adm}}_{G,\zeta}(\cO)^\kB$ is the full subcategory of $\mathrm{Mod}^{\mathrm{l\, adm}}_{G,\zeta}(\cO)$ consisting of representations with all irreducible subquotients in the block $\kB$. Write $\kC_{G,\zeta}(\cO)^\kB$ for the full subcategory of $\kC_{G,\zeta}(\cO)$ anti-equivalent to $\mathrm{Mod}^{\mathrm{l\, adm}}_{G,\zeta}(\cO)^\kB$ for the block $\kB$ via the Pontryagin dual. Then we also have a natural decomposition \[\kC_{G,\zeta}(\cO)\cong\prod_{\mathfrak{B}\in\mathrm{Irr}_{G,\zeta}/\sim}\kC_{G,\zeta}(\cO)^\kB.\]
    
    In our case, the blocks containing an absolutely irreducible representation have been well-understood. See the list in \cite[Appendix A.2]{Zhang:2024aa} (containing the case $p=2$) for the case $G= \GL_2(\Q_p)$ and combine \cite[Lemma 3.4.7]{pan2022fontaine}. For each $v|p$ and a block $\kB$, we can attach a semi-simple two-dimensional representation $(\bar{\rho}_{\kB})_v$ over $\F$ with determinant $\varepsilon\zeta|_{F_v} \modd \varpi$. Here we view $ \zeta|_{F_v}$ as a character of $ G_{F_v}$ by local class field theory. 
    
    Now fix a block $\kB$, and write $\pi_{\kB}=\bigoplus_{\pi\in\kB_i}\pi$, where $\kB_i$ is the set of isomorphism classes of elements of $\kB$. This is a finite set by the discussion in the previous paragraph. Let $\pi_\kB\hookrightarrow J_{\kB}$ be an injective envelope of $\pi_{\kB}$ in $\mathrm{Mod}^{\mathrm{l\, adm}}_{G,\zeta}(\cO)$. Then its Pontryagin dual $P_{\kB}:=J_{\kB}^\vee$ is a projective envelope of $\pi^\vee\cong\bigoplus_{\pi\in\kB_i}\pi^\vee$ in $\kC_{G,\zeta}(\cO)$, called a \textit{projective generator} of $\kB$. Write
    $E_{\kB}:=\End_{\kC_{G,\zeta}(\cO)}(P_{\kB})\cong \End_G(J_{\kB}).$
    This is a pseudo-compact ring (not necessarily commutative). Let $Z_{\kB}$ be the centre of $ E_{\kB}$. 
    
    \begin{thm} \label{Paskunas theory for G}
    	Write $R_v^{\ps}$ for the universal pseudo-deformation ring of $(\bar{\rho}_{\kB})_{v}$ with determinant $\varepsilon\zeta|_{F_v}$ (via class field theory). 
    	
    	1) There exists a natural finite map $\widehat{\bigotimes}_{v|p} R^{\ps}_{v}\to E_{\kB} $. In particular, $E_{\kB} $ is Noetherian.
    	
    	2) $ E_{\kB}$ is a finitely generated module over $Z_{\kB}$.
    \end{thm}
    
    \begin{proof}
    	See \cite[Corollary 3.4.8]{pan2022fontaine} and \cite[Proposition 3.3.2]{Zhang:2024aa}.
    \end{proof}
    
     The functor sending an object $N$ of the category $ \kC_{G,\zeta}(\cO)$ to  $\Hom_{\kC_{G,\zeta}(\cO)}(P_\kB, N)$ induces an anti-equivalence between $ \kC_{G,\zeta}(\cO)^\kB$ and the category of right pseudo-compact $ E_{\kB}$-modules. The inverse functor is given by $\mm \mapsto \mm \widehat{\otimes}_{E_{\kB}} P_{\kB}$.
  
    Let $\Pi$ be a unitary $K$-Banach space representation of $G$ and $\Theta$ an open bounded $G$-invariant lattice in $\Pi$. By \cite[Lemma 4.4]{Pa_k_nas_2013}, its Schikhof dual
    $ \Theta^d:=\Hom_{\cO}^{\cont}(\Theta, \cO)$ is an object of $\mathrm{Mod}_{G, \zeta}^{\mathrm{pro\,aug}}(\cO)$. 
    
    %Write $ \Ban_{G, \zeta}^{\adm}(K)_{\kB}$ for the category of admissible unitary $K$-Banach space representations $\Pi$ such that the Schikhof dual of each open bounded $G$-invariant lattice of $\Pi$ is an object of $\kC_{G,\zeta}(\cO)^\kB$.$N$ is an object in 
    
    \begin{rem}
    	Fix a tame level $U^p$ and a character $\psi$ as in the previous subsection. Take $\zeta=\psi$. We can conclude the following facts:
    	
    	1) $S_\psi(U^p, K/\cO)$ is a smooth admissible $\cO$-representation of $D_p^\times$.
    	
    	2) The complete homology $M_\psi(U^p)$ is an object in $ \kC_{G,\zeta}(\cO)$. Here we view $\psi$ as a character of $Z(G)$.
    	
    	3) Write $S_\psi(U^p)_K :=S_\psi(U^p) \otimes_\cO K$. Then $S_\psi(U^p)_K$ is a $K$-Banach space representation with open bounded $G$-invariant $\cO$-lattice $S_\psi(U^p) $.
    \end{rem}
    
    \subsection{Local-global compatibility}
    In this subsection, we collect some local-global compatibility results.
    
    We fix a tame level $U^p$ and character $\psi$. Recall that $F$ is a totally real field in which $p$ splits completely. Let $\Sigma$ be a finite set of places of $F$ containing all those lying above $p$ and places $v$ where either $D$ is ramified or $U_v$ is not a maximal open subgroup. Under this setting, we can define a big Hecke algebra $\T_{\psi}(U^p)$. Let $\km$ be a maximal ideal of $\T_{\psi}(U^p)$. Then by Proposition \ref{pseudo of T}, we get a two-dimensional pseudo-representation $T_\km : G_{F,\Sigma} \to \T_{\psi}(U^p)_\km$ with determinant $\psi\varepsilon^{-1}$ and its residual representation $\bar{\rho}_\km : G_{F,\Sigma} \to \GL_2(\F)$ in the sense of 
    Definition \ref{construction}. 
    
    For each $v|p$, we write $\bar{\rho}_{\km, v} := \bar{\rho}_\km|_{G_{F_v}} \otimes \varepsilon$. Assume that $\bar{\rho}_{\km, v}$ is semisimple for each $v|p$.  We can define a block $\kB_v$ of $\GL_2(F_v)$ corresponding to $\bar{\rho}_{\km, v}$ (see the list in \cite[Appendix A.2]{Zhang:2024aa}) for each $v|p$. Then by \cite[Lemma 3.4.7]{pan2022fontaine}, $\kB_\km := \otimes_{v|p} \kB_v$ defines a block of $D_p^\times=\prod_{v|p} \GL_2(F_v)$. 
    
    Let $P_{\kB_v}$ be a projective generator of $\kB_v$ for each $v|p$. Write $E_{\kB_v}:= \End_{\kC_{\GL_2(F_v),\psi}(\cO)}(P_{\kB_v})$.   Then by \cite[Lemma 3.4.5]{pan2022fontaine}, $P_{\kB_\km} := \widehat{\otimes}_{v|p} P_{\kB_v}$ defines a projective generator of the block $\kB_\km$. Also, by \cite[Proposition 3.3.2]{Zhang:2024aa}, we have $ E_{\kB_\km}:=\End_{\kC_{D_p^\times,\psi}(\cO)}(P_{\kB_\km}) \cong  \widehat{\otimes}_{v|p} E_{\kB_v}$.
    
     Write $$\mm_v := \Hom_{\kC_{\GL_2(F_v),\psi}(\cO)}(P_{\kB_v}, M_\psi(U^p)_\km),~~~ v|p,$$
    $$\mm:=\Hom_{\kC_{D_p^\times,\psi}(\cO)}(P_{\kB_\km}, M_\psi(U^p)_\km).$$
     
     Write $R_v^\ps$ for the universal pseudo-deformation ring of the  pseudo-representation $\tr \bar{\rho}_{\km, v}$ with determinant $\psi\varepsilon^{-1}$, and write $R_{F,p}^\ps := \widehat{\otimes}_{v|p} R_v^\ps$. Then by the universal property, there is a natural map $ R_{F,p}^\ps \to \T_{\psi}(U^p)_\km$ induced by $T_\km$.
     
     We can define two actions of $R_{F,p}^\ps $ on $\mm$. 
     \begin{enumerate}
     	\item \textbf{the Galois side} $\tau_{\Gal}$: which comes from the action of $\T_\psi(U^p)_\km$ on $M_\psi(U^p)_\km$ and the natural map $ R_{F,p}^\ps \to \T_{\psi}(U^p)_\km$.
     	\item \textbf{the automorphic side} $\tau_{\mathrm{Aut}}$: which comes from the action of $R_{F,p}^\ps $ (after twisting $\varepsilon^{-1}$ ) on $P_{\kB_\km}$ via $E_{\kB_\km}$ due to the map $ R_{F,p}^\ps \to E_{\kB_\km}$ (see Theorem \ref{Paskunas theory for G}).
     \end{enumerate}

     We first recall the main local-global compatibility results proved in \cite{pan2022fontaine} and \cite{Zhang:2024aa}.
     
     \begin{thm}[Local-global compatibility]\label{lgc}
     	Keep notations as above.
     	\begin{enumerate}
     	\item We have $M_\psi(U^p)_\km \in \kC_{\GL_2(F_v),\psi}(\cO)^{\kB_v}$, and hence  $M_\psi(U^p)_\km \in \kC_{D_p^\times,\psi}(\cO)^{\kB_\km}$.
     	
     	\item  Both actions $\tau_{\Gal},\tau_{\mathrm{Aut}}$ of $R_v^\ps $ on $\mm_v$ are the same. Hence, $\tau_{\Gal} = \tau_{\mathrm{Aut}}$.
     \end{enumerate}
     \end{thm}

       \begin{proof}
       	These assertions are proved in \cite[Theorem 3.5.3]{pan2022fontaine} and \cite[Theorem 3.3.3]{Zhang:2024aa}. 
       \end{proof}

       The following results are direct consequences from the previous theorem.
       
       \begin{prop}\label{finiteness from lgc}
       	Keep notations as above.
       	\begin{enumerate}
       	\item  $\mm$ is a faithful, finitely generated $\T_\psi(U^p)_\km$-module. 
       	\item  The natural map $ R_{F,p}^\ps \to \T_{\psi}(U^p)_\km$ is finite.
       \end{enumerate}
       	
       \end{prop}
  
  \begin{proof}
  	See \cite[Corollary 3.5.8]{pan2022fontaine} and \cite[Corollary 3.3.7]{Zhang:2024aa}.
  \end{proof}
  
    Now we fix a place $v|p$. Our next goal is to study the action of $ R_v^\ps$ on $\mm_v$. 
    
    \begin{defn}\label{r dR}
    	We say that a one-dimensional prime $\kp$ of $ R_v^\ps$ is \textit{regular de Rham} (resp. \textit{crystalline}) if it is the inverse image of a maximal ideal $\kp'$ of $\T_{\psi}(U^p)_\km[1/p] $ via the composite map $R_v^\ps \to \T_{\psi}(U^p)_\km \to  \T_{\psi}(U^p)_\km[1/p]$ such that the Galois representation $\rho(\kp')$ (in the sense of Definition \ref{construction}) is de Rham (resp. crystalline) with distinct Hodge-Tate weights at $v$.
    \end{defn}
    
    The remaining part is crucial in this article.
    
    \begin{lem}\label{intersection}
    	Write $I_v$ for the kernel of the map $ R_v^\ps \to \T_{\psi}(U^p)_\km$. Assume the character $\psi$ is crystalline at $v$ of Hodge-Tate weight $w_\psi$. Further assume that the tame level $U^p$ satisfies  $\psi|_{U^p\cap(\A_F^\infty)^\times}$ is trivial and that $U^pK_p$ is sufficiently small (in the sense of \cite[page 1046]{pan2022fontaine}). Then we have $$
    	I_v = \operatorname{Ann}_{R_v^\ps}(\mm_v) = \bigcap_{\kp~\textnormal{regular de Rham}} \kp = \bigcap_{\kp~\textnormal{regular crystalline}} \kp .$$
    	
    \end{lem}
    
    \begin{proof}
    	We prove the following inclusions:
    	$$
    	I_v \subseteq \bigcap_{\kp~\textnormal{regular de Rham}} \kp \subseteq \bigcap_{\kp~\textnormal{regular crystalline}} \kp \subseteq \operatorname{Ann}_{R_v^\ps}(\mm_v) = \operatorname{Ann}_{R_v^\ps} (\mm) \subseteq I_v .$$
    	Here the action of $ R_v^\ps$ on $ \mm$ is given by the action of  $R_{F,p}^\ps$ on $ \mm$ and the natural map $R_v^\ps \to R_{F,p}^\ps $.
    	
    	
    	The first inclusion is clear as $I_v \subset \kp$ for any regular de Rham prime $\kp$.
    	
    	The second one is also clear because crystalline representations must be de Rham as well.
    	
    	The proof of the third inclusion is based on the proof of Theorem \ref{lgc}, and we keep most notations and definitions as in the proof of \cite[Theorem 3.3.3]{Zhang:2024aa}. 
    	
    	Choose $U_p = \GL_2(\cO_{F_v})U^v$ with $U^v$ an open subgroup of  $\prod_{w\neq v, w|p}\GL_2(\cO_{F_w})$. Then for any $(\vec{k},\vec{w})\in \Z_{>1}^{\Hom(F,\overbar{\Q_p})}\times  \Z^{\Hom(F,\overbar{\Q_p})}$ such that $k_\sigma+2w_\sigma=w_\psi+2$ is independent of the embedding $\sigma: F \hookrightarrow \overbar{\Q_p}$, we get a natural surjection $\T_\psi(U^p)[\frac{1}{p}]\to\T_{(\vec{k},\vec{w}),\psi}(U^pU_p,K)$ sending $T_w$ to $T_w$ for $w\notin \Sigma$, where $ \T_{(\vec{k},\vec{w}),\psi}(U^pU_p,K)$ is Hecke algebra for classical automorphic forms (see \cite[3.3.2]{pan2022fontaine} for the precise definition). Let $\kp$ be a prime ideal of $\T_{(\vec{k},\vec{w}),\psi}(U^pU_p,K)\otimes_K\overbar{\Q_p}$, corresponding to an automorphic representation $\pi_\kp=\pi_\kp^\infty\otimes(\pi_\kp)_\infty$ on $\DAi$.
    	
    	Fix the level $U^v$ and let $(\pi_\kp)_v$ be the local representation of $\pi_\kp$ at place $v$. Write $d(\kp)>0$ for the multiplicity satisfying $(\pi_\kp^\infty)^{U^pU^v}\cong (\pi_\kp)_v^{\oplus d(\kp)} $, and choose a model $\pi^{K(\kp)}_v$ of $(\pi_\kp)_v$ over a finite extension $K(\kp) $ of $K$. Then we get a natural map
    	\[\Phi_{\kp}:W^{*}_{(\vec{k},\vec{w}),K}\otimes_K (\pi^{K(\kp)}_v)^{\oplus d(\kp)}\to (S_\psi(U^p)\otimes_{\cO} K(\kp))[\kp],\]
    	where $W^{*}_{(\vec{k},\vec{w}),K}$ is the dual of $W_{(\vec{k},\vec{w}),K}=\bigotimes_{\sigma:F\to K}(\Sym^{k_\sigma-2}(K^2)\otimes \det{}^{w_\sigma})$. We denote the closure of the image of $\Phi_{\kp} $ by $\Pi(\kp)$. 
    	
    	Write $\Pi_{\kB_{\km,v}}:=\Hom^{\cont}_\cO(P_{\kB_{\km,v}},K)$. This is a Banach space with unit ball $\Hom^{\cont}_\cO(P_{\kB_{\km,v}},\cO)$. Then by \cite[Lemma 3.3.4]{Zhang:2024aa}, the natural inclusions $\Pi(\kp) \hookrightarrow S_{\psi}(U^p)\otimes_{\cO} K(\kp)$ for all such $\kp$ induce an injective map:
    	$$t_v: \mm_v \otimes K \hookrightarrow\prod_{U^v}\prod_{(\vec{k},\vec{w})}\prod_\kp \Hom^{\cont}_{K(\kp)[\GL_2(F_v)]}(\Pi(\kp),\Pi_{\kB_{\km,v}}\otimes K(\kp)),$$
    	where $U^v$ runs over all open subgroups of $\prod_{w\neq v,w|p} \GL_2(\cO_{F_w})$, the pair $(\vec{k},\vec{w})$ runs over all elements  in $\Z_{>1}^{\Hom(F,\overbar{\Q_p})}\times  \Z^{\Hom(F,\overbar{\Q_p})}$ satisfying $k_\sigma+2w_\sigma=w_\psi+2$  for any $\sigma$, and $\kp$ runs over all the maximal ideals of $\T_{(\vec{k},\vec{w}),\psi}(U^p\GL_2(\cO_{F_v})U^v)_\km\otimes \overbar{\Q_p}$. 
    	
    	We claim that the intersection of all such $\kp$ is contained in  $\operatorname{Ann}_{R_v^\ps}(\mm_v) $. To show this, we prove that the action of $ R_v^\ps$ on each factor of the right hand side of the map $t_v$ factors through $R_v^\ps/\kp_v$, where $\kp_v$ is the inverse image of $\kp$ via the map $R_v^\ps \to \T_{\psi}(U^p)_\km \to  \T_{\psi}(U^p)_\km[1/p]$. 
    	
    	If $\pi_\kp$ factors through the reduced norm $N_{D/F}$, then by Jacquet-Langlands correspondence, the Galois representation $\rho(\kp)$ associated to $\kp$ has the form $\eta \oplus \eta\varepsilon^{-1}$ for some character $\eta$. Then by \cite[Proposition 4.37]{Pa_k_nas_2013} and \cite[Corollary 6.10]{Pa_k_nas_2021}, the action of $R_v^\ps$ on $\Hom^{\cont}_{K(\kp)[\GL_2(F_v)]}(\Pi(\kp),\Pi_{\kB_{\km,v}}\otimes K(\kp)) $ factors through $R_v^\ps/\kp_v$. Note that $\psi$ is crystalline at $v$. We know that $\kp_v$ is a regular crystalline prime.
    	
    	Now we assume $\pi_\kp$ does not factor through the reduced norm $N_{D/F}$, and by Jacquet-Langlands correspondence, $\pi_\kp$ corresponds to a regular algebraic cuspidal automorphic representation of $\GL_2(\A_F)$. Let $\Pi_v$ be the universal unitary completion of $(\Sym^{k_{\sigma_v-2}}(K^2)\otimes \det{}^{w_{\sigma_v}})^*\otimes \pi^{K(\kp)}_v$ as a $K$-representation of $\GL_2(F_v)$. Note that from our definition, the level of $\pi_\kp$ at $v$ is $\GL_2(\cO_{F_v}) $ and hence unramified. In this case, by \cite[Fact 4.17]{Gee_2022}, $\pi^{K(\kp)}_v$ is an irreducible principal series. By \cite[Lemma 3.5.7]{pan2022fontaine}, $\Pi(\kp)$ is a quotient of $\Pi_v^{\oplus d(\kp)'}$, where $ d(\kp)'$ is some multiple of $d(\kp) $.
    	
    	Let $\Pi_v^0$ be a $\GL_2(F_v)$-invariant bounded open ball of $\Pi_v$ and let $(\Pi_v^0)^d=\Hom_\cO(\Pi_v^0,\cO)$ be its Schikhof dual. Then by the previous arguments, we have
    	\[\Hom^{\cont}_{K[\GL_2(F_v)]}(\Pi(\kp),\Pi_{\kB_{\km,v}})\hookrightarrow  K^{\oplus d(\kp)'}\otimes \Hom_{\kC_{\GL_2(F_v),\psi|_{F_v^\times}}(\cO)}(P_{\kB_{\km,v}},(\Pi_v^0)^d).\]
    	Therefore, to prove our claim, we just need to prove that the action of $R_v^\ps$ (as the centre of $E_{\kB_v} $ by \cite[Theorem A.2.2]{Zhang:2024aa}) on $\Hom(P_{\kB_{\km,v}},(\Pi_v^0)^d)\otimes K$ factors through $R_v^\ps/\kp_v$.
    	
    	If $\rho(\kp)|_{G_{F_v}}$ is absolutely irreducible, then our claim follows from \cite[Proposition 4.37]{Pa_k_nas_2013} and \cite[Proposition 6.11]{Pa_k_nas_2021} and the classical local-global compatibility at $p$. If $\rho(\kp)|_{G_{F_v}}$ is reducible, then our claim follows from \cite[Proposition 4.37]{Pa_k_nas_2013} and \cite[Corollary 6.10]{Pa_k_nas_2021}. See \cite[page 1061]{pan2022fontaine} for more details.
    	
    	 By the classical local Langlands correspondence, the Galois representation $\rho(\kp)$ associated to $\kp$ (as a maximal prime of $\T_\psi(U^p)_\km[\frac{1}{p}]$) is crystalline of distinct Hodge-Tate weights at $v$. Note that $S_{\psi}(U^p)_\km$ is $p$-torsion free (by Proposition \ref{3.1.2}). We know that the intersection of all such $\kp_v$ is contained in $\operatorname{Ann}_{R_v^\ps}(\mm_v)$. Thus, the third inclusion follows directly.
    	
    	For the fourth result, note that we have $\End_{E_{\kB_v}}(\mm_v) \cong \End_{\kC_{\GL_2(F_v),\psi}(\cO)^{\kB_v}}(M_\psi(U^p)_\km)$ and $\End_{E_{\kB_\km}}(\mm) \cong \End_{\kC_{D_p^\times,\psi}(\cO)^{\kB_\km}}(M_\psi(U^p)_\km)$ by Theorem \ref{lgc}. The actions of $R_v^\ps$  on $\mm$ and $\mm_v$ factor through $\T_\psi(U^p)_\km$ and commute with the action of $D_p^\times$. Then the equality holds.
    	
    	The last one follows from (1) of Proposition \ref{finiteness from lgc}.
    \end{proof}
    
    \begin{rem}\label{de Rham}
    	1) If $\psi$ is de Rham at $v$, then after a finite twist, by the previous lemma and \cite[3.2.2]{pan2022fontaine}, we have 
    	$$
    	I_v = \operatorname{Ann}_{R_v^\ps}(\mm_v) = \bigcap_{\kp~\textnormal{regular de Rham}} \kp. $$
    	
    	2) We see that the local deformation problem defined by $I_v$ here not only encodes information about the $p$-adic Hodge type, but also comes from restrictions of global Galois representations.
    \end{rem}
    
    Now we assume $\psi$ is de Rham at $v$. We fix a subgroup $U'^p$ of (the tame level) $U^p$ such that $\psi|_{U'^p\cap(\A_F^\infty)^\times}$ is trivial and $U'^pK_p$ is sufficiently small. For each $v|p$, we write $R_v=R_v^\ps/I_v$, where $I_v$ is defined as in Lemma \ref{intersection}.
    
    \begin{cor}\label{finiteness version 1}
    	There exists a finite map $\widehat{\otimes}_{v|p} R_v \to \T_{\psi}(U^p)_\km$.
    \end{cor}
    
    \begin{proof}
    	By Lemma \ref{intersection} and 1) of Remark \ref{de Rham}, we get natural maps $R_v^\ps \twoheadrightarrow R_v \hookrightarrow \T_{\psi}(U'^p)_\km $. Thus, the map $R_{F,p}^\ps \to \T_{\psi}(U'^p)_\km$ (by the universal property) induces maps $\widehat{\otimes}_{v|p} R_v \to \T_{\psi}(U'^p)_\km \twoheadrightarrow  \T_{\psi}(U^p)_\km$. Then our result follows from (2) of Proposition \ref{finiteness from lgc}.
    \end{proof}
    
    The following result states some properties of $R_v$.
    
    \begin{prop}\label{R_v}  Assume $\psi$ is de Rham (resp. crystalline) at $v$.
    	\begin{enumerate}
    		\item The ring $R_v$ is reduced and $\cO$-flat.
    		\item  For each minimal prime $\kq$ of $R_v$, $\kq$ is the intersection of all regular de Rham (resp. crystalline) primes containing $\kq$.
    		\item  Assume that $\tr \bar{\rho}_{\km, v} $ is not of the form $\eta + \eta\omega$ for any character $\eta$ if $p=3$. Then we have $\dim R_v \le 3$.
    	\end{enumerate}
    \end{prop}
    
    \begin{proof}
    	By Lemma \ref{intersection} and 1) of Remark \ref{de Rham}, we know that $I_v$ is the intersection of all regular de Rham (resp. crystalline) primes. Then we get an injection $R_v \hookrightarrow \prod_\kp R_v/\kp$, where $\kp$ ranges over all regular de Rham (resp. crystalline) primes. As $ \prod_\kp R_v/\kp$ is reduced and $\cO$-flat, we get the first claim.
    	
    	Note that $I_v$ is radical. Write $I_v= \kq_1 \cap \cdots \cap \kq_t$, where $\kq_i$ are distinct  minimal primes of $R_v$. Write $ \mathfrak{r}_{\kq_i}$ for the intersection of all regular de Rham (resp. crystalline) primes containing $\kq_i $, and we have $I_v= \cap \mathfrak{r}_{\kq_i}$. We show that $ \mathfrak{r}_{\kq_i} = \kq_i$ for each $i$. As $ \cap \mathfrak{r}_{\kq_i}  =I_v \subset \kq_i$, by \cite[Lemma 10.15.1]{stacks-project}, we have $\kq_j \subset \mathfrak{r}_{\kq_j} \subset \kq_i$ for some minimal prime $\kq_j$. Hence, we have $\kq_i= \kq_j $ and $ \kq_i =  \mathfrak{r}_{\kq_i}$. This shows the second claim.
    	
    	 %If $R_v$ is not $\cO$-flat, then $\varpi$ is a zero divisor, and hence is contained in some minimal prime $\kq_i$. Combining our first claim, we know that $\varpi$ is contained in some regular de Rham prime, which is contrary to our construction. This proves the second claim.
    	
    	Note that $\mm_v$ is an $E_{\kB_v}$-module. We have isomorphisms
    	$$ \End_{E_{\kB_v}}(\mm_v) \cong \End_{\kC_{\GL_2(F_v),\psi}(\cO)^{\kB_v}}(M_\psi(U^p)_\km) \cong \End_{\GL_2(F_v), \psi}(S_\psi(U^p, K/\cO)_\km).$$
    	Here the first isomorphism follows from \cite[Proposition 4.19]{Pa_k_nas_2013}, and the second one follows from (1) of Theorem \ref{lgc} and the equivalence between $\mathrm{Mod}^{\mathrm{l\, adm}}_{\GL_2(F_v),\psi}(\cO)^{\kB_v} $ and $\kC_{\GL_2(F_v),\psi}(\cO)^{\kB_v} $ (see the discussions in Subsection \ref{pas sec}). By Lemma \ref{intersection} and 1) of Remark \ref{de Rham}, we have natural maps $R_v^\ps \twoheadrightarrow R_v \hookrightarrow \End_{E_{\kB_v}}(\mm_v)$. Here we use the fact that $R_v^\ps $ is the centre of $ E_{\kB_v}$ (see \cite[Theorem A.2.2]{Zhang:2024aa}). Then the last claim follows from \cite[Corollary 5.13]{Paskunas_2021} (taking $\tau = S_\psi(U^p, K/\cO)_\km$ there).
    \end{proof}
    
    \begin{rem}\label{rem}
    	1) The assumption $\tr \bar{\rho}_{\km, v} \ne \eta + \eta\omega$ when $p=3$ is necessary for \cite[Corollary 5.13]{Paskunas_2021}. The reason is that in the proof of Pašk\=unas' result, he needs the exact calculation of the ring $E_{\kB_v}$, which has been computed in \cite[Section 6-10]{Pa_k_nas_2013} for the other cases. As the behavior of the block in the excluded case is more complicated, the calculation is much harder. See \cite[Subsection 1.2]{Pa_k_nas_2021}. If one can compute the ring $E_{\kB_v}$ in that case and further prove an analogous result as \cite[Corollary 5.13]{Paskunas_2021}, then we can remove the assumption in (3) of Proposition \ref{R_v}.
    	
    	2) Combined with Proposition \ref{local pseudo}, the third part of Proposition \ref{R_v} is equivalent to the non-triviality of the ideal $I_v$. From the Zariski density of points with prescribed $p$-adic Hodge theoretic properties in the local deformation space (see \cite[Section 6]{B_ckle_2023} for example), we can find that not every local regular de Rham representation arises from the restriction of a global Galois representation. Hence, the ring $R_v$ is different from Kisin's potential semi-stable (or crystalline) deformation ring.
    	
    	3) It is a natural and rather intriguing question to ask which local regular de Rham primes actually occur as global ones, i.e. arise as the restriction to $G_{F_v}$ of some global (pseudo-)representation. The results above only provide a very coarse answer to this question, and it would be interesting to know whether there is a more concrete, purely local characterization of those global regular de Rham primes.
    \end{rem}
    
    From now on, we call the case that $\tr \bar{\rho}_{\km, v} $ is of the form $\eta + \eta\omega$ for some character $\eta$ if $p=3$ and $v|p$ \textit{exceptional}.
    
    \begin{thm}\label{dim of T}
    	Assume that we are not in the exceptional case. Then each irreducible component of $ \T_{\psi}(U^p)_\km$ is of characteristic zero and of dimension $1+2[F:\Q]$.
    \end{thm}
    
    \begin{proof}
    	By Proposition \ref{com alg} and Proposition \ref{R_v}, we know that the completed tensor product $\widehat{\otimes}_{v|p} R_v$ is of Krull dimension at most $1+2[F: \Q]$ (as $p$ splits completely in $F$). By Corollary \ref{finiteness version 1}, we know that the dimension of  $ \T_{\psi}(U^p)_\km$ is also at most $1+2[F: \Q]$. By \cite[Theorem 3.6.1]{pan2022fontaine} (or \cite[Corollary 3.3.10]{Zhang:2024aa}), each irreducible component is of characteristic zero and of dimension at least  $1+2[F: \Q]$. Thus, our result follows directly.
    \end{proof}
    
    \section{Pro-modularity}
    In this section, we recall some modularity or automorphy lifting theorems.
    
    Throughout this section, we fix an odd prime $p$, and $F$ an abelian totally real field in which $p$ is unramified. Let $\Sigma$ be a finite set of finite places containing $\Sigma_p$, the set of finite places $v$ of $F$ lying above $p$. 
    
    We denote by $K$  a $p$-adic local field with its ring of integers $\cO$, a uniformizer $\varpi$ and its residue field $\F$. Let $\bar{\chi}: G_{F, \Sigma} \to \F^\times$ be a continuous, totally odd character, and suppose that $\bar{\chi} $ can be extended to a character of $G_{\Q}$.
    
    \subsection{Fontaine-Mazur conjecture in the residually reducible case}\label{sec deformation ring}
    In this subsection, we recall some progress on the Fontaine-Mazur conjecture in the residually reducible case.
    
    Let $\chi: G_{F, \Sigma} \to \cO^\times$ be a continuous, totally odd character lifting $ \bar{\chi}$ such that $\chi$ is de Rham for any $v|p$. Let $\Lambda_F:=\widehat{\bigotimes}_{v|p}\cO[[\cO_{F_v}^\times(p)]] $ be the Iwasawa algebra for $F$, where $ \cO_{F_v}^\times(p)$ is the $p$-adic completion of $ \cO_{F_v}^\times$.
    
    \textbf{Case 1}: ordinary Fontaine-Mazur conjecture when $ \bar{\chi}|_{G_{F_v}} \ne \mathbf{1}$ for any $v|p$. 
    
    Let $R^{\ps,\ord}$ be the universal ordinary pseudo-deformation ring of the pseudo-representation $1+\bar{\chi}$ for $G_{F, \Sigma} $) with determinant $\chi$. Write $T^{\operatorname{univ}}: G_{F, \Sigma} \to R^{\ps,\ord}$ for the universal pseudo-deformation, and for each $v|p$, we assume $T^{\operatorname{univ}}|_{G_{F_v}}=\psi^{\operatorname{univ}}_{v,1}+\psi^{\operatorname{univ}}_{v,2}$, where $\psi^{\operatorname{univ}}_{v,1},\psi^{\operatorname{univ}}_{v,2}:G_{F_v}\to (R^{\ps,\ord})^\times$ are liftings of $\mathbf{1},\bar{\chi}|_{G_{F_v}}$, respectively.
    For each $v|p$, the character $ \psi^{\operatorname{univ}}_{v,1}$ induces a homomorphism $\cO[[\cO_{F_v}^\times(p)]]\to R^{\ps,\ord} $ via the local class field theory, and hence we get a natural map $\Lambda_F \to R^{\ps,\ord} $ by taking the completed tensor product.
    
    \begin{thm}[ordinary case 1] \label{ord1} Keep notations and assumptions as above.
    		\begin{enumerate}
    			\item $R^{\ps,\ord}$ is a finite $\Lambda_F$-algebra.
    			\item Let $\kp$ be any maximal ideal of $R^{\ps,\ord}[\frac{1}{p}]$. Let $\rho(\kp)$ be the Galois representation $G_{F,\Sigma}\to\GL_2(k(\kp))$ corresponding to $\kp$ (in the sense of Definition \ref{construction}). Assume 
    			\begin{itemize}
    				\item $\rho(\kp)$ is absolutely irreducible.
    				\item For any $v|p$, $\rho(\kp)|_{G_{F_v}}\cong\begin{pmatrix}\psi_{v,1} & *\\ 0 & \psi_{v,2}\end{pmatrix}$ such that $\psi_{v,1}$ is de Rham and has strictly less Hodge-Tate number than $\psi_{v,2}$ for any embedding $F_v\hookrightarrow \overbar{\Q_p}$.
    			\end{itemize}
    			Then $\rho(\kp)$ comes from a twist of a Hilbert modular form.
    		\end{enumerate}
    \end{thm}
    
    \begin{proof}
    	This is \cite[Theorem 5.1.1]{pan2022fontaine}.
    \end{proof}
    
    \textbf{Case 2}: ordinary Fontaine-Mazur conjecture when $ \bar{\chi}|_{G_{F_v}} = \mathbf{1}$ for any $v|p$.
    
    Let $R^{\ps,\ord}_1$ be the universal ordinary pseudo-deformation ring which pro-represents the functor from $\cOf$ to the category of sets sending $R$ to the set of tuples $(T;\{\psi_{v,1}\}_{v|p})$, where
    \begin{itemize}
    	\item $T:G_{F,\Sigma}\to R$ is a pseudo-representation lifting $1+\bar\chi$ with determinant $\chi$.
    	\item For any $v|p$, $\psi_{v,1}:G_{F_v}\to R^\times$ is a lifting of the trivial character, and we have
    	$T|_{G_{F_v}}=\psi_{v,1}+\psi_{v,2}$ (hence $\psi_{v,2}=\chi|_{G_{F_v}}(\psi_{v,1})^{-1}$) such that $T$ is $\psi_{v,1}$-ordinary in the sense of \cite[Definition 5.3.1]{pan2022fontaine}.
    \end{itemize}
    
    Write $(T_1^{\operatorname{univ}}, \{\psi^{\operatorname{univ}, \mathbf{1}}_{v,1}\}_{v|p}) $ for the universal object. The character $ \psi^{\operatorname{univ}, \mathbf{1}}_{v,1}$ induces a homomorphism $\cO[[\cO_{F_v}^\times(p)]]\to R^{\ps,\ord}_1 $ via the local class field theory, and further we get a natural map $\Lambda_F \to R^{\ps,\ord}_1 $ by taking the completed tensor product.
    
     \begin{thm}[ordinary case 2] \label{ord2} Keep notations and assumptions as above.
    	\begin{enumerate}
    		\item $R^{\ps,\ord}_1$ is a finite $\Lambda_F$-algebra.
    		\item Let $\kp$ be any maximal ideal of $R^{\ps,\ord}_1[\frac{1}{p}]$. Let $\rho(\kp)$ be the Galois representation $G_{F,\Sigma}\to\GL_2(k(\kp))$ corresponding to $\kp$ (in the sense of Definition \ref{construction}). Assume 
    		\begin{itemize}
    			\item $\rho(\kp)$ is absolutely irreducible.
    			\item For any $v|p$, $\rho(\kp)|_{G_{F_v}}\cong\begin{pmatrix}\psi_{v,1} & *\\ 0 & \psi_{v,2}\end{pmatrix}$ such that $\psi_{v,1}$ is de Rham and has strictly less Hodge-Tate number than $\psi_{v,2}$ for any embedding $F_v\hookrightarrow \overbar{\Q_p}$.
    		\end{itemize}
    		Then $\rho(\kp)$ comes from a twist of a Hilbert modular form.
    	\end{enumerate}
    \end{thm}
    
    \begin{proof}
    	This is \cite[Theorem 6.1.1]{pan2022fontaine}.
    \end{proof}
    
    \textbf{Case 3}: non-ordinary Fontaine-Mazur conjecture.
    
    In this case, we suppose that $p$ splits completely in $F$. 
    
    \begin{thm} [non-ordinary case] \label{non-ord} Keep notations and assumptions as above.
     Let
    	$\rho:G_{F, \Sigma} \to \GL_2(\cO)$
    	be a continuous irreducible representation with the following properties
    	\begin{itemize}
    		\item Let $\bar{\rho}$ be the reduction of $\rho$ modulo $\varpi$. We assume its semi-simplification has the form $\bar{\chi}_1\oplus\bar{\chi}_2$ with $\bar{\chi}_1\bar{\chi}_2^{-1}=\bar{\chi}$.
    		\item $\rho|_{G_{F_v}}$ is irreducible and de Rham of distinct Hodge-Tate weights for any $v|p$. 
    	\end{itemize}
    	Then $\rho$ comes from a twist of a Hilbert modular form.
    \end{thm}
    
    \begin{proof}
    	See \cite[Theorem 7.1.1]{pan2022fontaine} and \cite[Theorem 5.1.1]{Zhang:2024aa}.
    \end{proof}
    
    \subsection{Potential pro-modularity}\label{potential pro-modular}
    In this subsection, we recall the potential pro-modularity argument in \cite{Zhang:2024aa} and \cite{Zhang:2024ab} and prove a new case.
    
    Assume that $F$ is an abelian totally real field in which $p$ splits completely. Recall that $\bar{\chi}: G_{F, \Sigma} \to \F^\times$ is a continuous, totally odd character, which can be extended to a character of $G_\Q$. Let $\chi : G_{F, \Sigma} \to \cO^\times$ be a de Rham character as a lifting of $\bar{\chi}$.
    
    By Grunwald-Wang's theorem (see \cite[Theorem 5, page 80]{artin}), we can find a finite extension $F^1/F$ satisfying the following assumptions:
    	\begin{itemize}
    	\item $F^1$ is a totally real abelian extension of $\mathbb{Q}$ of even degree.
    	\item $p$ splits completely in $F^1$.
    	\item  For any $v \in \Sigma^1 \setminus \Sigma^1_p$, we have $p| \operatorname{Nm}(v)-1$.
    	\item $[F^1:\mathbb{Q}]>7|\Sigma^1 \setminus \Sigma^1_p|+3$.
    	\item  $\chi$ is unramified outside $ \Sigma^1_p$, and $\chi(\operatorname{Frob}_v) \equiv 1~\operatorname{mod}~\varpi$ for any $v \in \Sigma^1 \setminus \Sigma^1_p$,
    \end{itemize}
    where $ \Sigma^1$ (resp. $\Sigma^1_p $) is the finite set of finite places of $F^1$ lying above places in $\Sigma$ (resp. $\Sigma_p $).
    
    Let $R_{F^1}^{\ps, \chi}$ be the universal pseudo-deformation ring of the pseudo-representation $1+\bar{\chi}|_{G_{F^1, \Sigma^1}}$ with determinant $\chi|_{G_{F^1, \Sigma^1}}$. We say that a prime of $R_{F^1}^{\ps, \chi}$ is \textit{pro-modular}, if it comes from some big Hecke algebra (see \cite[Definition 4.1.3]{Zhang:2024aa}). Note that by Corollary \ref{lower}, we have $\dim R_{F^1}^{\ps, \chi} \ge 2[F^1:\Q]+1 $. Recall that $R_{F^1,0}^{\ps, \chi}= R_{F^1}^{\ps, \chi}/I_{F^1}$, where $ I_{F^1}$ is the intersection of all minimal primes of $R_{F^1}^{\ps, \chi} $ of dimension at least $2[F^1:\Q]+1 $. We say that a prime of $R_{F^1,0}^{\ps, \chi}$ is \textit{pro-modular}, if it comes from some pro-modular prime of $R_{F^1}^{\ps, \chi}$.
    
    \begin{prop} \label{potential R=T}
    	For each irreducible component $C$ of $R_{F^1}^{\ps, \chi} $ of dimension at least $1+2[F^1:\Q] $, its generic point is pro-modular. Equivalently, every prime of $R_{F^1,0}^{\ps, \chi}$ is pro-modular.
    \end{prop}
    
    \begin{proof}
    	For the case $\bar{\chi}|_{G_{F^1_v}} \ne \mathbf{1}$ for any $v|p$,  it is proved in \cite[Theorem 5.5.1]{Zhang:2024aa}.
    	
    	Now we assume $\bar{\chi}|_{G_{F^1_v}} = \mathbf{1}$ for any $v|p$. We sketch the main steps here.
    	
    	We first claim that if $\mathfrak{r}$ is an irreducible one-dimensional pro-modular prime  of $R_{F^1}^{\ps, \chi} $, then every irreducible component of $R_{F^1}^{\ps, \chi} $ containing $\mathfrak{r}$ is pro-modular.
    	
    	The proof is similar to the proof of \cite[Proposition 4.1.5]{Zhang:2024ab}. Let $Z$ be the set of all irreducible components of $\Spec R_{F^1}^{\ps, \chi}$. Write $Z= Z_1 \amalg Z_2$, where $Z_1$ is the subset of $Z$ consisting of all pro-modular irreducible components. We show that if $Z_1$ is not empty, then $Z_2$ is empty.
    	
    	If not, by 1) of Proposition \ref{irreducible point}, there exist an irreducible component $Y \in Z_2$ and a pro-modular prime $\mathfrak{r}' \in Y$ of dimension at least $2[F^1: \mathbb{Q}]-1$. Then arguing as in the proof of \cite[Proposition 4.1.5]{Zhang:2024ab}, we may find a prime $\mathfrak{r}'' \subset \mathfrak{r}'$ satisfying the following conditions:
    	\begin{itemize}
    		\item $\varpi \in \mathfrak{r}''$.
    		\item  The dimension of $\mathfrak{r}'' $ is at least $ 2[F^1: \mathbb{Q}]-7|\Sigma^1\setminus\Sigma^1_p|-2 \ge [F^1: \mathbb{Q}]+2$.
    		\item  For any irreducible one-dimensional prime $\mathfrak{t}$ containing $\mathfrak{r}''$, there exists a two-dimensional Galois representation $\rho(\mathfrak{t}) : G_{F^1, \Sigma^1} \to \operatorname{GL}_2(k(\mathfrak{t}))$ satisfying the conditions in (2) of \cite[Definition 4.1.1]{pan2022fontaine}. (As $\bar{\chi}|_{G_{F^1_v}} = \mathbf{1}$ for any $v|p$, the third one follows from \cite[Lemma 4.1.6 (2)]{pan2022fontaine}).
    	\end{itemize}
    	
    	Now choose an irreducible one-dimensional prime containing $\mathfrak{r}''$, and we know that it is actually a nice prime in the sense of \cite[Definition 4.1.3]{Zhang:2024aa}. Therefore, from our assumption on $F^1$, \cite[Proposition 3.1.5]{Zhang:2024ab} and \cite[Theorem 4.1.1]{Zhang:2024aa} imply that $Y$ is pro-modular, which is contrary to our assumption. Hence, we only need to show that $Z_1$ is not empty.
    	
    	Note that $\mathfrak{r}$ is pro-modular. Using \cite[Corollary 3.3.10]{Zhang:2024aa}, we know that there exists a pro-modular prime $\mathfrak{r}_1$ contained in $\mathfrak{r}$ of dimension at least $1+2[F^1: \mathbb{Q}]$. Arguing as above, we can find a nice prime containing $\mathfrak{r}_1$. Then \cite[Proposition 3.1.5]{Zhang:2024ab} and \cite[Theorem 4.1.1]{Zhang:2024aa} imply that  at least one of the irreducible component containing $\mathfrak{r}$ is pro-modular. Hence, $Z_1$ is not empty.
    	
    	Next we show that the irreducible component $C$ contains an irreducible one-dimensional pro-modular prime. 
    	
    	We use some arguments in \cite[Subsection 7.3]{pan2022fontaine}. Let $R^{\ps}_v $ be the universal pseudo-deformation ring of the pseudo-representation $\mathbf{1}+\mathbf{1}$ for $G_{F^1_v}$ with determinant $\chi|_{G_{F^1_v}}$, and by \cite[Corollary 9.1.3]{Pa_k_nas_2013}, $R^{\ps}_v $ is a power series ring of relative dimension $3$ over $\cO$. Let $R_{v, \mathbf{1}}^{\ps,\ord}$ be the universal deformation ring of the one-dimensional representation $\mathbf{1}$ (trivial character) of $G_{F^1_v}$, and it is a power series ring of relative dimension $2$ over $\cO$. For each character $\psi_{v,1}$ lifting $\mathbf{1}$, $\psi_{v,1}+\psi_{v,1}^{-1}\chi$ defines a pseudo-deformation of $\mathbf{1}+\mathbf{1}$, which induces a finite map $\phi: R_v^\ps \to R_{v, \mathbf{1}}^{\ps,\ord}$. Consider the surjective map:
    	\[\varphi: R^{\ps}_v\hat{\otimes}_{\cO} R^{\ps,\ord}_{v,\mathbf{1}}\to R^{\ps,\ord}_{v,\mathbf{1}}:~a\otimes b\mapsto \phi(a)b.\]
    	  By \cite[Lemma 10.106.4]{stacks-project}, the kernel of $\varphi$ can be generated by three elements.
    	
    	Let $R^{\ps,\ord}_{p,\mathbf{1}}$ (resp. $R_{p, F^1}^\ps$) be the completed tensor products of all $R^{\ps,\ord}_{v,\mathbf{1}} ,v|p$ (resp. $R_v^\ps, v|p$) over $\cO$, and let $R^{\ps,\ord}_1$ be the universal ordinary pseudo-deformation ring defined in the previous subsection. Then we have $R^{\ps,\ord}_1=R_{F^1}^{\ps, \chi}\otimes_{R^{\ps}_{p, F^1}}R^{\ps,\ord}_{p,\mathbf{1}}$. Write $\kP$ for the generic point of $C$, and let $C^{\ord, \mathbf{1}} = C \cap \Spec R^{\ps,\ord}_1$.  Then $C^{\ord,\mathbf{1}}$ is the underlying space of the spectrum of $(R_{F^1}^{\ps, \chi}/\kP)\otimes_{R^{\ps}_{p, F^1}}R^{\ps,\ord}_{p,\mathbf{1}}$, which we may rewrite as 
    	\[((R_{F^1}^{\ps, \chi}/\kP)\hat{\otimes}_{\cO} R^{\ps,\ord}_{p,\mathbf{1}})\otimes_{(R^{\ps}_{p, F^1} \hat{\otimes}_{\cO} R^{\ps,\ord}_{p,\mathbf{1}})}R^{\ps,\ord}_{p,\mathbf{1}}.\]
    	Note that $\dim C \ge 1+2[F^1:\Q]$. Then by the previous discussions and Krull's principal ideal theorem, we have $\dim C^{\ord,\mathbf{1}} \ge 1+[F^1:\Q]$. We choose an irreducible component $ C_1^{\ord,\mathbf{1}}$ of $C^{\ord,\mathbf{1}} $ with the largest dimension.
    	
    	Note that in our case, the Iwasawa algebra $\Lambda_{F^1}$ is regular of dimension $1+[F^1:\Q] $. Then by Theorem \ref{ord2} and the inequality $\dim C_1^{\ord,\mathbf{1}} \ge 1+[F^1:\Q]$, we get a finite surjective morphism $C_1^{\ord,\mathbf{1}} \to \Spec \Lambda_{F^1} $. As Leopoldt's conjecture holds for $F^1$, we know that the dimension of the reducible locus of $ R_{F^1}^{\ps, \chi}$ is at most $2 ~( < 1+[F^1: \Q])$. Hence, modular points are Zariski dense in $C_1^{\ord,\mathbf{1}} $. In particular, combining \cite[Proposition 3.1.5]{Zhang:2024ab} and \cite[Lemma 4.8.5]{pan2022fontaine}, we can construct a suitable big Hecke algebra with a maximal ideal defined by the pseudo-representation $1+\bar{\chi}|_{G_{F^1, \Sigma^1}}$. Furthermore, we can find an irreducible one-dimensional pro-modular prime in $C_1^{\ord,\mathbf{1}} $.
    	
    	Consequently, we obtain our result.
    	
    \end{proof}
    
    \section{Zariski density of modular points}
    
    In this section, we prove the Zariski density of modular points in some deformation space.
    
    For notations, we fix an odd prime $p$, and $F$ an abelian totally real field in which $p$ is unramified. Let $\Sigma$ be a finite set of finite places containing $\Sigma_p$, the set of finite places $v$ of $F$ lying above $p$. 
    
    We denote by $K$  a $p$-adic local field with its ring of integers $\cO$, a uniformizer $\varpi$ and its residue field $\F$. Let $\bar{\chi}: G_{F, \Sigma} \to \F^\times$ be a continuous, totally odd character, and suppose that $\bar{\chi} $ can be extended to a character of $G_{\Q}$. We say that the case $ \bar{\chi}|_{G_{F_v}}=\omega=\omega^{-1}$ for any $v \in \Sigma_p$ and $p=3$ is \textit{exceptional}.
    
    \subsection{A finiteness result (I)}
    
    In this subsection, we prove a finiteness result as an analogue to \cite[Theorem 1]{Allen_2014} in the residually reducible case.
    
    Let $\chi: G_{F, \Sigma} \to \cO^\times$ be a continuous, totally odd character lifting $\bar{\chi}$ such that $\chi$ is de Rham for any $v \in \Sigma_p$. 
    
    Let $R_{F}^{\ps, \chi}$ be the universal pseudo-deformation ring of the pseudo-representation $1+\bar{\chi}$ for $ G_{F, \Sigma}$ with determinant $\chi$. Let $R_v^\ps$ be the local universal pseudo-deformation ring of $1+\bar{\chi}|_{G_{F_v}}$  for $ {G_{F_v}}$ with determinant $\chi|_{G_{F_v}}$. Let $R_{p, F}^\ps$ be the completed tensor product (over $\cO$) of all $R^{\ps}_v, v|p$ . Then by the universal property, we get a natural map $R_{p, F}^\ps \to R_{F}^{\ps, \chi}$.
    
    \begin{prop}\label{finiteness1}
    	The map $R_{p, F}^\ps \to R_{F}^{\ps, \chi}$ is finite.
    \end{prop}
    
    \begin{proof}
    	Let $\mathfrak{m}_p$ be the maximal ideal of $R_{p, F}^\ps$. We show that $ R_{F}^{\ps, \chi}/(\mathfrak{m}_p)$ is an Artinian local ring. Assume it is not true. Then we can find a prime $\kp$ of $ R_{F}^{\ps, \chi}/(\mathfrak{m}_p)$ of dimension one and we deduce a contradiction.
    	
    	 Note that $\kp$ is of characteristic $p$. By the universal property, we know that the pseudo-representation $T(\kp)$ defined by $\kp$ satisfies that $T(\kp)|_{G_{F_v}} = 1+\bar{\chi}|_{G_{F_v}}$ for each $v|p$. In particular, $T(\kp)$ defines an ordinary pseudo-representation. As $\bar{\chi} $ can be extended to a character of $G_{\Q}$, we can divide the proof into the following two cases.
    	
    	If $\bar{\chi}|_{G_{F_v}} \ne \mathbf{1}$ for any $v|p$, let $R^{\ps, \ord}$ be the universal ordinary pseudo-deformation ring as defined in Subsection \ref{sec deformation ring}. Then we get a surjection $R^{\ps, \ord} \twoheadrightarrow R_{F}^{\ps, \chi}/\kp$ as both are quotients of $ R_{F}^{\ps, \chi}$. Thus, by Theorem \ref{ord1}, the composite map $\Lambda_F \to R^{\ps, \ord} \to R_{F}^{\ps, \chi}/\kp$ is finite. From our construction, the kernel of the map $ \Lambda_F \to R_{F}^{\ps, \chi}/\kp$ is the maximal ideal of the Iwasawa algebra $\Lambda_F $, and the map $\Lambda_F/\mathfrak{m}_{\Lambda_F} \to R_{F}^{\ps, \chi}/\kp$ is finite. This is impossible as $1=\dim R_{F}^{\ps, \chi}/\kp \le \dim \Lambda_F/\mathfrak{m}_{\Lambda_F} =0$.
    	
    	If $\bar{\chi}|_{G_{F_v}} = \mathbf{1}$ for any $v|p$, let $R_1^{\ps, \ord}$ be the universal ordinary pseudo-deformation ring as defined in Subsection \ref{sec deformation ring}. By Definition \ref{construction}, we can construct a Galois representation $\rho(\kp)$, and further by \cite[Lemma 5.3.2]{pan2022fontaine}, we may assume that $T(\kp)$ is $\psi_{v,1}$-ordinary (in the sense of \cite[Definition 5.3.1]{pan2022fontaine}) for each $v|p$ and some $\psi_{v,1}$ lifting the trivial character $\mathbf{1}$ of $G_{F_v}$. Thus, we get a surjection $R_1^{\ps, \ord} \to R_{F}^{\ps, \chi}/\kp $. By Theorem \ref{ord2}, a similar argument as in the previous paragraph leads to a contradiction as well. 
    	
    	Now we can conclude the proof by applying Nakayama's lemma.
    \end{proof}
    
    \begin{rem}
    	1) This result is not used in our main proof, but we may compare it with the finiteness result proved in the next subsection. See Remark \ref{rem on fin}.
    	
    	2) In the proof of Proposition \ref{finiteness1}, we use a similar strategy as in \cite{Allen_2014}. Also, we can prove a finiteness result about the unramified pseudo-deformation ring as in \cite[Theorem 1]{Allen_2014} using the following arguments.
    	
    	Keep all the assumptions on $F$ and $\bar{\chi}$, and further suppose that $\bar{\chi}$ is unramified at $v|p$. Let $\chi$ be a finite order (hence de Rham) character lifting $\bar{\chi}$, unramified at each $v|p$. Then we can define a universal unramified pseudo-deformation ring $R^{\ps, \operatorname{unr}}$ of the pseudo-representation $1+\bar{\chi}$ for $G_{F, \Sigma\setminus\Sigma_p}$ with determinant $\chi$. Then by Theorem \ref{ord1} (or Theorem \ref{ord2}), the finite map $\Lambda_F \to R^{\ps, \ord}$ (or $\Lambda_F \to R_1^{\ps, \ord}$) induces a finite map $\cO \to R^{\ps, \operatorname{unr}}$. Moreover, we may use such finiteness result to study the unramified Fontaine-Mazur conjecture in the residually reducible case.
    \end{rem}
    
    
    \subsection{A finiteness result (II)}\label{sec finite 2}
    In this subsection, we prove another finiteness result based on our local-global compatibility arguments.
    
    Throughout this subsection, we assume that $p$ splits completely in $F$, and let  $\chi: G_{F, \Sigma} \to \cO^\times$ be a continuous, totally odd character such that $\chi|_{G_{F_v}}$ is de Rham for any $v \in \Sigma_p$. 
    
    As in Subsection \ref{potential pro-modular}, by Grunwald-Wang's theorem \cite[Theorem 5, page 80]{artin}, we can find a finite extension $F^1/F$ satisfying the following assumptions:
    \begin{itemize}
    	\item $F^1$ is a totally real abelian extension of $\mathbb{Q}$ of even degree.
    	\item $p$ splits completely in $F^1$.
    	\item  For any $v \in \Sigma^1 \setminus \Sigma^1_p$, we have $p| \operatorname{Nm}(v)-1$.
    	\item $[F^1:\mathbb{Q}]>7|\Sigma^1 \setminus \Sigma^1_p|+3$.
    	\item  $\chi|_{G_{F^1}}$ is unramified outside $ \Sigma^1_p$, and $\chi(\operatorname{Frob}_v) \equiv 1~\operatorname{mod}~\varpi$ for any $v \in \Sigma^1 \setminus \Sigma^1_p$.
    \end{itemize}
    Here $ \Sigma^1$ (resp. $\Sigma^1_p $) is the finite set of finite places of $F^1$ lying above places in $\Sigma$ (resp. $\Sigma_p $). For simplicity, we sometimes write $\chi|_{G_{F^1}}$ for $\chi$ as well.
    
    Let $D$ be a quaternion algebra over $F^1$ ramified exactly at all infinite places. We fix an isomorphism between $(D \otimes_{F^1}\A_{F^1}^\infty )^\times$ and $\GL_2(\A_{F^1}^\infty)$. We write $\psi:=\chi\varepsilon$ and view it as a character of $(\A_{F^1}^\infty )^\times/(F^1)^\times_{+}$ via global class field theory. Define a tame level $U^p=\prod_{v\nmid p}U_v$ as follows: $U_v=\GL_2(\cO_{F^1_v})$ if $v\notin \Sigma^1$ and 
    \[U_v=\mathrm{Iw}_v:=\{g\in\GL_2(\cO_{F^1_v}),g\equiv \begin{pmatrix}*&*\\0&*\end{pmatrix}\mod \varpi_v\}\]
    otherwise. Then we can define a big Hecke algebra $\T:=\T_\psi(U^p)$. By the proof of Proposition \ref{potential R=T}, we know that $1+\bar{\chi}|_{G_{F^1, \Sigma^1}}$ defines a maximal ideal $\km$ of $\T $.
    
    Let $R_{F^1}^{\ps, \chi}$ be the universal pseudo-deformation ring of the pseudo-representation $1+\bar{\chi}|_{G_{F^1, \Sigma^1}}$ with determinant $\chi|_{G_{F^1, \Sigma^1}}$. Then by the universal property, we have a natural surjection $R_{F^1}^{\ps, \chi} \twoheadrightarrow \T_\km$. We say that a prime of $R_{F^1}^{\ps, \chi}$ is \textit{pro-modular}, if it comes from the big Hecke algebra $\T_\km$.
    
     \begin{rem}
    	In \cite{pan2022fontaine} and \cite{Zhang:2024aa}, a pro-modular prime is defined to be the one coming from some big Hecke algebra $\T_{\psi, \xi}(U^p)_\km$, where $\xi_v: k(v)^\times \to \cO^\times, v\in \Sigma^1\setminus\Sigma_p^1$ are some characters of $p$-power order and we view $\xi:={\xi_v}$ as some character of $U^p$ via class field theory. Under our setting, combining \cite[Proposition 3.1.5]{Zhang:2024ab} and classical local-global compatibility at $v$, we know that for any minimal prime $\kp$ of $\T_\km$, the natural surjection $\T_\km \twoheadrightarrow \T_\km/\kp  $ factors through $\T_{\psi, \xi}(U^p)_\km$ for some $\xi$. Hence, in this subsection, a pro-modular prime means that it comes from a prime of $\T_\km$.
    \end{rem}
    
    Note that by Corollary \ref{lower}, we have $\dim R_{F^1}^{\ps, \chi} \ge 2[F^1:\Q]+1 $. Recall that $R_{F^1,0}^{\ps, \chi}= R_{F^1}^{\ps, \chi}/I_{F^1}$, where $ I_{F^1}$ is the intersection of all minimal primes of $R_{F^1}^{\ps, \chi} $ of dimension at least $2[F^1:\Q]+1 $. Consider the natural composite map $f: R_{F^1}^{\ps, \chi} \twoheadrightarrow \T_\km \twoheadrightarrow (\T_\km)^{\operatorname{red}}$. By \cite[Corollary 3.3.10]{Zhang:2024aa}, every minimal prime of $(\T_\km)^{\operatorname{red}}$ is of dimension at least $1+2[F^1:\Q]$. Thus, the map $f$ factors through $R_{F^1,0}^{\ps, \chi}$, and we get a surjection $R_{F^1,0}^{\ps, \chi} \twoheadrightarrow (\T_\km)^{\operatorname{red}}$.
    
    \begin{prop}\label{R=T}
    	The surjection $R_{F^1,0}^{\ps, \chi} \twoheadrightarrow (\T_\km)^{\operatorname{red}}$ is an isomorphism.
    \end{prop}
    
   \begin{proof}
   	Note that by Proposition \ref{def of 0}, both $R_{F^1,0}^{\ps, \chi} $ and $(\T_\km)^{\operatorname{red}}$ are reduced. Thus, we only need to show that each prime ideal of $R_{F^1,0}^{\ps, \chi}$ is pro-modular, and it follows directly from Proposition \ref{potential R=T}.
   \end{proof}
    
    For each $v \in \Sigma_p^1$ (resp. $v \in  \Sigma_p$), let $R_v^\ps$ be the local universal pseudo-deformation ring of $1+\bar{\chi}|_{G_{F^1_v}}$ (resp. $1+\bar{\chi}|_{G_{F_v}}$)  for $ {G_{F^1_v}}$ (resp. $ {G_{F_v}}$) with determinant $\chi|_{G_{F^1_v}}$ (resp. $\chi|_{G_{F_v}}$). Choose a subgroup $U^p_1 \subseteq U^p$ such that  $U^p_1$ satisfies the assumptions in Lemma \ref{intersection}. For each $v \in \Sigma_p^1$, we get an ideal $I_v$ of $R_v^\ps$. Write $R_v = R_v^\ps/I_v $.  
    
    Consider a place $v \in \Sigma_p$, and choose a place $v_1 \in \Sigma_p^1$ lying above $v$. Then by the universal property, we have a natural map $f_1: R_{v_1}^\ps \to R_v^\ps$. This is actually an isomorphism via $\sigma_1: F_{v_1}^1 \cong F_v (\cong \Q_p)$. Let $I_v$ be the image of $I_{v_1}$ and write $R_v = R_v^\ps/I_v $.  
    
    Note that for any $v', v'' \in \Sigma_p^1$ lying above $v$, there exists an element $\sigma \in \Gal(F^1/F)$ such that $\sigma(v')=v''$ and there exists an isomorphism $ R_{v'}^\ps \cong R_{v''}^\ps$ induced by $\sigma$. As $\sigma(F^1) = F^1$, we know that for any regular de Rham prime $\kp_{v'}$ of $R_{v'}^\ps$, there exists a regular de Rham prime $\kp_{v''}$ of $R_{v''}^\ps$ such that $\sigma(\kp_{v'})=\kp_{v''}$. Combining Lemma \ref{intersection}, we know that $\sigma$ also induces an isomorphism $R_{v'} \cong R_{v''}$. Furthermore, the natural surjection $\widehat{\otimes}_{v'| v} R_{v'}^\ps \twoheadrightarrow R_{v}^\ps$ induces a surjection $ \widehat{\otimes}_{v'| v} R_{v'} \twoheadrightarrow R_v$.
    
    For simplicity, we write $R_{F^1,p}^\ps$ (resp. $R_{F,p}^\ps$) for $ \widehat{\otimes}_{v \in \Sigma^1_p} R_{v}^\ps$ (resp. $\widehat{\otimes}_{v \in \Sigma_p} R_{v}^\ps$).
    \begin{lem}\label{com dia1} 
    	1) We have the following commutative diagram:
    	\[\begin{tikzcd}
    		R_{F^1,p}^\ps  \arrow[r, two heads] \arrow[d, two heads]& \widehat{\otimes}_{v \in \Sigma^1_p} R_{v} \arrow[d, two heads]\\
    		R_{F,p}^\ps \arrow[r, two heads] & \widehat{\otimes}_{v \in \Sigma_p} R_{v}.
    	\end{tikzcd}\]
    	
    	2) We have $(\widehat{\otimes}_{v \in \Sigma^1_p} R_{v})\widehat{\otimes}_{R_{F^1,p}^\ps} R_{F,p}^\ps = \widehat{\otimes}_{v \in \Sigma_p} R_{v}$.
    \end{lem}
    
    \begin{proof}
     	For the first part, it is direct from our discussions above.
     	
     	For the second part, consider a place $v \in \Sigma_p$ and let $\Sigma^1_v$ be the subset of $\Sigma_p^1$ consisting of places lying above $v$. We get a natural map $f_v: \widehat{\otimes}_{v' \in \Sigma^1_v} R_{v'}^\ps \twoheadrightarrow R_v^\ps$. By 3) of Proposition \ref{com alg} and Proposition \ref{R_v}, the kernel of the surjection $ \widehat{\otimes}_{v' \in \Sigma^1_v} R_{v'}^\ps \twoheadrightarrow \widehat{\otimes}_{v' \in \Sigma^1_v} R_{v'}$ is the ideal $(I_{v'}, \dots)$, $v' \in \Sigma_v^1$. We claim that the image of this ideal via the map $f_v$ is $I_v$. Note that $f_v$ is given by the norm map. In other words, for a non-zero element $x_{v'} \otimes x_{v''} \otimes \cdots \in \widehat{\otimes}_{v' \in \Sigma^1_v} R_{v'}^\ps$, the image of it via $f_v$ is non-zero if and only if for any $v', v'' \in \Sigma_v^1 $, $x_{v'}=x_{v''}^\sigma$, where $\sigma$ is an isomorphism between $F_{v'}^1$ and $F_{v''}^1$. As $I_{v'}\cong I_{v''}$ (induced by $\sigma$), we prove the claim. Now taking the completed tensor product (for all $v \in \Sigma_p$) and combining 3) of Proposition \ref{com alg}, we get the desired result.
    \end{proof}
    
    Similarly to $R_{F^1}^{\ps, \chi}$ (resp. $R_{F^1,0}^{\ps, \chi} $), we can also define the universal pseudo-deformation ring $R_{F}^{\ps, \chi} $ (resp. $R_{F, 0}^{\ps, \chi} $ using Corollary \ref{lower}). Then we have natural maps $R_{F,p}^\ps \to R_{F}^{\ps, \chi} \twoheadrightarrow R_{F, 0}^{\ps, \chi}$.
    
    \begin{lem}\label{com dia2}
    	We have the following diagram:
    	\[\begin{tikzcd}
    		R_{F^1,p}^\ps  \arrow[r] \arrow[d, two heads]& R_{F^1,0}^{\ps, \chi} \arrow[d]\\
    		R_{F,p}^\ps \arrow[r] & R_{F, 0}^{\ps, \chi}.
    	\end{tikzcd}\]
    \end{lem}
    
    \begin{proof}
    	By Corollary \ref{finite for 0}, the map $R_{F^1,0}^{\ps, \chi} \to R_{F, 0}^{\ps, \chi} $ is well-defined and induced by the map $ R_{F^1}^{\ps, \chi} \to R_{F}^{\ps, \chi}$. Our result follows from the universal property.
    \end{proof}
    
    Now we can conclude our finiteness result in this subsection.
    
     \begin{prop}\label{finiteness 2}
    	There exists a finite map $\widehat{\otimes}_{v \in \Sigma_p} R_{v} \to R_{F, 0}^{\ps, \chi} $.
    \end{prop}
    
    \begin{proof}
    	We consider the following diagram:
    		\[\begin{tikzcd}
    		R_{F^1,p}^\ps  \arrow[r, two heads] \arrow[d, two heads]& \widehat{\otimes}_{v \in \Sigma^1_p} R_{v} \arrow[d, two heads]\arrow[r, dashed] & R_{F^1,0}^{\ps, \chi} \arrow[d] \\
    		R_{F,p}^\ps \arrow[r, two heads] & \widehat{\otimes}_{v \in \Sigma_p} R_{v} \arrow[r, dashed] & R_{F, 0}^{\ps, \chi}.
    	\end{tikzcd}\]
    	
    	We first show that this diagram is commutative. The commutativity of the left square is from 1) of Lemma \ref{com dia1}. Using Proposition \ref{R=T}, we have an isomorphism $R_{F^1,0}^{\ps, \chi} \cong (\T_\km)^{\operatorname{red}}$. By the proof of Corollary \ref{finiteness version 1}, we know that the map $R_{F,p}^\ps \to \T_\km$ factors through $\widehat{\otimes}_{v \in \Sigma_p^1} R_v$. Therefore, the dashed map in the first horizontal line exists. The existence of the other dashed one and the commutativity follow from 2) of Lemma \ref{com dia1} and Lemma \ref{com dia2}.
    	
    	Note that by Corollary \ref{finiteness version 1} and Corollary \ref{finite for 0}, both maps $\widehat{\otimes}_{v \in \Sigma^1_p} R_{v} \to R_{F^1, 0}^{\ps, \chi} $ and $R_{F^1,0}^{\ps, \chi} \to R_{F, 0}^{\ps, \chi} $ are finite. Our result holds by using the commutative diagram and \cite[Lemma 10.36.15]{stacks-project}.
    \end{proof}
    
    Recall that we say the case $ \bar{\chi}|_{G_{F_v}}=\omega=\omega^{-1}$ for any $v \in \Sigma_p$ and $p=3$ is \textit{exceptional}.
    
    \begin{prop}\label{finite same dim}
    	Assume that we are not in the exceptional case.
    	
    	1) The ring $R_{F, 0}^{\ps, \chi}$ is equidimensional of dimension $1+2[F:\Q]$. In particular, $\dim R_{F}^{\ps, \chi}=1+2[F:\Q]$.
    	
    	2) Let $R_{F}^{\ps}$ be the universal pseudo-deformation ring of the pseudo-representation $1+\bar{\chi}$ for $G_{F, \Sigma}$ (without fixed determinant). Then we have $\dim R_{F}^{\ps}=2+2[F:\Q]$.
    	
    	3) For each $v|p$, we have $\dim R_v=3$.
    \end{prop}
    
    \begin{proof}
    	As we are not in the exceptional case, by Proposition \ref{R_v}, we know that for each $v|p$, $R_v$ is $\cO$-flat of Krull dimension at most $3$. Further, by Proposition \ref{com alg}, the ring $\widehat{\otimes}_{v \in \Sigma_p} R_{v}$ is of Krull dimension at most $1+2[F: \Q]$. On the other hand, by Proposition \ref{def of 0} and Corollary \ref{lower}, each irreducible component of $ R_{F, 0}^{\ps, \chi}$ has Krull dimension at least $1+2[F:\Q]$. Combining Proposition \ref{finiteness 2}, we have $$
    	1+2[F:\Q] \le \dim R_{F, 0}^{\ps, \chi} \le \dim \widehat{\otimes}_{v \in \Sigma_p} R_{v} \le 1+2[F:\Q].$$
    	This shows both the first and the last assertions.
    	
    	For the second one, it is direct from the first one and Proposition \ref{det} as Leopoldt's conjecture holds in our case.
    	
    \end{proof}
    
    \begin{rem}\label{rem on fin}
    	So far, we have proved two independent but similar finiteness results between the completed tensor product of all local pseudo-deformation rings at $v|p$ and the global deformation ring, using different methods.
    	Both of them are analogous to Pan's ordinary finiteness results (Theorem \ref{ord1} and Theorem \ref{ord2}) in the non-ordinary case, substituting the Iwasawa algebra $\Lambda_F$ with some suitable local pseudo-deformation ring. Although Proposition \ref{finiteness1} applies for more totally real fields, Proposition \ref{finiteness 2} is much stronger so that we can control the Krull dimension of the global pseudo-deformation ring (as in Proposition \ref{finite same dim}) and further use the going-up property to find modular points in the global deformation space (see the following subsection). Note that by Proposition \ref{local pseudo}, $\widehat{\otimes}_{v \in \Sigma_p} R_{v}^\ps$ has Krull dimension $1+3[F:\Q]$, which is much larger than the Krull dimension of $R_{F}^{\ps, \chi}$ (expected to be $1+2[F:\Q]$). Thus, Proposition \ref{finiteness1} is insufficient for our purposes.
    	
    	 More precisely, we actually find a non-trivial local deformation problem defined by $I_v$ such that the natural map $R_v^\ps \to R_F^{\ps, \chi} $ factors through $R_v^\ps/I_v$ by using local-global compatibility results at $p$ and potential pro-modularity. In general, it is unclear whether the map $R_v \to R_F^{\ps, \chi} $ is well-defined. The ring $R_v$ seems to be quite mysterious, and we do not know whether it is equidimensional. Although the definition of $R_v$ comes from the big Hecke algebra for the totally real field $F^1$ in our context, we do not know whether it is actually independent of the choice of such abelian base change.
    	 
    	 In addition, if Pašk\=unas' result (\cite[Corollary 5.13]{Paskunas_2021}) holds in the exceptional case, then by our proof, Proposition \ref{finite same dim} also holds in that case without more arguments.
    	 \end{rem}
    
        \subsection{Zariski closure of modular points in the global deformation space}\label{5.3} In this subsection, we study the Zariski closure of modular points in the Galois deformation space.
        
        We keep all the assumptions on the number field $F$ and the character $\chi$ as in the previous subsection. Recall that $R_{F}^{\ps, \chi}$ is the universal pseudo-deformation ring of the pseudo-representation $1+\bar{\chi}$ for $G_{F, \Sigma}$ with determinant $\chi$, and let $I_{F}$ be the intersection of all minimal primes of dimension at least $1+2[F:\Q]$. Write $ R_{F, 0}^{\ps, \chi}= R_{F}^{\ps, \chi}/I_{F}$.
        
        Assume that we are not in the exceptional case. Recall that by Proposition \ref{finiteness 2}, there exist local deformation rings $R_v, v \in \Sigma_p$ (as quotients of the local pseudo-deformation rings $R_v^\ps, v \in \Sigma_p$) such that the map $\widehat{\otimes}_{v \in \Sigma_p} R_{v} \to R_{F, 0}^{\ps, \chi} $ is well-defined and finite.
        
        For a place $v \in \Sigma_p$ and a prime $\kp$ of $R_v^\ps$, we say that $\kp$ is \textit{ordinary} if the pseudo-representation defined by $\kp $ is ordinary, i.e., a sum of two characters. Otherwise, we say $\kp$ is \textit{non-ordinary}. Let $R_v^{\ps, \ord}$ be the local pseudo-deformation ring parametrizing all ordinary liftings of $1+\bar{\chi}|_{G_{F_v}}$ with determinant $\chi$.
        
        \begin{lem}\label{dim of v ord}
        	The ring $R_v^{\ps, \ord}$ is a power series ring over $\cO$ of relative dimension $2$. Moreover, there exists at most one minimal prime of $R_v$ which is ordinary of dimension $3$.
        \end{lem}
    
    \begin{proof}
    	We show the first assertion. For the case $\bar{\chi}|_{G_{F_v}} \ne \mathbf{1}, \omega^{\pm 1}$, one can see \cite[Remark B.28]{Pa_k_nas_2013}. For the case $\bar{\chi}|_{G_{F_v}} = \mathbf{1}$, we can prove it by \cite[Proposition 9.12 \& Corollary 9.13]{Pa_k_nas_2013}. For the case $\bar{\chi}|_{G_{F_v}}= \omega $ and $p \ge 5$, this is \cite[Corollary B.6]{Pa_k_nas_2013}.
    	
    	For the second sentence, it is clear from the first one and 3) of Proposition \ref{finite same dim}.
    \end{proof}
    
    \begin{rem}
    	The proof of this lemma relies on the explicit calculation of the local pseudo-deformation ring $R_v^{\ps}$, but such a result is not known in the exceptional case. However, we can also prove this result in that case in the following way.
    	
    	Let $(R_1, \psi_1^{\operatorname{univ}})$ and $(R_{\bar{\chi}}, \psi_{\bar{\chi}}^{\operatorname{univ}})$ be the universal deformation of the $1$-dimensional character $\mathbf{1}$ and $\bar{\chi}|_{G_{F_v}}$ respectively. As $F_v \cong \Q_p$, it is well known that both $R_1$ and $ R_{\bar{\chi}}$ are power series rings over $\cO$ of relative dimension $2$. On the one hand, by the universal property, we have natural maps $ R_v^{\ps, \ord} \to R_1 \to R_v^{\ps, \ord}$. The first one is defined by the pseudo-representation $ \psi_1^{\operatorname{univ}}+\chi|_{G_{F_v}}( \psi_1^{\operatorname{univ}})^{-1}$. If $T=\psi_1+\psi_{\bar{\chi}}$ is the universal ordinary pseudo-deformation with determinant $\chi|_{G_{F_v}}$, then the second one is defined by $\psi_1$ (lifting the trivial character). Hence, the composite map is an identity, which implies that the map $R_1 \to R_v^{\ps, \ord} $ is a surjection. Using similar discussions, we can also show that the map $R_v^{\ps, \ord} \to R_1 $ is a surjection as well. This shows that $R_v^{\ps, \ord} $ is isomorphic to $ R_1$, hence a power series ring over $\cO$ of relative dimension $2$.
    	
    	%On the other hand, let $R_v^{\ps, \operatorname{red}}$ be the universal ordinary pseudo-deformation ring (without fixed determinant). By \cite[Proposition 2.3.2]{Zhang:2024aa}, this is a power series ring over $\cO$ of relative dimension $4$. By definition, we have $R_v^{\ps, \ord}= R_v^{\ps, \operatorname{red}} \otimes_{R_{\bar{\chi}}} \cO$, where the map $R_{\bar{\chi}} \to \cO $ is defined by $\chi|_{G_{F_v}}$. Then by Krull's principal ideal theorem, we know that $R_v^{\ps, \ord}$ is of dimension at least $3$. Thus, $R_v^{\ps, \ord}$ is a power series ring over $\cO$ of relative dimension $2$.
        \end{rem}
    
        \begin{lem}\label{lem of Rv}
        	For a minimal prime $\kq$ of $R_v$ of dimension $3$, if $\kq$ is non-ordinary, then $\kq$ is the intersection of all non-ordinary regular de Rham primes containing $\kq$.
        \end{lem}
    
       \begin{proof}
       	By Proposition \ref{R_v}, we know that $\kq$ is the intersection of all regular de Rham primes containing $\kq$. Then we can write $\kq = \mathfrak{r}_1 \cap \mathfrak{r}_2$, where $\mathfrak{r}_1$ is the intersection of all non-ordinary regular de Rham primes and $\mathfrak{r}_2$ is the intersection of all ordinary ones. Then by \cite[Lemma 10.15.1]{stacks-project}, we have $\kq = \mathfrak{r}_1 $ or $\kq=\mathfrak{r}_2$.
       	
       	Let $\kp^{\ord}$ be the kernel of the surjection $R_v^\ps \twoheadrightarrow R_v^{\ps, \ord}$. By Proposition \ref{local pseudo} and Lemma \ref{dim of v ord}, it is a prime ideal of height $1$. If $\kq=\mathfrak{r}_2 $, then we have $\kq = \mathfrak{r}_2 \supset \kp^{\ord}$. This implies $ \kq = \kp^{\ord} $, and hence $\kq$ is ordinary, which is contrary to our assumption.
         \end{proof}
    
    Let $\kp$ be a minimal prime of $R_{F, 0}^{\ps, \chi} $, and let $T(\kp) $ be the pseudo-representation of $G_{F, \Sigma}$ defined by $\kp$. Recall that we have well-defined maps $R_v \to \widehat{\otimes}_{v |p} R_{v} \to R_{F, 0}^{\ps, \chi}$. Write $\kp_v$ for the inverse image of $\kp$ via the composite map.
    
    \begin{lem}\label{char 0}
    	1) The ring $R_{F, 0}^{\ps, \chi}/\kp $ is of characteristic $0$.
    	
    	2) For each $v|p$, $\kp_v$ is a minimal prime of dimension $3$.
    \end{lem}
    
    \begin{proof}
    	We prove the first assertion. If not, then for each $v|p$,  we have $\varpi \in \kp_v$. By (1) of Proposition \ref{R_v}, we know that $\kp_v$ is not a minimal prime of $R_v$, and hence of dimension at most $2$. Consider the composite map $\widehat{\otimes}_{v |p} R_{v} \to R_{F, 0}^{\ps, \chi} \twoheadrightarrow R_{F, 0}^{\ps, \chi}/\kp$, which is finite by Proposition \ref{finiteness 2}. From our definition of $\kp_v$, we know that this finite map factors through $\widehat{\otimes}_{v |p} R_{v}/\kp_v$, which is of Krull dimension at most $2[F:\Q]$. By Proposition \ref{finite same dim}, $\kp$ is of dimension $1+2[F: \Q]$. We get a contradiction.
    	
    	From the first assertion, we know that $\kp_v$ is of characteristic $0$. Then by Proposition \ref{com alg}, the ring $\widehat{\otimes}_{v |p} R_{v}/\kp_v$ is of Krull dimension $1+ 2[F: \Q] - \sum_{v|p} \textnormal{ht}(\kp_v)$. This shows the second assertion.
    \end{proof}
    
    \begin{lem}\label{min is non-ord}
    	The pseudo-representation $T(\kp) $ is non-ordinary at each $v|p$.
    \end{lem}
    
    \begin{proof}
    	The method of this lemma is similar to the one of Proposition \ref{finiteness1}. Assume that there exists a place $v\in \Sigma_p$ such that $T(\kp)|_{G_{F_v}}$ is ordinary, and we derive a contradiction.
    	
    	 By Lemma \ref{dim of v ord} and Lemma \ref{char 0}, we know that $R_v/\kp_v\cong R_v^{\ps, \ord}$. Write $\kp'$ for the kernel of the map $\widehat{\otimes}_{v |p} R_{v}/\kp_v \to R_{F, 0}^{\ps, \chi}/\kp$. By our previous discussions, $\kp'$ is also a minimal prime. Write $\km_v$ for the maximal ideal of $R_v$.
    	
    	 Recall that we have a finite map $$\widehat{\otimes}_{v |p} R_{v}/\kp_v= \widehat{\otimes}_{v' |p, v' \ne v} (R_{v'}/\kp_{v'}) \widehat{\otimes}_{\cO} R_v^{\ps, \ord} \to R_{F, 0}^{\ps, \chi}/\kp.$$
    	Consider the natural map $ \widehat{\otimes}_{v' |p, v' \ne v} (R_{v'}/\kp_{v'}) \to \widehat{\otimes}_{v' |p, v' \ne v} (R_{v'}/\kp_{v'}) \widehat{\otimes}_{\cO} R_v^{\ps, \ord}$, and let $\kp''$ be the inverse image of $\kp'$ via this map. Then we get a natural surjection $$(\widehat{\otimes}_{v' |p, v' \ne v} (R_{v'}/\kp_{v'}))/\kp'' \widehat{\otimes}_{\cO} R_v^{\ps, \ord} \twoheadrightarrow (\widehat{\otimes}_{v |p} R_{v}/\kp_v)/\kp'. $$
    	Combining Proposition \ref{com alg} and Lemma \ref{dim of v ord}, we know that both rings are integral domains (as $ R_v^{\ps, \ord}$ is a power series ring) with Krull dimension $1+2[F:\Q]$. Hence, it is an isomorphism. This implies that the ring $\widehat{\otimes}_{v' |p, v' \ne v} (R_{v'}/\km_{v'}) \widehat{\otimes}_{\cO} R_v^{\ps, \ord}/(\varpi)$, which is actually a power series ring over $\F$ of relative dimension two by Lemma \ref{dim of v ord}, is a quotient of $(\widehat{\otimes}_{v |p} R_{v}/\kp_v)/\kp' $, and we write $\kq'$ for its kernel. Consider the chain $\kp' \subset \kq'$. As $(\widehat{\otimes}_{v |p} R_{v}/\kp_v)/\kp' \to R_{F, 0}^{\ps, \chi}/\kp$ is a finite map between integral domains, by the going-up property (see \cite[Theorem 5.E]{matsumura80}), there exists a prime $\kq$ of $R_{F, 0}^{\ps, \chi}$ containing $ \kp$ lying above $\kq'$.
    	
    	We first consider the case $\bar{\chi}|_{G_{F_v}} \ne \mathbf{1}$ for any $v|p$. Let $T(\kq')$ be the pseudo-deformation corresponding to the prime $\kq'$. It has the property that for any $v' \in \Sigma_p, v' \ne v$, $T(\kq')|_{G_{F_{v'}}} $ is just the trivial lifting, and for the place $v$, $T(\kq')|_{G_{F_{v}}} $ corresponds to the local pseudo-representation defined by $ R_v^{\ps, \ord}/(\varpi)$. In particular, $T(\kq')$ is an ordinary pseudo-deformation lifting $1+\bar{\chi}$, and hence we get a surjection $R^{\ps, \ord} \twoheadrightarrow R_{F, 0}^{\ps, \chi}/\kq'$, where $R^{\ps, \ord}$ is the universal ordinary pseudo-deformation ring defined in Subsection \ref{sec deformation ring}. Then by Theorem \ref{ord1}, the composite map $\Lambda_F \to R^{\ps, \ord} \twoheadrightarrow R_{F, 0}^{\ps, \chi}/\kq'$ is finite. From our construction of the prime $\kq'$, the kernel of the composite map contains the ideal $(\varpi, X_{v'})$ for all $v' \ne v$, where $X_{v'} $ corresponds to the generator of $\cO_{F_{v'}}^\times(p)$. In other words, we get a finite map $\F[[X_v]] \to R_{F, 0}^{\ps, \chi}/\kq'$. We get a contradiction as $\dim \F[[X_v]] =1$ and $\dim R_{F, 0}^{\ps, \chi}/\kq' =\dim R_v^{\ps, \ord}/(\varpi) =2$.
    	
    	Now consider the case $\bar{\chi}|_{G_{F_v}} = \mathbf{1}$ for any $v|p$. By Definition \ref{construction}, the pseudo-deformation $T(\kq') \otimes k(\kq')$ defines a Galois representation $\rho(\kq') : G_{F, \Sigma} \to \GL_2(k(\kq'))$. Combining \cite[Lemma 5.3.2]{pan2022fontaine}, we get a surjection $R_1^{\ps, \ord} \twoheadrightarrow R_{F, 0}^{\ps, \chi}/\kq'$, where $R_1^{\ps, \ord}$ is the universal ordinary pseudo-deformation ring defined in Subsection \ref{sec deformation ring}. Then we can use the same arguments as in the previous paragraph and get a contradiction as well, replacing Theorem \ref{ord1} by Theorem \ref{ord2}.
    \end{proof}
    
    We say that a one-dimensional prime $\mathfrak{r}$ of $\widehat{\otimes}_{v |p} R_{v}/\kp_v$ is \textit{non-ordinary regular de Rham}, if for any $v|p$, the inverse image of $\mathfrak{r}$ via the composite map $R_v \to \widehat{\otimes}_{v |p} R_{v}/\kp_v \twoheadrightarrow (\widehat{\otimes}_{v |p} R_{v}/\kp_v)/\mathfrak{r}$ is non-ordinary and regular de Rham.
    
    Let $\kp'$ be the inverse image of the prime $\kp$ via the finite map $\widehat{\otimes}_{v |p} R_{v} \to R_{F, 0}^{\ps, \chi}$. Then $\kp'$ is a minimal prime of $\widehat{\otimes}_{v |p} R_{v}$, and also a minimal prime of the ideal $(\kp_{v_1}, \cdots, \kp_{v_d})$, where $\{v_i: 1 \le i \le d \} = \Sigma_p$ and $d=[F: \Q]$.
    
    \begin{lem}\label{non-ord regular de Rham}
    	The prime $\kp'$ is the intersection of all non-ordinary regular de Rham primes containing $\kp'$.
    \end{lem}
    
    \begin{proof}
    	Write $\sqrt{(\kp_{v_1}, \cdots, \kp_{v_d})} = \kp_1' \cap \cdots \cap \kp_s'$, and we prove that for each $\kp_i', 1 \le i \le s$, it is the intersection of all non-ordinary regular de Rham primes containing itself.
    	
    	By Lemma \ref{lem of Rv} and Lemma  \ref{min is non-ord}, we know that for each $ \kp_{v_i}$, it is the intersection of all non-ordinary regular de Rham primes containing $ \kp_{v_i}$. Then by 3) and 5) of Proposition \ref{com alg}, we have $$
    	(\kp_{v_1}, \cdots, \kp_{v_d})= \cap (\kq_{v_1}, \cdots, \kq_{v_d}),$$
    	where $\kq_{v_i}$ ranges over all non-ordinary regular de Rham primes containing $ \kp_{v_i}$.
    	
    	As $R_v/\kq_v$ is $\cO$-flat and one-dimensional, by 4) of Proposition \ref{com alg}, we know that the ideal $(\kq_{v_1},  \cdots, \kq_{v_d})$ is radical. If we write $(\kq_{v_1},  \cdots, \kq_{v_d}) =\sqrt{(\kq_{v_1},  \cdots, \kq_{v_d})}= \mathfrak{r}_1 \cap \cdots \cap \mathfrak{r}_t$ for the primary decomposition, then by our definition, each $\mathfrak{r}_i$ is a non-ordinary regular de Rham prime. Then we have the following relations.
    	$$\kp_1' \cap \cdots \cap \kp_s'=\sqrt{(\kp_{v_1}, \cdots, \kp_{v_d})}=\sqrt{\cap (\kq_{v_1}, \cdots, \kq_{v_d})}= \sqrt{\cap \mathfrak{r}} = \cap \mathfrak{r},$$
        where $\mathfrak{r}$ ranges over all non-ordinary regular de Rham primes containing the ideal $(\kp_{v_1}, \cdots, \kp_{v_d})$.
        
        Write $\mathfrak{a}_i, 1 \le i \le s$ for the intersection of all non-ordinary regular de Rham primes containing the minimal prime $\kp_i'$. Then by \cite[Lemma 10.15.1]{stacks-project}, we know that there exists $j, 1 \le j \le s$, such that $\kp_i' \supset \mathfrak{a}_j \supset \kp_j'$. Therefore, we know that $i=j$, and $\kp_i'=\mathfrak{a}_i$. This shows the result.
    \end{proof}
    
    By a \textit{modular point} in the deformation space $\Spec R_{F}^{\ps, \chi}$, we mean that the corresponding Galois representation arises from a twist of a Hilbert modular form, or equivalently, a regular algebraic cuspidal automorphic representation of $\GL_2(\A_F)$.
    
    \begin{thm}\label{main thm}
    	The Zariski closure of the set of all modular points in the deformation space $\Spec R_{F}^{\ps, \chi}$ is $\Spec R_{F, 0}^{\ps, \chi}$, i.e. the union of all irreducible components of dimension $1+2[F: \Q]$. If further $\chi$ is crystalline at each $v|p$, then the Zariski closure of all modular points which are crystalline at each $v|p$ in the deformation space $\Spec R_{F}^{\ps, \chi}$ is $\Spec R_{F, 0}^{\ps, \chi}$
    \end{thm}
    
    \begin{proof}
    	Let $\mathcal{X} \subset \Spec R_{F}^{\ps, \chi}$ be the set consisting of all modular points and $\overline{\mathcal{X}}$ be its Zariski closure in $\Spec R_{F}^{\ps, \chi} $.
    	
        Let $C$ be an irreducible component of $\Spec R_{F, 0}^{\ps, \chi}$ and let $\kp$ be the corresponding generic point. Let $\kp'$ be the inverse image of $\kp$ via the finite map $\widehat{\otimes}_{v |p} R_{v} \to R_{F, 0}^{\ps, \chi}$. It is a minimal prime of $\widehat{\otimes}_{v |p} R_{v}$, and we get a finite injection $(\widehat{\otimes}_{v |p} R_{v})/\kp' \hookrightarrow R_{F, 0}^{\ps, \chi}/\kp$. By Lemma \ref{non-ord regular de Rham}, $\kp'$ is the intersection of all non-ordinary regular de Rham primes containing $\kp'$. For any non-ordinary regular de Rham prime $\mathfrak{r}' \supset \kp'$, by the going-up property, we can find a prime
    	$\mathfrak{r} \supset \kp$ lying above $\mathfrak{r}'$.
    	
    	We show that the Galois representation $\rho(\mathfrak{r})$ corresponding to $\mathfrak{r}$ (in the sense of Definition \ref{construction}) arises from a twist of a Hilbert modular form. First, from our construction on $\mathfrak{r}$, we know that $\rho(\mathfrak{r})|_{G_{F_v}}$ is absolutely irreducible for each $v|p$. In particular, $\rho(\mathfrak{r})$ is absolutely irreducible. Second, for each $v \in \Sigma_p$, by our definition of regular de Rham prime (see Definition \ref{r dR}), we know that there exists a big Hecke algebra (for some abelian extension $F^1$ of $F$) $\T_{\psi}(U^p)_\km$ and a maximal prime $\mathfrak{t}$ of $\T_{\psi}(U^p)_\km[1/p]$ such that the Galois representation $\rho(\mathfrak{t})|_{G_{F^1_{v'}}}$ for some $v'$ lying above $v$ is absolutely irreducible regular de Rham and $\tr \rho(\mathfrak{t})|_{G_{F^1_{v'}}} = \tr \rho(\mathfrak{r})|_{G_{F_v}}$ via the isomorphism $F^1_{v'} \cong F_v $. By Brauer-Nesbitt theorem (see \cite[Fact 2.5]{Gee_2022}) for example), we know that for each $v \in \Sigma_p$,  $\rho(\mathfrak{r})|_{G_{F_v}}$ is absolutely irreducible and de Rham with distinct Hodge-Tate weights. Then by the non-ordinary Fontaine-Mazur conjecture (Theorem \ref{non-ord}), $\mathfrak{r}$ is a modular point.
    	
    	Write $C^{\operatorname{non-ord}} \subset C$ for the set of all non-ordinary modular points obtained in the above way contained in $C$. We claim that $C^{\operatorname{non-ord}} $ is Zariski dense in $C$. Write $\mathfrak{I}=\cap \mathfrak{r}$, where $\mathfrak{r}$ ranges over all points in $C$ that we obtain in the above way. By Lemma \ref{non-ord regular de Rham}, the inverse image of $\mathfrak{I}$ via the finite injection $f: (\widehat{\otimes}_{v |p} R_{v})/\kp' \hookrightarrow R_{F, 0}^{\ps, \chi}/\kp$ is $0$. Write  $\mathfrak{I}= \kp_1 \cap \cdots \cap \kp_c $ as its primary decomposition, and we have $0= f^{-1} (\kp_1) \cap \cdots \cap f^{-1}(\kp_c)$. It implies $ f^{-1} (\kp_i)=0$ for some $1 \le i \le c$. By the incomparability property (see \cite[Lemma 10.36.20]{stacks-project}), we have $\mathfrak{I}=\kp_i=\kp$. This proves the claim.
    	 
    	 In conclusion, we know that the Zariski closure of $\cup C^{\operatorname{non-ord}}$ contains $\cup C =\Spec R_{F, 0}^{\ps, \chi}$. In particular, we have $\overline{\mathcal{X}} \supset \Spec R_{F, 0}^{\ps, \chi}$. On the other hand, as a modular point defines an irreducible Galois representation, by Proposition \ref{def of 0}, we have $ \mathcal{X} \subset \Spec R_{F, 0}^{\ps, \chi}$. Consequently, we have $\overline{\mathcal{X}} = \Spec R_{F, 0}^{\ps, \chi}$.
    	 
    	 If further $\chi$ is crystalline at each $v|p$, then combining Lemma \ref{intersection}, we can use the same proof.
    \end{proof}
    
    Let $R_F^\ps$ be the universal pseudo-deformation ring of the pseudo-representation $1+\bar{\chi}$ for $G_{F, \Sigma}$ (without fixed determinant). We can also study the Zariski closure of modular points in the deformation space $\Spec R_F^\ps$.
    
    \begin{thm}\label{main thm without det}
    	The Zariski closure of the set of all modular points which are crystalline at each $v|p$ in the deformation space $\Spec R_F^\ps$ is the union of all irreducible components of Krull dimension $2+2[F:\Q]$.
    \end{thm}
    
    \begin{proof}
    	As $\bar{\chi}$ can be viewed as a character of $G_\Q$ and $F$ is an abelian totally real field in which $p$ splits completely, we know that the set of crystalline liftings of $\bar{\chi}$ is Zariski dense in its Galois deformation space $\Spec \cO[[G_{F, \Sigma}^{\operatorname{ab}}(p)]]$. Then our result follows from Proposition \ref{det} and Theorem \ref{main thm}.
    \end{proof}
    
    \begin{rem}
    	We observe that if Pašk\=unas' result (\cite[Corollary 5.13]{Paskunas_2021}) is also true in the exceptional case, then our theorems can be proved with the same proof.
    \end{rem}
    
    \section{Applications}
    In this section, we study some applications of our results.
    
    \subsection{Emerton's conjecture}\label{emerton's conj}
   In this subsection, we prove Emerton's conjecture about the Krull dimension of the big Hecke algebra for $p \ge 5$. Our main reference here is \cite{Emerton2011}.
   
   For notations, let $p \ge 5$ be an odd prime. We denote by $K$  a $p$-adic local field with its ring of integers $\cO$, a uniformizer $\varpi$ and its residue field $\F$. We fix a positive integer $N \ge 1$ such that $p$ does not divide $N$. Write $\Sigma$ for the set of primes of $\Q$ dividing $Np$.
   
    For an integer $k \ge 1 $, write $\mathcal{M}_k(N)$ for the space of modular forms of weight $k$ and level $\Gamma_1(N)$ (or level $N$ for simplicity). Recall that the map $\begin{pmatrix}
    	a & b\\
    	c& d
    \end{pmatrix}  \mapsto d~ \modd N$ induces an isomorphism $ \Gamma_0(N)/\Gamma_1(N) \cong (\Z/N\Z)^\times$. If $d \in (\Z/N\Z)^\times$, this map induces an automorphism of $\mathcal{M}_k(N)$ by $\langle d \rangle$, the \textit{diamond operator}.
    
    Let $l$ be a prime not dividing $N$. We define the automorphism $S_l$ of $\mathcal{M}_k(N)$ via the formula $S_l = \langle l \rangle l^{k-2}$. We also define the endomorphism $T_l$ of  $\mathcal{M}_k(N)$ via the formula 
    $$ (T_l f)(\tau)= \sum_{n=0}^{\infty} c_{nl}(f)q^n + \sum_{n=0}^{\infty} l c_n(S_lf)q^{nl},~~ \tau \in \mathcal{H},~~ q=\exp(2\pi i \tau),$$
    where $c_n(f)$ is the $n$-th Fourier coefficient of the modular form $f$.
    
    The \textit{Hecke algebra} $\T_k(N)$ (or just $\T_k$) is defined to be the $\Z$-subalgebra of $\End(\mathcal{M}_k(N))$ generated by the Hecke operators $lS_l$ and $T_l$ as $l$ ranges over primes not dividing $N$.
    
    Let $\T_{\le k}^{(p)} (N)$ (or just $\T_{\le k}^{(p)}$) be the $\Z$-algebra of endomorphisms of $\oplus_{i=1}^k \mathcal{M}_i(N)$ generated by the operators $lS_l$ and $T_l$ as $l$ ranges over primes not dividing $Np$. From the definition, we can find that if $k' > k$, then there is a surjection $\T_{\le k'}^{(p)} \twoheadrightarrow \T_{\le k}^{(p)}$. The $p$\textit{-adic Hecke algebra} $\T(N)$ (or just $\T$) is defined to be the projective limit
    $$ \T:= \varprojlim_k \Z_p \otimes_\Z \T_{\le k}^{(p)}$$. 
    
    \begin{thm}\label{decomposition}
      The ring $\T$ is $p$-adically complete, Noetherian $\Z_p$-algebra, and is the product of finitely many noetherian local $\Z_p$-algebras.
    \end{thm}
   
   \begin{proof}
   	See \cite[Theorem 2.7]{Emerton2011}.
   \end{proof}
   
   \begin{conj}\cite[Conjecture 2.9]{Emerton2011}\label{conj}
   	The ring $\T$ is equidimensional of Krull dimension $4$.
   \end{conj}
    
    Using the theory of eigencurves, we can obtain the lower bound of $\dim \T$.
    
    \begin{prop}\label{lower bound}
    	The Krull dimension of each irreducible component of $\Spec \T$ is at least $4$.
    \end{prop}
    
    \begin{proof}
    	See \cite[Corollary 2.28]{Emerton2011}.
    \end{proof}
    
     If we replace $\Z_p$ by $\cO$ in all definitions above, we can define a $p$-adic Hecke algebra $\T'$ over $\cO$. It is clear that we have $\T' = \T \otimes_{\Z_p} \cO$. From now on, we assume that our Hecke algebra $\T$ is defined over $\cO$. 
    
    By Theorem \ref{decomposition}, we know that $\T$ is semi-local. If we assume $\km_1, \cdots, \km_r$ are all maximal ideals of $\T$, combining \cite[24.C]{matsumura80}, we have
    $$ \T \cong \T_{\km_1} \times \cdots \times \T_{\km_r},$$
    where each $ \T_{\km_i}$ is $\km_i$-adically complete.
    
    Let $\km$ be a maximal ideal of $\T$. It is classical that there exists a pseudo-representation $T_\km: G_{\Q, \Sigma} \to \T_\km$ such that $T_\km(\Frob_l)=T_l$ and $\det (\Frob_l)=lS_l$ (see \cite[Lemma 4.1]{Deo2023} for example). Let $R_{\Q}^\ps$ be the universal pseudo-deformation ring of $T_\km \modd \km$ for $G_{\Q, \Sigma}$. Then by the universal property, we get a surjection $R_{\Q}^\ps \twoheadrightarrow \T_\km$.
    
    \begin{thm}\label{main thm for emerton}
    	Conjecture \ref{conj} holds for $p \ge 5$.
    \end{thm}
    
    \begin{proof}
    	Let $\km$ be a maximal ideal of $\T$. Combining Proposition \ref{lower bound}, we only need to show that $\dim \T_{\km}=4$. By Definition \ref{construction}, we can divide the proof into two cases, depending on the reducibility of $T_\km \modd \km$. 
    	
    	If $T_\km \modd \km$ is reducible ($\km$ is Eisenstein), combining Proposition \ref{finite same dim} and Proposition \ref{lower bound} (after a finite twist), we have $4 \le \dim \T_\km \le \dim R_{\Q}^\ps =4$.
    	
    	If $T_\km \modd \km$ is irreducible ($\km$ is non-Eisenstein), then the universal deformation ring and the universal pseudo-deformation ring are the same (see \cite[Theorem 2.4.1, page 44]{Berger_2013}). Combining Proposition \ref{lower bound} and Corollary \ref{dim in irr case}, we have $4 \le \dim \T_\km \le \dim R_{\Q}^\ps =4$.
    \end{proof}
    
    \subsection{Big $R=\T$ theorem (I)}
    In this subsection, we prove a big $R=\T$ theorem in the residually reducible case in the setting of completed cohomology of modular curves.
    
    For notations, let $p \ge 3$ be an odd prime. We denote by $K$  a $p$-adic local field with its ring of integers $\cO$, a uniformizer $\varpi$ and its residue field $\F$. Write $\Sigma$ for a set of primes of $\Q$ containing $p$ and write $\Sigma= \Sigma_0 \amalg \{p\}$. Let $\bar{\chi} : G_{\Q, \Sigma} \to \F^\times$ be a continuous, odd character. Assume that $\bar{\chi}|_{G_{\Q_p}} \ne \omega$ if $p=3$.
    
    We first recall the definition of completed cohomology of modular curves. Our main references are \cite[Section 5]{emerton11} and \cite[Section 2]{Nakamura2023}.
    
    For any compact open subgroup $K_f$ of $\GL_2(\A_\Q^\infty)$, we let $Y(K_f)$ denote the adelic modular curve
    $$Y(K_f) := \GL_2(\Q)\setminus ((\mathbb{C}\setminus \mathbb{R}) \times \GL_2(\A_\Q^\infty))/K_f.$$
    We may also view $Y(K_f)$ as an algebraic curve over $\Q$. We write $H^1(K_f)_A:= H_{\operatorname{et}}^1(Y(K_f)_{\bar{\Q}}, A)$, where $A$ denotes one of $K, \cO, $ or $\cO/\varpi^s\cO$ for some $s>0$.
    
    If $K^p$ is some fixed compact open subgroup of $\GL_2(\A_f^p)$, where $ \A_f^p = \Q \otimes \prod_{l \ne p} \Z_l$, we write $H^1(K^p)_A := \varinjlim_{K_p} H^1(K_pK^p)_A$. Here the direct limit is taken over all compact open subgroups $K_p$ of $\GL_2(\Q_p)$. Write $\widehat{H}^1(K^p)_\cO := \varprojlim_s H^1(K^p)_\cO/\varpi^sH^1(K^p)_\cO$ for the $\varpi$-adic completion of $H^1(K^p)_\cO $. Write $ \widehat{H}^1(K^p)_K := \widehat{H}^1(K^p)_\cO \otimes_\cO K$.
    
    Write $G_{\Sigma_0} := \prod_{l \in \Sigma_0} \GL_2(\Q_l)$ and $K_0^\Sigma := \prod_{l \notin \Sigma} \GL_2(\Z_l)$. If $K_{\Sigma_0}$ is an open subgroup of $G_{\Sigma_0}$, we write $\widehat{H}^1(K_{\Sigma_0})_A := \widehat{H}^1(K_{\Sigma_0}K_0^\Sigma)_A$ for $A=K$ or $\cO$. Also write $\widehat{H}^1_{A, \Sigma}:= \varinjlim_{K_{\Sigma_0}} \widehat{H}^1(K_{\Sigma_0})_A$. It is well-known that $\widehat{H}^1_{A, \Sigma} $ are equipped with a continuous action of $G_\Q \times \GL_2(\Q_p)\times G_{\Sigma_0}$.
    
    Let $K_\Sigma$ be a compact open subgroup of $G_\Sigma = \prod_{l \in \Sigma} \GL_2(\Q_l)$. For any prime $l \notin \Sigma$, we write $S_l$ and $T_l$ for the Hecke operators acting on $H^1(K_\Sigma K_0^\Sigma)_\cO$ corresponding to the double cosets of $K_l$ in $\GL_2(\Q_l)$ represented by $\begin{pmatrix}
    	l & 0\\
    	0& l
    \end{pmatrix} $ and $\begin{pmatrix}
    l & 0\\
    0& 1
    \end{pmatrix} $ respectively. We let $\T(K_\Sigma)_\cO$ denote the $\cO$-subalgebra of $\End_\cO(H^1(K_\Sigma K_0^\Sigma)_\cO)$ generated by $S_l$ and $T_l$ for all $\ \notin \Sigma$. If $K_{\Sigma_0}$ is an open subgroup of $G_{\Sigma_0}$, and $K_p' \subset K_p  $ are compact open subgroups of $\GL_2(\Q_p) $, then there is a natural surjection $\T(K_p'K_{\Sigma_0})_\cO \twoheadrightarrow \T(K_pK_{\Sigma_0})_\cO$. Define $\T(K_{\Sigma_0}) = \varprojlim_{K_p}\T(K_pK_{\Sigma_0})_\cO$, which is an $\cO$-algebra topologically generated by the Hecke operators $S_l$ and $T_l$ for all $l \notin \Sigma$ and acts faithfully on $ \widehat{H}^1(K_{\Sigma_0})_\cO$.
    
    We consider the pseudo-representation $\bar{T} = \bar{\chi}_1 + \bar{\chi}_2 : G_{\Q, \Sigma} \to \F$, where $\bar{\chi}_1, \bar{\chi}_2 : G_{\Q, \Sigma} \to \F^\times$ are continuous characters satisfying $\bar{\chi}_1\bar{\chi}_2^{-1}=\bar{\chi}$. It is classical that $\bar{T}$ is modular (see \cite[Lemma 2.23]{Deo2023} or \cite[Lemma 7.3]{Fakhruddin2022} for example). In other words, there exist a compact open subgroup $K_{\Sigma_0} \subset G_{\Sigma_0}$ and a maximal ideal $\km$ of $\T(K_{\Sigma_0})$ satisfying
    $$ \bar{T}(\Frob_l) \equiv T_l \modd \km,  ~~\det(\bar{T}(\Frob_l)) \equiv lS_l \modd \km$$
    for all $l \notin \Sigma$. We say such $K_{\Sigma_0}$ is an \textit{allowable level} for the pseudo-representation $\bar{T}$.
    
    For an allowable level $K_{\Sigma_0}$, there exists a pseudo-deformation $T_\km: G_{\Q, \Sigma} \to \T(K_{\Sigma_0})_\km$ satisfying 
    $$ T_\km(\Frob_l)=T_l ,  ~~\det(T_\km(\Frob_l)) = lS_l $$
    for all $l \notin \Sigma$. If $K_{\Sigma_0}' \subset K_{\Sigma_0}$ is an inclusion of allowable levels, then $\T(K_{\Sigma_0}')_\km \twoheadrightarrow \T(K_{\Sigma_0})_\km$ is a surjection. Define $\T_\km := \varprojlim_{K_{\Sigma_0}} \T(K_{\Sigma_0})_\km$, where $K_{\Sigma_0} $ ranges over all allowable levels.
    
    Let $R_\Q^\ps$ be the universal pseudo-deformation ring of $\bar{T}$.  Combining Proposition \ref{det} and Proposition \ref{finite same dim} (after a finite twist), we can define $R_{\Q, 0}^\ps := R_\Q^\ps/I_\Q$, where $I_\Q $ is the intersection of all minimal prime ideals of dimension $4$. Hence, $R_{\Q, 0}^\ps$ is equidimensional.
    
    Note that by the universal property, we get a natural surjection $R_\Q^\ps \twoheadrightarrow \T_\km $. The following result is an analogue to the big $R=\T$ theorem in the residually irreducible case proved by B\"ockle and Emerton (see \cite[Theorem 2.1]{Nakamura2023} for example).
    
    \begin{thm}\label{big R=T modular curve}
    	The kernel of the surjection $R_\Q^\ps \twoheadrightarrow \T_\km $ is contained in 
    	$I_\Q $. Moreover, every prime of $R_{\Q, 0}^\ps$ comes from some big Hecke algebra.
    \end{thm}
    
    \begin{proof}
    	For each modular point $\kp \in \Spec R_\Q^\ps$, by our definition, the natural surjection $R_\Q^\ps \twoheadrightarrow R_\Q^\ps/\kp$ factors through $\T_\km$. By Theorem \ref{main thm without det}, the kernel of $R_\Q^\ps \twoheadrightarrow \T_\km $ is contained in the intersection of all modular points, which is just $I_\Q $.
    \end{proof}
    
    \begin{rem}
    	1) If we fix a central character $\zeta : (\A_\Q^\infty)^\times/\Q_{+}^\times \to \cO^\times$ and study the subspace of  $ \widehat{H}^1(K_{\Sigma_0})_\km$ for some allowable level $K_{\Sigma_0} $ on which the centre  $(\A_\Q^\infty)^\times $ acts via $\zeta$, we can also define the universal pseudo-deformation ring $R_\Q^{\ps, \zeta}$ and the big Hecke algebra $\T_\zeta$, whose determinants are fixed corresponding to the central character $\zeta$. (See \cite[Section 6]{Pan_2022} for the details of such settings.) In this case, we can also prove a similar big $R=\T$ theorem using Theorem \ref{main thm}.
    	
    	%In that case, we may follow the proof the local-global compatibility results in the Eisenstein case (see \cite[Theorem 6.3.5]{Pan_2022}) and the method of \cite[Theorem 3.6.1]{pan2022fontaine} to obtain the lower bound of the Krull dimension of each irreducible component of $\T_\zeta$, and further show that $\Spec R_{\Q,0}^{\ps, \zeta} = \Spec \T_\zeta$ by using Theorem \ref{main thm}.
    	
    	2) We may also prove some similar results in the setting of completed cohomology of Shimura curves.
    \end{rem}
    
    \subsection{Big $R=\T$ theorem (II)}
    In this subsection, we prove a big $R=\T$ theorem in the residually reducible case in the setting of completed cohomology of quaternionic forms. 
    
    For notations, let $p \ge 3$ be an odd prime. We denote by $K$  a $p$-adic local field with its ring of integers $\cO$, a uniformizer $\varpi$ and its residue field $\F$. Let $F$ be an abelian totally real field of even degree such that $p$ splits completely in $F$. Write $\Sigma$ for a set of primes of $F$ containing $\Sigma_p$ (the set of places of $F$ lying above $p$). Let $\bar{\chi} : G_{F, \Sigma} \to \F^\times$ be a continuous, totally odd character such that it can be extended to a character of $G_\Q$. Assume that $\bar{\chi}|_{G_{\F_v}} \ne \omega$ for each $v|p$ if $p=3$. Let $\chi : G_{F, \Sigma} \to \cO^\times$ be a continuous, de Rham character lifting $\bar{\chi}$.
    
    Note that $F$ is a totally real field of even degree. There exists a quaternion algebra $D$ over $F$ ramified exactly at all infinite places. Fix an isomorphism between $(D \otimes_{F}\A_{F}^\infty )^\times$ and $\GL_2(\A_{F}^\infty)$. Write $\psi = \chi \varepsilon$ for a character of $(\A_F^\infty)^\times/F_{+}^\times $ via global class field theory.
    
    Let $U^p = \prod_{v \nmid p} U_v$ be a tame level satisfying for each $v \notin \Sigma$, $U_v = \GL_2(\cO_{F_v})$. As in Subsection \ref{sec for T}, we can define the big Hecke algebra $\T_{\psi}(U^p)$. 
    
    Consider the pseudo-representation $\bar{T} = \bar{\chi}_1 + \bar{\chi}_2 : G_{F, \Sigma} \to \F$, where $\bar{\chi}_1, \bar{\chi}_2 : G_{F, \Sigma} \to \F^\times$ are continuous characters satisfying $\bar{\chi}_1\bar{\chi}_2^{-1}=\bar{\chi}$. As $\bar{\chi}$ can be extended to a character of $G_\Q$, we know that $\bar{T}$ is modular. Hence, there exist a tame level $U^p$, a maximal ideal $\km$ of $\T_{\psi}(U^p)$ and a pseudo-representation $T_\km : G_{F, \Sigma} \to \T_{\psi}(U^p)_\km$ sending $\Frob_v$ to the Hecke operator $T_v$ for $v \notin \Sigma$ such that $ \bar{T} \equiv T_\km \modd \km$. We say that such a tame level $U^p$ is an \textit{allowable level}.
    
    Let $R_F^{\ps, \chi}$ be the universal pseudo-deformation ring of $\bar{T}$ with determinant $\chi$. For an allowable level $U^p$, by the universal property, we get a natural surjection $ R_F^{\ps, \chi} \twoheadrightarrow \T_{\psi}(U^p)_\km$. If ${U^p}' \subset U^p$ is another allowable level, then we have the following commutative diagram
    \[\begin{tikzcd}
    	R_F^{\ps, \chi} \arrow[r, two heads] \arrow[rd, two heads] & \T_{\psi}({U^p}')_\km \arrow[d, two heads]\\
      ~ &	\T_{\psi}(U^p)_\km.
    \end{tikzcd}\]
    Hence, if we write $\T_{\km} := \varprojlim_{U^p} \T_{\psi}(U^p)_\km$, where the inverse limit is taken over all allowable tame levels $U^p$, then we get a surjection $ R_F^{\ps, \chi} \twoheadrightarrow \T_{\km}$ induced by the universal property.
    
    Recall that $R_{F, 0}^{\ps, \chi} = R_F^{\ps, \chi}/I_F$, where $I_F$ is the intersection of all minimal primes of dimension $1+2[F:\Q]$. 
    
    \begin{thm}\label{big R=T quaternion}
    	The kernel of the surjection $ R_F^{\ps, \chi} \twoheadrightarrow \T_{\km}$ is contained in $I_F$. Moreover, every prime of $R_{F, 0}^{\ps, \chi}$ comes from some big Hecke algebra.
    \end{thm}
    
    \begin{proof}
    	For a modular point $\kp \in \Spec R_{F}^{\ps, \chi}$, from our construction, we know that the natural surjection $ R_{F}^{\ps, \chi} \twoheadrightarrow R_{F}^{\ps, \chi}/\kp$ factors through $ \T_{\km} $. Combining Theorem \ref{main thm}, the kernel of the surjection $ R_{F}^{\ps, \chi} \twoheadrightarrow ( \T_{\km})^{\operatorname{red}}$ is contained in $I_F$. This shows the result.
    \end{proof}
    
    \subsection{Serre's modularity conjecture in the residually reducible case}  In this subsection, we discuss the relation between our previous results and Serre's modularity conjecture in the residually reducible case.
    
    Serre's modularity conjecture states that for an odd, irreducible, continuous Galois representation $\bar{\rho}: G_{\Q} \to \GL_2(\F)$ with a finite field $\F$ of characteristic $p$ (including the case $p=2$), it arises from a newform. It has been proved in \cite{Khare_2009a} and \cite{Khare_2009}. If $\bar{\rho}$ is reducible, then we say it is \textit{modular} if there exist a Galois representation $\rho_{f}: G_{\Q} \to \GL_2(\cO_f)$ associated to some newform $f$ and a suitable choice of lattice over $\cO_f$ such that $\bar{\rho} \cong \rho_{f} \modd \mathfrak{m}_{\cO_f}$.
    
    To state it more precisely, we give some notations first. Let $p \ge 5$ be an odd prime, and let $F$ be an abelian totally real field in which $p$ splits completely. Let $\Sigma$ be a finite set of finite places containing $\Sigma_p$, the set of places $v$ of $F$ lying above $p$. We denote by $K$  a $p$-adic local field with its ring of integers $\cO$, a uniformizer $\varpi$ and its residue field $\F$. Let $\bar{\chi}: G_{F, \Sigma} \to \F^\times$ be a continuous, totally odd character, and suppose that $\bar{\chi} $ can be extended to a character of $G_{\Q}$.
    
    Let $x \in H^1(F_\Sigma/F, \F(\bar{\chi}^{-1}))$ be a non-zero cohomology class. It is easy to see that $x$ defines a two-dimensional representation (up to equivalence)
    $$ \bar{\rho}_x : G_{F, \Sigma} \to \GL_2(\F), ~~ g \mapsto \begin{pmatrix}
    	1 & x(g)\bar{\chi}(g)\\
    	0& \bar{\chi}(g)
    \end{pmatrix}. $$
    For simplicity, we choose a complex conjugation $c \in G_{F, \Sigma}$ and an element $\sigma \in G_{F, \Sigma}$, and assume $x(c)=0, x(\sigma)=1$ for the definition of $ \bar{\rho}_x$. Then it is conjectured that the Galois representation $ \bar{\rho}_x$ arises from some newform $f$. In \cite{Fakhruddin2022}, they solved most cases when $F=\Q$ (also including $x=0$) by a Galois theoretic method. 
    
     Let $\chi : G_{F, \Sigma} \to \cO^\times$ be a de Rham lifting of $\bar{\chi}$. As $x \ne 0$, we can define the universal deformation ring $R_x^{\chi}$ of the non-split Galois representation $ \bar{\rho}_x$ with determinant $\chi$. Let $ R_{F}^{\ps, \chi}$ be the universal pseudo-deformation ring parametrizing all pseudo-deformations of $1+\bar{\chi}$ with determinant $\chi$. By the universal property, we have a natural map $R_{F}^{\ps, \chi} \to R_x^{\chi} $.
     
     Let $\kp$ be a one-dimensional prime of $R_x^{\chi}$ such that the Galois representation $\rho(\kp)$ associated to $\kp$ is irreducible. Here we do not require
     that $R_x^{\chi}/\kp$ is of characteristic $0$.
     
     \begin{prop}\label{dim of x}
     	For any irreducible component of $\Spec R_x^{\chi}$ containing $\kp$, it is of Krull dimension $1+2[F:\Q]$.
     \end{prop}
    
    \begin{proof}
    	Let $C$ be an irreducible component of $\Spec R_x^{\chi}$ with generic point $\kq$. Let $\kq'$ be the kernel of the composite map $R_{F}^{\ps, \chi} \to R_x^{\chi} \twoheadrightarrow R_x^{\chi}/\kq$. We show that $C$
        is of Krull dimension	$1+2[F:\Q]$.
    	
    	By the global characteristic formula (see \cite[Proposition 3.1.3]{Zhang:2024ab}), we know that each irreducible component of $ R_x^{\chi}$ has Krull dimension at least $1+2[F:\Q]$. In particular, $\dim C \ge 1+2[F:\Q]$.
    	
    	Note that $\rho(\kp)$ is irreducible, which implies that the Galois representation $\rho(\kq)$ associated to $\kp$ is also irreducible. By \cite[Corollary 2.3.6 (2)]{pan2022fontaine}, we have $\dim R_x^{\chi}/\kq \le \dim R_{F}^{\ps, \chi}/\kq' $. Combining Proposition \ref{finite same dim}, we prove our result.
    \end{proof}
    
    Keep notations in the proof of Proposition \ref{dim of x}. Combining Proposition \ref{finite same dim}, we can find that $\kq'$ defines an irreducible component $C'$ of $\Spec R_{F}^{\ps, \chi}$. By Theorem \ref{main thm}, the set of modular points contained in $C'$ is Zariski dense. Then it is natural to pose the following conjecture.
    
    \begin{conj}\label{conj for x}
    	For an irreducible component $C$ of $\Spec R_x^{\chi}$, if it contains an irreducible one-dimensional point (an irreducible lifting of $ \bar{\rho}_x$, not necessarily of characteristic $0$), then the set of modular points contained in $C$ is Zariski dense.
    \end{conj}
    
    In general, we do not know whether there exists an irreducible lifting of $\bar{\rho}_x$ unramified outside $\Sigma$. In \cite{Fakhruddin2022}, their method first requires to enlarge the finite set $\Sigma$ appropriately, and then finds a geometric lifting of the residual Galois representation. If our conjecture is true, then it is clear that the existence of an irreducible lifting implies Serre's modularity conjecture in this case.
    
    \begin{rem}
    	We may prove Conjecture \ref{conj for x} under the assumption $H^1(F_\Sigma/F, \F(\bar{\chi}^{-1}))=1$ and $F=\Q$. Note that by Definition \ref{construction}, for every one-dimensional irreducible point of $\Spec R_{\Q}^{\ps, \chi}$, it must be a lifting of $\bar{\rho}_x$ up to a possible conjugation. As the natural map $R_{\Q}^{\ps, \chi} \to R_x^{\chi} $ is a surjection (by \cite[Corollary 1.4.4]{kisin2009fontaine}), Theorem \ref{main thm} implies Conjecture \ref{conj for x} in this case. Here, we also give another proof of \cite[Theorem A]{Deo2023}.
    	
    	Note that by Definition \ref{construction} and \cite[Lemma 2.7]{Skinner_1999}, each one-dimensional irreducible point of $\Spec R_{F}^{\ps, \chi}$ determines a unique non-zero cohomology class in $H^1(F_\Sigma/F, \F(\bar{\chi}^{-1}))$. In general, it is unclear whether the one-dimensional irreducible points lying on a given irreducible component of $\Spec R_{F}^{\ps, \chi}$ all correspond to the same cohomology class, or how to distinguish points associated to different classes in $\Spec R_{F}^{\ps, \chi}$. Also, the natural map $R_{F}^{\ps, \chi} \to R_x^{\chi} $ is not finite in general. Hence, we need some more ideas for Conjecture \ref{conj for x}.
    \end{rem}
    
    
    \appendix
    \section{A finiteness result in the residually irreducible case}\label{A}
    In this appendix, we prove a finiteness result in the residually irreducible case. It is prepared for Subsection \ref{emerton's conj}. An advantage of our result here is that we do not need the Taylor-Wiles hypothesis.
    
    Let $p \ge 5$ be an odd prime. Let $\Sigma$ be a finite set of finite places of $\Q$ containing $p$. We denote by $K$  a $p$-adic local field with its ring of integers $\cO$, a uniformizer $\varpi$ and its residue field $\F$.
    
    Let $\bar{\rho} : G_{\Q, \Sigma} \to \GL_2(\F)$ be a continuous, odd, irreducible Galois representation . By Serre's modularity conjecture (see \cite{Khare_2009a} and \cite{Khare_2009}), we know that $\bar{\rho}$ arises from a regular algebraic cuspidal automorphic representation $\pi_0$ of $\GL_2(\A_\Q)$.
    
    Let $\bar{\chi}=\det \bar{\rho}$ and let $\chi$ be the lifting of $\bar{\chi}$ corresponding to the central character of $\pi_0$. By \cite[Fact 4.27]{Gee_2022}, we can find a finite extension $F^1/\Q$ satisfying the following conditions:
     \begin{itemize}
    	\item $F^1$ is a totally real soluble extension of $\mathbb{Q}$ of even degree such that $p$ splits completely.
    	\item  For any $v \in \Sigma^1 \setminus \Sigma^1_p$, we have $p| \operatorname{Nm}(v)-1$, and  $\bar{\rho}|_{G_{F^1_v}}$ is trivial.
    	\item $[F^1:\mathbb{Q}] \ge 4|\Sigma^1 \setminus \Sigma^1_p|+2$.
    	\item  $\bar{\rho}|_{G_{F^1, \Sigma^1}}$ is irreducible.
    	\item  $\pi_0$ is unramified everywhere except places lying above $p$ (as an automorphic representation of $\GL_2(\A_{F^1})$),
    	
    \end{itemize}
    Here $ \Sigma^1$ (resp. $\Sigma^1_p $) is the finite set of finite places of $F^1$ lying above places in $\Sigma$ (resp. $\Sigma_p $).
    
    Let $R_{F^1}^{\operatorname{univ}}$ be the universal deformation ring parametrizing all deformations of $\bar{\rho}|_{G_{F^1, \Sigma^1}}$ with determinant $\chi|_{G_{F^1, \Sigma^1}}$. Note that from our assumptions on $F^1$, for any $v \in \Sigma^1 \setminus \Sigma^1_p$, $\chi|_{G_{F^1, \Sigma^1}}$ is unramified at $v$ and $ \chi|_{G_{F^1, \Sigma^1}}(\operatorname{Frob}_v) \equiv 1 \modd \varpi$.
    
     Write $\psi=\chi\varepsilon$ . Let $D$ be the quaternion algebra over $F^1$ which is ramified exactly at all infinite places. Fix an isomorphism between $(D \otimes_{F^1}\A_{F^1}^\infty )^\times$ and $\GL_2(\A_{F^1}^\infty)$. We define a tame level $U^p=\prod_{v\nmid p}U_v$ as follows: $U_v=\GL_2(\cO_{F^1_v})$ if $v\notin \Sigma^1$ and 
    \[U_v=\mathrm{Iw}_v:=\{g\in\GL_2(\cO_{F^1_v}),g\equiv \begin{pmatrix}*&*\\0&*\end{pmatrix}\mod \varpi_v\}\]
    otherwise. Then we can define a big Hecke algebra $\T:=\T_\psi(U^p)$, and the existence of $\pi_0$ implies that $\bar{\rho}|_{G_{F^1, \Sigma^1}}$ defines a maximal ideal $\km$ of $\T$. As $\bar{\rho}|_{G_{F^1, \Sigma^1}}$ is irreducible, there exists a two-dimensional representation $\rho_\km: G_{F^1, \Sigma^1} \to \GL_2(\T_\km)$. By the universal property, we have a surjection $R_{F^1}^{\operatorname{univ}} \twoheadrightarrow \T_\km $, which induces a surjection $ (R_{F^1}^{\operatorname{univ}})^{\operatorname{red}} \twoheadrightarrow (\T_\km)^{\operatorname{red}}$.
    
    \begin{thm}[potential pro-modularity]\label{potential pro-modularity irr}
    	The surjection $ (R_{F^1}^{\operatorname{univ}})^{\operatorname{red}} \twoheadrightarrow (\T_\km)^{\operatorname{red}}$ is an isomorphism.
    \end{thm}
    
    \begin{proof}
    	It is a direct consequence of \cite[Proposition 4.3.3]{Zhang:2024aa} and \cite[Proposition 3.1.5]{Zhang:2024ab}.
    \end{proof}
    
    For each $v \in \Sigma_p^1$ (resp. $v=p$), let $R_v^\ps$ (resp. $R_p^\ps$) be the local universal pseudo-deformation ring of the pseudo-representation $\operatorname{tr} \bar{\rho}|_{G_{F^1_v}}$ (resp. $\operatorname{tr} \bar{\rho}|_{G_{\Q_p}}$)  for $ {G_{F^1_v}}$ (resp. $ {G_{\Q_p}}$) with determinant $\chi|_{G_{F^1_v}}$ (resp. $\chi|_{G_{\Q_p}}$). Choose a subgroup $U^p_1 \subseteq U^p$ such that  $U^p_1$ satisfies the assumptions in Lemma \ref{intersection}, and for each $v \in \Sigma_p^1$, we get an ideal $I_v$ of $R_v^\ps$. Write $R_v = R_v^\ps/I_v $.  
    
    Similarly to the discussions in Subsection \ref{sec finite 2}, we get an ideal $I_p$ of $R_p^\ps$ such that for any $v \in \Sigma_p^1$, we have an isomorphism $R_v \cong R_p^\ps/I_p$ (defined as $R_p$) induced by the isomorphism $F^1_v \cong \Q_p$. Then we get the following lemma.
    
    \begin{lem}\label{com dia irr2} Write $R_{F^1, p}^\ps$ for the completed tensor product $ \widehat{\otimes}_{v \in \Sigma^1_p} R_{v}^\ps$.
    	
    	1) We have the following commutative diagram:
    	\[\begin{tikzcd}
    		R_{F^1, p}^\ps  \arrow[r, two heads] \arrow[d, two heads]& \widehat{\otimes}_{v \in \Sigma^1_p} R_{v} \arrow[d, two heads]\\
    		 R_{p}^\ps \arrow[r, two heads] &  R_{p}.
    	\end{tikzcd}\]
    	
    	2) We have $ (\widehat{\otimes}_{v \in \Sigma^1_p} R_{v})\widehat{\otimes}_{R_{F^1, p}^\ps} R_{p}^\ps = R_{p}. $
    \end{lem}
    
    Let $R_\Q^{\operatorname{univ}}$ be the universal deformation ring of $\bar{\rho}$ with determinant $\chi$. Then by the universal property, we get a natural map $f: R_{F^1}^{\operatorname{univ}} \to R_\Q^{\operatorname{univ}} $.
    
    \begin{prop}\label{finiteness 3}
    	The map $f$ induces a well-defined finite map $(R_{F^1}^{\operatorname{univ}})^{\operatorname{red}} \to (R_\Q^{\operatorname{univ}})^{\operatorname{red}}$.
    \end{prop}
    
    \begin{proof}
    	By \cite[Proposition 3.26]{Gee_2022}, the map $f$ is finite. Then our result is clear.
    \end{proof}
    
    Similarly to Lemma \ref{com dia2}, we also have the following lemma using the universal property.
    
    \begin{lem}\label{com dia irr}
    	We have the following diagram:
    	\[\begin{tikzcd}
    		R_{F^1,p}^\ps  \arrow[r] \arrow[d, two heads]& (R_{F^1}^{\operatorname{univ}})^{\operatorname{red}} \arrow[d]\\
    		 R_{p}^\ps \arrow[r] & (R_\Q^{\operatorname{univ}})^{\operatorname{red}} .
    	\end{tikzcd}\]
    \end{lem}
    
    Combining these results, we can prove a finiteness result in the residually irreducible case.
    
    \begin{prop}\label{finiteness irr}
    	There exists a finite map $ R_{p} \to (R_\Q^{\operatorname{univ}})^{\operatorname{red}}$.
    \end{prop}
    
    \begin{proof}
    	We consider the following diagram:
    	\[\begin{tikzcd}
    		R_{F^1,p}^\ps  \arrow[r, two heads] \arrow[d, two heads]& \widehat{\otimes}_{v \in \Sigma^1_p} R_{v} \arrow[d, two heads]\arrow[r, dashed] & (R_{F^1}^{\operatorname{univ}})^{\operatorname{red}} \arrow[d] \\
    		 R_{p}^\ps \arrow[r, two heads] & R_{p} \arrow[r, dashed] & (R_\Q^{\operatorname{univ}})^{\operatorname{red}}.
    	\end{tikzcd}\]
    	
    	We show that this diagram is commutative. The commutativity of the left square is just 1) of Lemma \ref{com dia irr2}. The existence of the dashed map in the first horizontal line follows from Theorem \ref{potential pro-modularity irr} and Corollary \ref{finiteness version 1}. The existence of the other dashed one and the commutativity follow from 2) of Lemma \ref{com dia irr2} and Lemma \ref{com dia irr}.
    	
    	Similarly to the proof of Proposition \ref{finiteness 2}, our result holds by using the commutative diagram, Corollary \ref{finiteness version 1}, Proposition \ref{finiteness 3} and \cite[Lemma 10.36.15]{stacks-project}.
    \end{proof}
    
    \begin{cor}\label{dim in irr case}
    	1) The ring $R_\Q^{\operatorname{univ}}$ is a local complete intersection ring of dimension $3$. 
    	
    	2) Let $R_{\Q}$ be the universal deformation ring of $\bar{\rho}$ for $G_{\Q, \Sigma}$ (without fixed determinant). Then we have $\dim R_{\Q}=4$.
    \end{cor}
    
    \begin{proof}
    	For the first assertion, it follows from Proposition \ref{R_v}, Proposition \ref{finiteness irr} and the global characteristic formula.
    	
    	For the second one, it is a direct consequence of the first one and \cite[Lemma 6.1.2]{Allen_2019}.
    \end{proof}
    
    \begin{rem}\label{rem irr}
    	It seems likely that a variant of the strategy of Subsection \ref{5.3} also works in the residually irreducible case, without imposing the Taylor–Wiles hypothesis, to show the Zariski density of non-ordinary modular points in  $\Spec R_\Q^{\operatorname{univ}}$ and $\Spec R_{\Q}$. The only missing input is an analogue of the finiteness statement used in Lemma \ref{min is non-ord}: in the residually reducible case we invoke Pan’s finiteness results from the Iwasawa algebra to the universal ordinary pseudo-deformation ring (see Theorem \ref{ord1} and Theorem \ref{ord2}). 
    	By combining our arguments with the methods of \cite{skinnerwiles01}, one should be able to establish corresponding finiteness results in the residually irreducible setting, and hence obtain an analogue of Lemma \ref{min is non-ord}. We do not pursue all the details here.
    \end{rem}
    

\bibliography{non-ordinary}
\bibliographystyle{alpha}~
\end{document}











The following result studies the reducible locus of $R_x^{\chi} $ as an analogue to .



Now we can state the main result in this subsection.

\begin{thm}\label{serre conj}
	Keep assumptions and notations above. If $[F: \Q] \ge 1+ \frac{1}{2} \dim H^1(F_\Sigma/F, \F(\bar{\chi}^{-1}))$, then for any non-zero cohomology class $x \in H^1(F_\Sigma/F, \F(\bar{\chi}^{-1}))$, $\bar{\rho}_x$ is modular. Further, the set of modular points in $\Spec R_x^{\chi}$  is Zariski dense.
\end{thm}

\begin{proof}
	Let $x \in H^1(F_\Sigma/F, \F(\bar{\chi}^{-1}))$ be a non-zero cohomology class. Then by Proposition \ref{dim of x} and Proposition \ref{red locus}, for every minimal prime of $ R_x^{\chi}$, the corresponding Galois representation is irreducible. 
	
	Let $\kp$ be a minimal prime of $R_x^{\chi} $ and let $\kp'$ be the inverse image of $\kp$ via the map $R_{F}^{\ps, \chi} \to R_x^{\chi}/\kp $. Combining Proposition \ref{finite same dim} and Proposition \ref{dim of x}, $\kp'$ is a minimal prime of $ R_{F}^{\ps, \chi}$.
	
	Now we show that the natural injection $f: R_{F}^{\ps, \chi}/\kp' \hookrightarrow R_x^{\chi}/\kp $ is finite. If not, by Nakayama's lemma, we know that the ring $R_x^{\chi}/(\kp + \mathfrak{m}^\ps) $ is not artinian, where  $\mathfrak{m}^\ps$ is the maximal ideal of $ R_{F}^{\ps, \chi}$. Assume that $\rho(\kp)$ is of the following form
	$$ \rho(\kp): G_{F, \Sigma} \to \GL_2(R_x^{\chi}/\kp), ~~ g \mapsto \begin{pmatrix}
		a_g & b_g\\
		c_g& d_g
	\end{pmatrix}.$$
	After a conjugation, we may assume $\rho(\kp)(c) = \begin{pmatrix}
		1 & 0\\
		0& -1
	\end{pmatrix}$ and $b_\sigma=1$.
	Under this settings, we have the following relations (see Subsection \ref{sec for pseudo def} for the definitions of $a,d,y$):
	$$ f(a(g))=a_g,~ f(d(g))=d_g, ~ f(y(\sigma, g))=c_g, ~ f(y(g, \tau))=b_gf(y(\sigma, \tau)), g, \tau \in G_{F, \Sigma}.$$
	
	
	
\end{proof}



